% Part: part core
\maketitle
\noteprovenance

\begin{abstract}
For coprime integers $a>b\ge1$, set $\theta=\log b/\log a$.
We specialise Zudilin's 2004 linear forms to show that
$F(a/b)=\sum_{n\ge1}((a/b)^n-1)^{-1}$ is irrational when
$\theta<\theta^*=0.4056830213840605\ldots$, with
\[
 \mu_{\rm irr}\!\left(F((a/b)^r)\right)
 \le\frac{1-\theta}{\theta^*-\theta}\qquad(r\ge1\text{ an integer}).
\]
Here $\mu_{\rm irr}$ is the irrationality exponent. In particular,
the bound is less than $301$ for every positive integral power of
$31/4$. The construction and its constants are Zudilin's; the proof
cancels common cyclotomic factors before rational specialisation and
accounts for the denominator introduced by the remaining degree.

For the distinct Hankel construction of 2016, with its auxiliary parameters
$x=z=1$, we compute
the first nonzero term of the normalised determinant at every rank:
\[
 \operatorname{ord}_q V_N^*=\frac{N(N-1)(2N-1)}6,
 \qquad [q^{\operatorname{ord}_q V_N^*}]V_N^*
       =\frac{(N!)^2(N+1)!}{2^N}.
\]
A separate positive-measure argument proves positivity and a size
estimate for each fixed real $0<q<1$. We give the full proofs,
compare the two constructions with earlier work, and distinguish
these all-rank results from finite computations about coefficient
matrices and common factors. The final sections explain the limits
of the stated clearing and congruence arguments at $3/2$. No small
nonzero sequence of divided remainders is constructed at that base,
and the full rational-base conjecture remains open.
\end{abstract}

\section{Introduction}\label{long1049:sec:problem}

Let $t>1$ be a rational number and let $\tau(n)$ count the divisors of $n$.
Erd\H{o}s Problem~\#1049 asks whether
\[
 F(t)=\sum_{n\ge1}\frac{1}{t^{n}-1}=\sum_{n\ge1}\frac{\tau(n)}{t^{n}}
\]
is irrational~\cite[p.~102]{erdos1988}.

The two series agree by expanding
$(t^{n}-1)^{-1}=\sum_{k\ge1}t^{-nk}$ and collecting equal exponents.
Nonnegativity justifies the rearrangement, while $\tau(m)\le m$ proves
absolute convergence for every real $t>1$. With $q=1/t$, the same
function is the $q$-harmonic series $\sum_{n\ge1}q^n/(1-q^n)$.
Chowla conjectured that the answer is yes for every rational $t>1$.
Erd\H{o}s proved the integer case $t\ge2$ in 1948~\cite{erdos1948};
the rational-base question is recorded in the catalogue~\cite{erdosproblems}.
The full conjecture remains open.

We specialise Zudilin's 2004 construction~\cite{zudilin2004} to prove
irrationality whenever $a>b\ge1$ are coprime and
$\log b/\log a<\theta^{*}=0.40568302138406054\ldots$.
This includes $31/4$ and every positive integral power of $31/4$.
Earlier results already treat noninteger rational bases, including
$7/2$; the point here is the explicit condition obtained from these
particular forms, not a first noninteger example.

For a reduced positive fraction $r/s$ with $r>s\ge1$, its height is
$\max(r,s)=r$. Among
noninteger rational numbers greater than one, $3/2$ has the smallest
height. It is also outside both the region proved here and the region
of Bundschuh and V\"a\"an\"anen
\cite[Thm.~2, p.~177; hypotheses pp.~175--176]{bv1994}, which contains
$7/2$. We use $3/2$ to examine why clearing denominators can destroy an
otherwise small approximation error.

The denominator cost can be read from the degrees. If $U,V\in\Z[X]$ have
degree at most $W$, then $b^WU(a/b)$ and $b^WV(a/b)$ are integers.
For this clearing argument, a small remainder
$U(a/b)F(a/b)-V(a/b)$ must remain small after multiplication by $b^W$.
Cancelling a common
polynomial factor first can reduce this cost. For example, the pair
$(X+1)(X+2),(X+1)(X+3)$ gives the cleared pair $(35,45)$ at $3/2$
with degree bound $2$; after cancelling $X+1$, degree bound $1$ gives $(7,9)$.
These illustrative polynomials are not approximants to $F$.

For the actual polynomials, Section~\ref{long1049:sec:region} proves
\[
 \log\!\left(b^{W_n}\bigl(U_n(a/b)F(a/b)-V_n(a/b)\bigr)\right)
 =\bigl(C_1\log b-C_0\log a\bigr)n^2+o(n^2),
\]
with a positive expression inside the logarithm. Here $W_n=\deg U_n$
and $\deg V_n=W_n-1$; the constants are defined in that section. A negative coefficient
of $n^2$ makes the positive values of these integer-coefficient forms
tend to zero. If $F(a/b)$
were rational, multiplication by its fixed denominator would give
positive integers tending to zero, a contradiction.

The Hankel determinant uses a different remainder sequence, from
Zudilin's 2016 construction. Its proof begins in
Section~\ref{long1049:sec:hankel-order}. We first compute its least
nonzero power of $q$ by row operations, then prove a separate estimate
at fixed real $q$. Neither argument transfers the cyclotomic factors
of the 2004 polynomials to this determinant.

The shorter companion, \emph{Zudilin's Forms at Rational Bases and
the Exact Normalised Hankel Order}, gives the principal proofs in
Sections~2 and~3; neither uses its supplementary Section~4. Here
Section~\ref{long1049:sec:coefficient-questions} contains separate
coefficient-moment and spectral calculations,
not inputs to the order formula. Sections~\ref{long1049:sec:primitive}--\ref{long1049:sec:endpoints}
compare rescaling integer rows with imposing congruences by addition;
Sections~\ref{long1049:sec:corridor}--\ref{long1049:sec:tail} examine partial sums.
Sections~\ref{long1049:sec:sevenhalves}--\ref{long1049:sec:pade} treat $7/2$
and a specified denominator-exponent model.
Section~\ref{long1049:sec:open} distinguishes a reformulation of
irrationality from questions about specified families, with an example
of a nonzero remainder that does not decay after division. Literature
comparisons follow, and Appendix~\ref{long1049:app:index} separates
formal proofs from finite computations.

All logarithms are natural. Empty products and determinants equal
$1$. A subscripted constant in $O_x(\cdot)$ may depend on the fixed
base $x$. We introduce the notation for each construction locally.

\medskip
\noindent\textbf{Keywords.} irrationality; Lambert series; rational base;
Pad\'e approximation; Lean~4.
\qquad
\textbf{MSC 2020.} 11J72 (primary); 11J82, 68V20 (secondary).

\section{The region \texorpdfstring{$b^{\mu}<a$}{b\textasciicircum mu < a} and the base \texorpdfstring{$31/4$}{31/4}}\label{long1049:sec:region}

We use Zudilin's parameter ratios $(14,12,14;27)$
\cite[Sec.~5, pp.~161--162]{zudilin2004}. Their degree calculation gives
\[
 C_1=(14+12+14)27-\tfrac12(12^2+14^2+27^2)=\frac{1091}{2},
 \qquad
 C_0=266-\frac{3}{\pi^{2}}\,(225-J),
\]
\[
 J=\sum_{i=1}^{13}\bigl(\psi_1(u_i)-\psi_1(v_i)\bigr),
 \qquad
 \psi_1(x)=\sum_{k\ge0}\frac1{(k+x)^{2}},
\]
where $\psi_1$ is the trigamma function, used only for positive arguments,
and the thirteen half-open intervals $[u_i,v_i)$ are
$[1/14,1/12)$, $[1/7,1/6)$, $[3/14,1/4)$, $[2/7,1/3)$, $[5/14,2/5)$,
$[3/7,7/15)$, $[1/2,8/15)$, $[4/7,3/5)$, $[9/14,2/3)$, $[5/7,11/15)$,
$[11/14,4/5)$, $[6/7,13/15)$, $[13/14,14/15)$.
The intervals are disjoint and lie in $[1/14,1)$, so
$0\le J\le\psi_1(1/14)-\psi_1(1)<196$: in the defining series,
$\psi_1(1/14)=196+\sum_{k\ge1}(k+1/14)^{-2}<196+\psi_1(1)$.
Since $\pi>3$, it follows that $191<C_0<266<C_1/2$.
In particular, the reciprocal constants below are positive without
using decimal approximations. The accompanying rational certificate gives
\[
\begin{aligned}
77.94318445520543&<J<77.94318449473569,\\
221.30008815898543&<C_0<221.30008817100119,\\
0.40568302137302&<\theta^*:=C_0/C_1<0.40568302139506.
\end{aligned}
\]
Put $\mu=1/\theta^*$. These parameter ratios, intervals and constants, including
$C_1/C_0=2.46497868\ldots$, occur in \cite[p.~162]{zudilin2004}; that ratio
is the integer-base exponent bound in its Theorem~1, not a new optimised
constant. The displayed interval endpoints are rounded outwards from
rational bounds; no digits beyond their certified accuracy are required.

The proof below needs $J$ in its combinatorial form as well.  Put
\[
 \omega(\xi)=\max\bigl\{0,\;
 \lfloor14\xi\rfloor+\lfloor13\xi\rfloor-\lfloor12\xi\rfloor-\lfloor15\xi\rfloor,\;
 2\lfloor14\xi\rfloor-\lfloor13\xi\rfloor-\lfloor15\xi\rfloor\bigr\},
\]
the weight that \cite[p.~162]{zudilin2004} attaches to the six-tuple
$c=(13,14,12,14,15,13)$ for the chosen parameter ratios; it gives the exponents
$\nu_l=\omega(n/l)$ of the source's~(22) and enters its limit~(26).  Zudilin
lists the support of $\omega$ in the same place; the next lemma verifies that
list by exact evaluation. Since $14+13=12+15$ and $2\cdot14=13+15$,
each floor difference is unchanged by $\xi\mapsto\xi+1$. Thus $\omega$
has period $1$, and the lemma also determines every value $\omega(n/l)$
needed for the cyclotomic exponents.

\begin{lemma}[the weight is an indicator]\label{long1049:res:omega-indicator}
On $[0,1)$ the function $\omega$ takes only the values $0$ and $1$, and
$\omega=1$ exactly on the union of the thirteen half-open intervals listed above.
\end{lemma}
\leannote{Lean: \lproof{ErdosProblems/Erdos1049/PaperOmegaIndicatorR7.lean}{1297}{omega_indicator}.}

\begin{proof}
Each of the three expressions inside the maximum is constant on every interval
between consecutive elements of
$\{k/c: c\in\{12,13,14,15\},\ 0\le k\le c\}$, since those are the only points at
which one of the four floor functions changes.  There are $48$ such intervals in
$[0,1)$; evaluating $\omega$ at the midpoint of each gives the value $1$ on the
thirteen listed half-open intervals and $0$ elsewhere, and the thirteen listed
intervals are unions of consecutive cells. Each floor function is
right-continuous, so a cell has the same value at its left endpoint as
at its midpoint. This also verifies the half-open endpoint convention.
All evaluations use exact rational arithmetic.
\end{proof}

\noindent Lemma~\ref{long1049:res:omega-indicator} and termwise integration of the positive series
$-\psi_1'(\xi)=\sum_{k\ge0}2(k+\xi)^{-3}$ give the two forms of $J$ used below,
\[
 J=\int_0^1\omega(\xi)\sum_{k\ge0}\frac{2}{(k+\xi)^{3}}\,d\xi
  =\sum_{i=1}^{13}\bigl(\psi_1(u_i)-\psi_1(v_i)\bigr),
\]
the first being the integral $\int_0^1\omega\,d(-\psi_1)$ of
\cite[Lemma~2, p.~155, and (26), p.~162]{zudilin2004}.  Since $\omega$ vanishes
near $0$, both are finite.

For $b=1$ the logarithmic ratio is zero, so the condition includes every
integer base. For $b>1$ it requires $a>b^{2.46497868\ldots}$, substantially
more than $a>b$. There are still infinitely many admissible coprime
numerators for each fixed denominator. The base $31/4$ is included,
whereas $3/2$ is excluded. Taking positive integral powers does not
change the ratio. At equality with the cutoff, the estimates give no
conclusion.

\begin{theorem}[rational-base region]\label{long1049:res:region}
Let $a>b\ge1$ be coprime integers with
\[
 b^{\mu}<a,
 \qquad\text{equivalently}\qquad
 \frac{\log b}{\log a}<\theta^{*}=0.40568302138406054\ldots
\]
Then $F(a/b)=\sum_{m\ge1}(( a/b)^{m}-1)^{-1}$ is irrational.
\end{theorem}
\leannote{Lean: \lproof{ErdosProblems/Erdos1049/PaperCompleteR21/PrintedContourConstants.lean}{418}{printed_contour}, \lproof{ErdosProblems/Erdos1049/PaperCompleteR21/PrintedContourConstants.lean}{487}{contour_enclosure}, \lproof{ErdosProblems/Erdos1049/PaperCompleteR21/PrintedContourConstants.lean}{334}{zudilinJ_enclosure}, \lproof{ErdosProblems/Erdos1049/PaperCompleteR21/PrintedContourConstants.lean}{383}{zudilinC0_enclosure}, and 3 further declarations in the \hyperref[long1049:sec:coverage]{coverage section}.}

The proof has three steps. We cancel the common polynomial factors
using Zudilin's integrality argument \cite[pp.~156--161]{zudilin2004}.
We then compute the degree of the remaining polynomials, using his
cyclotomic limits and the reciprocal-interval method of
\cite[Lemma~1, p.~466]{zudilin2002}. Finally, we evaluate at $a/b$
and show that the positive remainder still decays after multiplication
by the required power of $b$. These steps occupy
Sections~\ref{long1049:sec:source-forms}--\ref{long1049:sec:integer-forms}.
The construction, direction and constants are Zudilin's; the rational
specialisation and its degree and remainder estimates are proved below.

\subsection{The polynomial forms and their common factors}\label{long1049:sec:source-forms}

We first form $U_n$ and $V_n$ in $\Z[X]$, before choosing a rational
base. This is the step that permits cancellation to reduce the
power of the denominator needed later.

For each $n\ge1$, the coefficient formulas below are rational functions
of an indeterminate $X$. We evaluate them at real $x>1$ and put $q=x^{-1}$
when using the convergent series. Set
\[
\begin{gathered}
 a_0=14n+1,\quad a_1=12n+1,\quad a_2=14n+1,\\
 \beta_n=27n+2,\quad N=15n,\quad M_n=266n^{2}+34n+1 .
\end{gathered}
\]
The symbol $\beta_n$ is the parameter $b$ of \cite[Sec.~3, p.~155]{zudilin2004},
the exponent in the lower parameter $q^{b}$ of the Heine series~(13) on p.~157;
the letter $b$ continues to
denote the denominator of the base. In this construction alone,
$N=15n$ is the cyclotomic cutoff. The exponent $M_n$ is the integer
of that paper's~(16) on p.~157. Write
$(z;q)_m=\prod_{j=0}^{m-1}(1-zq^{j})$ for the $q$-Pochhammer symbol
(the shifted $q$-factorial). Zudilin's positive remainder is
\begin{equation}\label{long1049:eq:positive-source-form}
 H_n(x)=\sum_{t\ge0}
 \frac{(q^{t+1};q)_{a_1-1}}{(q;q)_{a_1-1}}\,
 \frac{(q;q)_{\beta_n-a_2-1}}{(q^{a_2+t};q)_{\beta_n-a_2}}\,q^{a_0t}.
\end{equation}
The denominator length $\beta_n-a_2$ agrees with the gamma
expression~(7) and residues~(8) of \cite[p.~156]{zudilin2004}.
The unnumbered display of $R(T)$ on that page instead has length
$\beta_n-a_2-1$; we use the normalisation fixed by the numbered identities.
Then \cite[(8)--(11), pp.~156--157]{zudilin2004} gives $H_n=A_nF-B_n$ with
$A_n,B_n\in\Q(X)$. We need their explicit formulas to compute degrees
and bound coefficients. Let
$\genfrac{[}{]}{0pt}{}{m}{r}_X=\prod_{j=1}^{r}(1-X^{m-r+j})/(1-X^j)$
for $0\le r\le m$. This Gaussian binomial, or $q$-binomial, polynomial has degree
$r(m-r)$ and nonnegative coefficients summing to $\binom{m}{r}$. For
$a_2\le k\le\beta_n-1$ put
\begin{equation}\label{long1049:eq:c-explicit}
 c_{n,k}(X)=(-1)^{a_1+a_2+k+1}X^{e_{n,k}}
 \genfrac{[}{]}{0pt}{}{k-1}{a_1-1}_X
 \genfrac{[}{]}{0pt}{}{\beta_n-a_2-1}{\beta_n-k-1}_X,
\end{equation}
\[
 e_{n,k}=\frac{a_1(a_1-1)}2-\frac{(\beta_n-a_2)(\beta_n-a_2-1)}2
 +\frac{(\beta_n-k)(\beta_n-k-1)}2 .
\]
Then
\begin{align}
 A_n(X)&=\sum_{k=a_2}^{\beta_n-1}c_{n,k}(X)X^{a_0k},\label{long1049:eq:A-explicit}\\
 B_n(X)&=\sum_{k=a_2}^{\beta_n-1}c_{n,k}(X)X^{a_0k}
 \biggl(\sum_{l=1}^{k-a_1}\frac1{X^{l}-1}
 +\sum_{j=1}^{a_0-1}\frac{X^{-j(k-a_1)}}{X^{j}-1}\biggr).\label{long1049:eq:B-explicit}
\end{align}
Put
\begin{equation}\label{long1049:eq:cyclotomic-products}
 D_N(X)=\prod_{l=1}^{N}\Phi_l(X),\qquad
 \nu_{n,l}=\omega(n/l),\qquad
 \Omega_n(X)=\prod_{l=2}^{N}\Phi_l(X)^{\nu_{n,l}},
\end{equation}
with $\Phi_l$ the $l$th cyclotomic polynomial; by
Lemma~\ref{long1049:res:omega-indicator} every $\nu_{n,l}$ is $0$ or $1$, and
$\nu_{n,l}=0$ for $l>15n$ because the intervals all begin at or above $1/14$.

Zudilin's Lemma~7 \cite[p.~161, (23)]{zudilin2004} applies at the parameters above: the tuple is
$c=n\cdot(13,14,12,14,15,13)$ with maximum $m(c)=15n$ and difference $s(c)=n>0$,
and the source's~(14) on p.~157 holds because $12n+1\le14n+1$ and
$26n+2\le27n+2\le28n+2$.  Its conclusion is the inclusion
\begin{equation}\label{long1049:eq:integer-polynomial-pair}
 U_n:=X^{-M_n}\frac{D_N}{\Omega_n}A_n\in\Z[X],\qquad
 V_n:=X^{-M_n}\frac{D_N}{\Omega_n}B_n\in\Z[X].
\end{equation}
Two points about that citation matter.  First, the conclusion used is the one
in $\Z[X]$, not integrality at integer arguments: an integer-valued polynomial
such as $X(X-1)/2$ shows that the two statements differ.  The source's proof supplies the polynomial form. Its Lemma~4
(pp.~157--159) clears the two rational coefficients separately by
Gaussian-binomial identities in $\Z[X]$. Lemma~5 (p.~160) makes the
normalised series invariant under the order-twelve group generated by
Heine's parameter transformation and the exchange of $a_1$ and $a_2$.
For the divisibility argument, however, only its order-six subgroup
preserving $s>0$ is used: one application of Heine's transformation
reverses the sign of $s$. Under that subgroup the normalising factorial
product takes the three forms displayed in the source. Comparing their cyclotomic
valuations gives the maximum of three terms in $\omega$.
Lemma~7 (p.~161) removes the resulting common cyclotomic factors by
polynomial divisibility; the divisor $\Omega_n$ is monic. These steps
take place before the integer specialisation in~(24) on p.~162.

Second, the module inclusion gives a polynomial coefficient pair,
and uniqueness identifies it with the displayed cancelled source pair.
If two pairs differed, subtracting their identities would express $F$
as a rational function of $X$. This is impossible: as $h\downarrow0$,
splitting the original Lambert sum $F(e^h)=\sum_{m\ge1}(e^{hm}-1)^{-1}$
at $L=\lfloor1/h\rfloor$ gives an initial sum
$h^{-1}\sum_{m\le L}m^{-1}+O(L)$, since
$y^{-1}-1\le(e^y-1)^{-1}\le y^{-1}$ for $0<y\le1$.
The remaining sum is at most
$e^{-h(L+1)}/((1-e^{-1})(1-e^{-h}))=O(h^{-1})$.  Hence
$F(e^h)=h^{-1}\log(1/h)+O(h^{-1})$, so
$(X-1)F(X)\to\infty$ and $(X-1)^{2}F(X)\to0$ as $X\downarrow1$, which no
rational function does. Thus the source's module inclusion yields the
specific pair in~\eqref{long1049:eq:integer-polynomial-pair}, not merely
some unspecified polynomial representation of the same function.
This functional nonrationality says nothing about an individual value.
For $x>1$, write the cancelled remainder as
\[
 \Lambda_n(x)=U_n(x)F(x)-V_n(x)
 =x^{-M_n}\frac{D_N(x)}{\Omega_n(x)}H_n(x).
\]

\subsection{Exact degrees}\label{long1049:sec:degrees}

The denominator cost is the degree left after cancellation. We first
identify the unique highest-degree term of $A_n$, then subtract the
degrees of the known factors.

Let $d_{n,k}=\deg\bigl(c_{n,k}(X)X^{a_0k}\bigr)$.  The Gaussian-binomial degree
formula gives
\[
 d_{n,k}=e_{n,k}+a_0k+(a_1-1)(k-a_1)+(\beta_n-k-1)(k-a_2)
 =\frac{-k^{2}+80kn+3k-340n^{2}-26n}{2},
\]
\[
 d_{n,k+1}-d_{n,k}=40n+1-k>0\qquad(a_2\le k\le\beta_n-2),
\]
the inequality because $k\le\beta_n-2=27n$ and $27n<40n+1$.  The summand of
top degree is therefore the one at $k=\beta_n-1$ alone, no cancellation occurs
there, and
\begin{equation}\label{long1049:eq:exact-degree}
 \begin{gathered}
 K_n:=\deg A_n=d_{n,\beta_n-1}=\frac{1091n^{2}+81n+2}{2},
 \\
 W_n:=\deg U_n=K_n-M_n+\sum_{l\le N}\varphi(l)-\sum_{2\le l\le N}\nu_{n,l}\varphi(l),
\end{gathered}
\end{equation}
the second equality because $\deg D_N=\sum_{l\le N}\varphi(l)$ and
$\deg\Omega_n=\sum_{2\le l\le N}\nu_{n,l}\varphi(l)$, and because
$U_n$ is a polynomial by~\eqref{long1049:eq:integer-polynomial-pair}, so subtracting
$M_n$ from the degree of $D_NA_n/\Omega_n$ is legitimate.  In particular
$U_n\ne0$. There is also no integer leading-coefficient factor to
cancel: the top summand has leading coefficient
$(-1)^{a_1+a_2+\beta_n}=(-1)^n$, and both Gaussian factors and
$D_N/\Omega_n$ are monic. Thus the leading coefficient of $U_n$ is
$(-1)^n$.

For a reduced fraction $a/b$, the unit leading coefficient also gives
the exact denominator of $U_n(a/b)$, not just an upper bound. For every prime $p\mid b$,
\[
 b^{W_n}U_n(a/b)\equiv(-1)^n a^{W_n}\not\equiv0\pmod p.
\]
No prime dividing $b$ cancels from the cleared numerator, so the
reduced denominator is exactly $b^{W_n}$. In particular every common
integer divisor of the cleared row is coprime to $b$. This excludes
cancelling a factor of $b$ from this row, not cancellation at other
primes or a saving from a different pair.

The formal arithmetic checks cover the
\lword{Erdos1049/RationalBaseContour.lean}{317}{two_mul_zudilinM}{value of $2M_n$ from the source's formula~(16)},
the
\lword{Erdos1049/RationalBaseContour.lean}{337}{twoMulZudilinSummandDegree_succ_sub}{degree difference $d_{n,k+1}-d_{n,k}=40n+1-k$},
its
\lword{Erdos1049/RationalBaseContour.lean}{350}{twoMulZudilinSummandDegree_step_pos}{positivity for $a_2\le k\le\beta_n-2$},
and the
\lword{Erdos1049/RationalBaseContour.lean}{344}{twoMulZudilinSummandDegree_top}{top degree $2K_n=1091n^2+81n+2$}.
These statements check the polynomial degree calculation, not the
analytic estimates used later.

For the second coefficient, fix $n$ and let $x\to\infty$. For $x\ge2$,
every finite product in~\eqref{long1049:eq:positive-source-form} lies
between $\prod_{j\ge1}(1-2^{-j})>0$ and $1$, independently of $t$.
Also $\sum_{t\ge0}x^{-a_0t}\le2$. These bounds give $H_n(x)=O(1)$
for the whole sum, not just each summand. Also
$F(x)=x^{-1}+O(x^{-2})$, because $\tau(1)=1$ and
$\sum_{m\ge2}m x^{-m}=O(x^{-2})$. The monic normalising factor at
infinity gives $\Lambda_n(x)=O(x^{W_n-K_n})$. Since $K_n\ge2$ and
$U_n$ has leading coefficient $(-1)^n$,
\[
 V_n(x)=U_n(x)F(x)-\Lambda_n(x)
       =(-1)^n x^{W_n-1}+O(x^{W_n-2}).
\]
The degree formula gives $W_n\ge K_n-M_n=(559n^2+13n)/2>0$.
Polynomial integrality was established earlier, so the displayed
asymptotic proves the exact statement
\begin{equation}\label{long1049:eq:V-degree}
 \deg V_n=W_n-1,\qquad \operatorname{lc}V_n=(-1)^n.
\end{equation}
In particular $V_n$ is not the zero polynomial. Applying the same
prime-by-prime reduction at its own degree shows that $V_n(a/b)$ has
reduced denominator exactly $b^{W_n-1}$. Thus $b^{W_n}$ is the least
common clearing denominator, although the second coordinate alone
needs one fewer power of $b$. In the common normalisation,
\[
 \gcd\!\bigl(b^{W_n}V_n(a/b),b^{W_n}\bigr)=b.
\]
Indeed $b^{W_n}V_n(a/b)$ is $b$ times an integer coprime to $b$.
These facts concern the fixed source pair, not all approximants.
The degree calculation fixes $n$ and varies $x$; the decay calculation
fixes $x$ and varies $n$.


\subsection{The cyclotomic limit}\label{long1049:sec:cyclotomic-limit}

The limits proved in this subsection are those of
\cite[Lemmas~1--2, p.~155]{zudilin2004}.  The proof splits the condition
$\{n/l\}\in[u,v)$ into reciprocal blocks and applies the summatory totient
estimate on each, as in the proof of \cite[Lemma~1, p.~466]{zudilin2002}, and
adds an explicit truncation estimate.

Write $\Sigma_n=\sum_{2\le l\le N}\nu_{n,l}\varphi(l)$.  The elementary
summatory estimate is $\sum_{l\le y}\varphi(l)=3\pi^{-2}y^{2}+O(y\log(2y))$.
Fix one half-open interval $[u,v)$ of Lemma~\ref{long1049:res:omega-indicator}.  The condition
$\{n/l\}\in[u,v)$ holds exactly on the blocks
\[
 \frac{n}{k+v}<l\le\frac{n}{k+u},\qquad k=0,1,2,\dots,
\]
so the summatory estimate gives, for each fixed $k$,
$n^{-2}\sum_{l\text{ in block }k}\varphi(l)\to3\pi^{-2}\bigl((k+u)^{-2}-(k+v)^{-2}\bigr)$.
For $L\ge1$, all blocks with $k\ge L$ lie below $n/(L+u)$. For every
$y\ge0$,
\[
 \sum_{l\le y}\varphi(l)\le y^2.
\]
For $y<1$ the sum is empty; otherwise it is at most
$\lfloor y\rfloor(\lfloor y\rfloor+1)/2\le y^2$. Hence the normalised
contribution of the omitted blocks is at most $(L+u)^{-2}$, uniformly
in $n$. This is the bound needed to pass from finitely many blocks to
all of them. Taking $n\to\infty$ first and then $L\to\infty$ justifies the
infinite block sum, and summing the thirteen half-open intervals gives
$n^{-2}\Sigma_n\to3J/\pi^{2}$.  With $\sum_{l\le15n}\varphi(l)=3\pi^{-2}225n^{2}+O(n\log n)$,
equation~\eqref{long1049:eq:exact-degree} yields
\begin{equation}\label{long1049:eq:degree-limits}
 \frac{K_n}{n^{2}}\to C_1,
 \qquad
 \frac{K_n-W_n}{n^{2}}=\frac{M_n-\sum_{l\le N}\varphi(l)+\Sigma_n}{n^{2}}
 \to266-\frac{3}{\pi^{2}}(225-J)=C_0,
\end{equation}
and hence $W_n/n^{2}\to C_1-C_0$.  These are the rational-base degree limits
built on Zudilin's cyclotomic limits \cite[Lemmas~1--2, p.~155; (26),
p.~162]{zudilin2004}.

For the real size estimate, fix $x>1$. The M\"obius product formula
and reindexing the divisors give
\[
 \sum_{l\le N}|\log\Phi_l(x)-\varphi(l)\log x|
 \le \sum_{d\le N}\lfloor N/d\rfloor[-\log(1-x^{-d})]
 \le N\sum_{d\ge1}\frac{-\log(1-x^{-d})}{d}.
\]
The bound $-\log(1-x^{-d})\le x^{-d}/(1-x^{-1})$ shows that the
series on the right is finite for each fixed $x>1$.
Since $\nu_{n,l}\in\{0,1\}$, this proves the bound
\begin{equation}\label{long1049:eq:cyclotomic-size}
 \log\frac{D_N(x)}{\Omega_n(x)}
 =\Bigl(\sum_{l\le N}\varphi(l)-\Sigma_n\Bigr)\log x+O_x(N).
\end{equation}
Here $N=15n$, so the error is $O_x(n)=o(n^2)$. The bracket equals
$W_n-K_n+M_n$ by~\eqref{long1049:eq:exact-degree}. No uniform constant as
$x\downarrow1$ is asserted.

\subsection{Positive linear forms with integer coefficients}\label{long1049:sec:integer-forms}

Fix the base $x=a/b>1$, put $q=1/x$ and write $P_q=(q;q)_\infty>0$.  Each of the four finite
$q$-Pochhammer products in~\eqref{long1049:eq:positive-source-form} is a product of
factors $1-q^{j}$ with $j\ge1$, hence lies in $[P_q,1]$, so each of the two
ratios lies in $[P_q,P_q^{-1}]$, every summand is positive, and
\[
 P_q^{2}\le H_n(x)\le\frac{P_q^{-2}}{1-q^{a_0}}\le\frac{P_q^{-2}}{1-q},
\]
whence $H_n(x)>0$ and $\log H_n(x)=O_x(1)$ uniformly in $n$.
The cyclotomic factors are positive on $(1,\infty)$, so the cancelled
remainder $\Lambda_n(x)$ is positive as well, by
\eqref{long1049:eq:integer-polynomial-pair}.
By~\eqref{long1049:eq:exact-degree} and~\eqref{long1049:eq:V-degree} both polynomial degrees are
at most $W_n$, so
\[
 \widehat U_n=b^{W_n}U_n(a/b),\qquad \widehat V_n=b^{W_n}V_n(a/b)
\]
are integers, and $\widehat\Lambda_n:=b^{W_n}\Lambda_n(a/b)
=\widehat U_nF(a/b)-\widehat V_n$ is positive and lies in $\Z F(a/b)+\Z$.  This
follows by multiplying each polynomial value by $b^{W_n}$.
The integer coefficients and the earlier cancellation of a common
cyclotomic factor are separate inputs.
Combining the size estimates,
\[
 \begin{aligned}
 \log\widehat\Lambda_n
 &=W_n\log b-M_n\log x+\Bigl(\sum_{l\le N}\varphi(l)-\Sigma_n\Bigr)\log x+O_x(n)\\
 &=K_n\log b-(K_n-W_n)\log a+o(n^{2}),
\end{aligned}
\]
the second equality by $\log x=\log a-\log b$ and~\eqref{long1049:eq:exact-degree}.
With~\eqref{long1049:eq:degree-limits} this proves
\begin{equation}\label{long1049:eq:final-limit}
 \lim_{n\to\infty}\frac{\log\widehat\Lambda_n}{n^{2}}=C_1\log b-C_0\log a .
\end{equation}

\begin{proof}[Proof of Theorem~\ref{long1049:res:region}]
The hypothesis $\log b/\log a<\theta^{*}=C_0/C_1$ makes the right side
of~\eqref{long1049:eq:final-limit} negative, so $\widehat\Lambda_n>0$ and
$\widehat\Lambda_n\to0$.  Suppose $F(a/b)=P/Q$ with integers $P$ and $Q\ge1$.
Then $Q\widehat\Lambda_n=P\widehat U_n-Q\widehat V_n$ is a positive integer for
every $n$ and tends to $0$, which is impossible.
\end{proof}

\noindent Above the threshold~\eqref{long1049:eq:final-limit} says that these same
homogenised forms grow like $e^{cn^{2}}$ with $c>0$, and at equality the limit
decides nothing.  Both statements are about the displayed family.

\noindent The hypothesis $b\ge1$ admits $b=1$, where the statement reduces to
the known integer-base theorem. Coprimality fixes the reduced
representation of the base. It is used in the exact-denominator
assertion above, but is not needed for the sufficient irrationality
implication once the displayed logarithmic inequality holds.
Negative bases are excluded,
because the proof uses positivity for real bases greater than $1$.

\begin{sloppypar}
\begin{theorem}[the base $31/4$ outside the Bundschuh--V\"a\"an\"anen
region]\label{long1049:res:31over4}
$F(31/4)$ is irrational, and so is $F\bigl((31/4)^{r}\bigr)$ for every integer
$r\ge1$.  Here
\[
 \frac{\log4}{\log31}=0.4036981731641997\ldots<\frac{81}{200}<\theta^{*},
\]
while
\[
 4^{\mu}=30.483515\ldots<31<4^{\mu_{\mathrm{BV}}}=32.369642\ldots
\]
with
$\mu_{\mathrm{BV}}=2\pi^{2}/(\pi^{2}-2)=2.508284761994\ldots$, so $31/4$ lies
outside the region $\log b/\log a<1/2-1/\pi^{2}=0.3986788163576622\ldots$ of
\cite[Thm.~2, p.~177]{bv1994} and inside the region of
Theorem~\ref{long1049:res:region}.
\end{theorem}
\leannote{Lean: \lproof{ErdosProblems/Erdos1049/PaperCompleteR21/PrintedLogConstants.lean}{118}{printed_log_ratio}, \lproof{ErdosProblems/Erdos1049/PaperCompleteR21/PrintedLogConstants.lean}{195}{printed_bv_cutoff}, \lproof{ErdosProblems/Erdos1049/PaperCompleteR21/PrintedLogConstants.lean}{218}{printed_bvMu}, \lproof{ErdosProblems/Erdos1049/PaperCompleteR21/PrintedLogConstants.lean}{308}{printed_four_rpow_bvMu}, and 9 further declarations in the \hyperref[long1049:sec:coverage]{coverage section}.}
\end{sloppypar}

\begin{proof}
The comparison $\log4/\log31<81/200$ is the integer certificate
$4^{200}<31^{81}$, checked as
\lword{Erdos1049/ZudilinHeightRegion.lean}{26}{thirtyoneFour_power_certificate}{the
power certificate}, and the membership it yields is
\lword{Erdos1049/ZudilinHeightRegion.lean}{45}{thirtyoneFour_mem_zudilinHeightRegion}{$31/4$ satisfies the inequality};
the ratio $\log b/\log a$ is invariant under $(a,b)\mapsto(a^{r},b^{r})$, which
gives the
\lword{Erdos1049/ZudilinHeightRegion.lean}{59}{thirtyoneFour_power_mem_zudilinHeightRegion}{power
family}, and $(31^{r},4^{r})$ are coprime.  The exclusion from the earlier
region is
\lword{Erdos1049/ZudilinHeightRegion.lean}{91}{thirtyoneFour_outside_bundschuhVaananenHeightRegion}{$31/4$ is outside the earlier region}.
That exclusion also has a two-line rational certificate: $31^{2}<4^{5}$ gives
$\log4/\log31>2/5$, and $\pi^{2}<10$ gives $1/2-1/\pi^{2}<2/5$, so
\[
 \frac12-\frac1{\pi^{2}}<\frac25<\frac{\log4}{\log31}<\frac{81}{200}<\theta^{*}.
\]
The remaining comparison $81/200<\theta^{*}$ is proved in
Section~\ref{long1049:sec:regionbracket}.  Theorem~\ref{long1049:res:region} then applies.
\end{proof}

\begin{corollary}[an irrationality measure uniform over powers]
\label{long1049:cor:rational-base-measure}
For coprime $a>b\ge1$ with $\theta=\log b/\log a<\theta^*$ and every
integer $r\ge1$,
\[
 \mu_{\rm irr}\!\left(F((a/b)^r)\right)
 \le\frac{1-\theta}{\theta^*-\theta}.
\]
Here $\mu_{\rm irr}(\xi)$ is the supremum of the exponents $\nu$ for which
$|\xi-p/q|<q^{-\nu}$ has infinitely many reduced rational solutions.
In particular, $\mu_{\rm irr}(F((31/4)^r))<301$ for every $r\ge1$.
\end{corollary}
\leannote{Lean: \lproof{ErdosProblems/Erdos1049/PaperCompleteR21/RationalBaseThreshold.lean}{194}{rational_base_measure_uniform}, \lproof{ErdosProblems/Erdos1049/PaperCompleteR21/RationalBaseThreshold.lean}{205}{thirtyoneFour_power_measure_lt_301}, \lproof{ErdosProblems/Erdos1049/PaperR17/SourceConsumers.lean}{167}{rational_base_power_measure}, \lproof{ErdosProblems/Erdos1049/PaperR17/SourceConsumers.lean}{176}{thirtyone_four_power_measure_lt_301}.}
\begin{proof}
The coefficient $A_n$ in~\eqref{long1049:eq:A-explicit} is a sum of
$O(n)$ Laurent monomials times two Gaussian binomial polynomials.
Each Gaussian polynomial has nonnegative coefficients summing to at
most $2^{27n+2}$. Hence the sum of the absolute coefficients of $A_n$
is $\exp(O(n))$. Since its largest exponent is $K_n$,
$|A_n(x)|\le x^{K_n}\exp(O(n))$ for each fixed $x>1$.
The cyclotomic estimate in Section~\ref{long1049:sec:integer-forms}
bounds the multiplier $x^{-M_n}D_N(x)/\Omega_n(x)$ by
$x^{W_n-K_n}\exp(O_x(n))$. Therefore
\[
 |U_n(x)|\le x^{W_n}\exp(O_x(n))\qquad(x>1\text{ fixed}).
\]
Set $\xi=F(a/b)$, $Q_n=b^{W_n}U_n(a/b)$ and
$P_n=b^{W_n}V_n(a/b)$.  With
\[
 \alpha=(C_1-C_0)\log a,\qquad
 \tau=C_0\log a-C_1\log b>0,
\]
we have $\log(Q_n\xi-P_n)=-\tau n^2+o(n^2)$ and
$|Q_n|\le\exp(\alpha n^2+o(n^2))$.

For integers $A,B,p,q$ with $q>0$, if $L=A\xi-B$ and $2q|L|\le1$, then
\[
 |L|\le |A|\,|\xi-p/q|.
\]
When $Ap-Bq=0$ this is equality.  Otherwise the nonzero integer $Ap-Bq$
gives $1/q\le |L|+|A|\,|\xi-p/q|$, which proves the inequality.
Fix $0<\eta<\tau$. For all sufficiently large $n$, the estimates above give
\[
 e^{-(\tau+\eta)n^2}\le Q_n\xi-P_n\le e^{-(\tau-\eta)n^2},
 \qquad |Q_n|\le e^{(\alpha+\eta)n^2}.
\]
For every sufficiently large denominator $q$, choose
$n=\lceil\sqrt{\log(2q)/(\tau-\eta)}\rceil$. Then
$2q(Q_n\xi-P_n)\le1$, so the preceding integer argument applies to
$(A,B)=(Q_n,P_n)$ for every numerator $p$. Also $Q_n\ne0$: otherwise
$Q_n\xi-P_n$ would be an integer strictly between $0$ and $1$.
It follows that
\[
 |\xi-p/q|\ge e^{-(\alpha+\tau+2\eta)n^2}
             =q^{-(\alpha+\tau+2\eta)/(\tau-\eta)-o(1)},
\]
since $n^2=\log(2q)/(\tau-\eta)+O_\eta(\sqrt{\log q}+1)$.

Let $\eta\downarrow0$.  The resulting bound is
$1+\alpha/\tau=(1-\theta)/(\theta^*-\theta)$.
Taking a common power multiplies $\alpha,\tau$ by $r$, leaving this quotient
unchanged; the constants in the approximation inequality may depend on $r$.
For $31/4$, the exact rational interval calculation below gives
\[
\begin{gathered}
 0.40568302137302<\theta^*<0.40568302139506,\\
 0.40369817316419<\frac{\log4}{\log31}<0.40369817316420.
\end{gathered}
\]
For each trigamma difference, sum $k=0,\ldots,255$ exactly and bound the
remaining decreasing positive summand $f_{u,v}(x)=(x+u)^{-2}-(x+v)^{-2}$
by its integral from $256$ to infinity and that integral plus $f_{u,v}(256)$.
Bound $\pi$ with Machin's identity and alternating arctangent series,
and logarithms with the positive $\operatorname{arctanh}$ series and a
geometric tail bound. Outward rational interval operations then place
$(1-\rho)/(\theta^*-\rho)$, where $\rho=\log4/\log31$,
between $300.4269130$ and $300.4269164$, hence below $301$.
All endpoints are rational. For the bound $301$ alone, the coarser
inequalities $\rho<0.4036982$ and $\theta^*>0.40568$ suffice:
\[
 \frac{1-\rho}{\theta^*-\rho}
 <\frac{1-0.4036982}{0.40568-0.4036982}
 =\frac{2981509}{9909}<301.
\]
Here the quotient is increasing in $\rho$ and decreasing in $\theta^*$
because $\rho<\theta^*<1$. The finer enclosure also certifies the
displayed decimal approximation; neither calculation is a new Lean theorem.
\end{proof}

\subsection{The rational bracket around \texorpdfstring{$\theta^{*}$}{theta*}}\label{long1049:sec:regionbracket}

To check a strict comparison with $\theta^{*}$, we enclose the constant
between rational numbers. The coarse bounds needed for $31/4$ and
$3/2$ follow from finite rational inequalities, without using a decimal
expansion.

For the lower bound, keep only the $k=0$ term of each of the thirteen
differences $\psi_1(u_i)-\psi_1(v_i)$.  All later terms are positive, so
\[
 J\ \ge\ \sum_{i=1}^{13}\Bigl(\frac1{u_i^{2}}-\frac1{v_i^{2}}\Bigr)
  =\frac{2015640690251}{25971865920}
  >\frac{776}{10}.
\]
Since $\pi>157/50$ and $225-J<225$,
\[
 C_0>266-\frac{3\,(225-776/10)}{(157/50)^{2}}=\frac{5451134}{24649}
 >\frac{88371}{400}=\frac{81}{200}\,C_1,
\]
so $81/200<\theta^{*}$.  For the upper bound the thirteen half-open intervals are
disjoint and ordered and $\psi_1$ is positive and decreasing, so the sum
telescopes below its first term:
\[
 J<\psi_1(1/14)=196+\sum_{k\ge1}\frac1{(k+1/14)^{2}}<196+\frac{\pi^{2}}6<198<225,
\]
whence $C_0<266$ and $\theta^{*}<532/1091<1/2$.  So
\[
 \frac{81}{200}<\theta^{*}<\frac12 .
\]

The module \texttt{RationalBaseContour}, imported by the library root at
revision~\commitshort, checks the definitions of $C_0$, $C_1$ and
$\theta^*$ together with the
\lword{Erdos1049/RationalBaseContour.lean}{159}{zudilinJ_ge}{lower bound on $J$},
\lword{Erdos1049/RationalBaseContour.lean}{180}{zudilinC0_gt}{rational lower bound on $C_0$},
and the two comparisons
\lword{Erdos1049/RationalBaseContour.lean}{201}{eightyOne_twoHundredths_lt_zudilinContour}{$81/200<\theta^*$}
and
\lword{Erdos1049/RationalBaseContour.lean}{251}{zudilinContour_lt_half}{$\theta^*<1/2$}.
It also checks that
\lword{Erdos1049/RationalBaseContour.lean}{266}{thirtyoneFour_mem_zudilinContourRegion}{$31/4$ belongs to the region},
that \lword{Erdos1049/RationalBaseContour.lean}{279}{thirtyoneFour_power_mem_zudilinContourRegion}{every positive integral power does too},
and that \lword{Erdos1049/RationalBaseContour.lean}{296}{threeHalves_outside_zudilinContourRegion}{$3/2$ is excluded}.
These comparisons use the defined constant; they do not supply the
analytic estimates in Theorem~\ref{long1049:res:region}. The sharper
interval in the preceding proof instead comes from the stated
exact-arithmetic certificate with $256$ summands per trigamma
difference. It does not rely on the historical thousand-term decimal
evaluation reported in Section~\ref{long1049:sec:receipts}.

The shorter interval above suffices for the bound $301$ but does not
certify all digits printed in the cutoff. For those digits, a separate
integer-arithmetic calculation gives
\[
\begin{aligned}
0.4056830213840605403&<\theta^*<0.4056830213840605417,\\
2.4649786835749750334&<\mu<2.4649786835749750415.
\end{aligned}
\]
Here is the complete error bound used in
\path{computations/check_threshold_digits.py}. For each of the thirteen
intervals, put $f_{u,v}(k)=(k+u)^{-2}-(k+v)^{-2}$.
Round each of the first $M=65536$ positive rational summands down to a
multiple of $10^{-50}$. Their true sum is between this rounded sum and
that sum plus $M10^{-50}$. Since $f_{u,v}$ is positive and decreasing,
the remaining tail lies between its integral from $M$ to infinity
and that integral plus $f_{u,v}(M)$. The integral is
$(M+u)^{-1}-(M+v)^{-1}$. Summing these rational intervals and using
Machin's alternating-series bounds for $\pi$ gives the displayed
bounds on $\theta^*$ and its reciprocal. The coefficient
$225-J$ is positive, so the interval endpoints for $C_0$ have the
claimed order. The file \path{computations/threshold_digits.json}
records the intervals. This verifies the displayed decimal truncations;
the exact definitions, not the decimal values, are used in the proofs.

\subsection{Proofs, earlier work and limitations}

\paragraph{Formal verification and computations.}
Theorems~\ref{long1049:res:region} and~\ref{long1049:res:31over4} are ordinary proofs.  They
use Zudilin's construction: his Lemma~7 in its polynomial reading, together with the
inputs of that lemma's own proof, namely his Lemma~3, the exponent $M(a;b)$
of~(16) proved by his Lemma~4, the group stability of his Lemma~5, and the
identity~(9)--(11), all in \cite[pp.~156--161]{zudilin2004}, and his direction and
constants \cite[p.~162]{zudilin2004}.  Proved here are
the nonrationality of $F$ as a function, the identification of the
polynomial coefficients in Lemma~7, the exact degrees of $U_n$ and $V_n$,
positivity, the Archimedean size estimate, and the limit
$(K_n-W_n)/n^{2}\to C_0$.  That last limit rests on Zudilin's Lemmas~1 and~2
\cite[p.~155]{zudilin2004}, whose method goes back to
\cite[Lemma~1, p.~466]{zudilin2002}; Section~\ref{long1049:sec:cyclotomic-limit}
proves them by that method with an explicit truncation estimate.  The supplied public sources prove the irrationality and measure
statements for the constructed polynomials. They also contain the
finite comparisons, the rational bracket in
Section~\ref{long1049:sec:regionbracket} and the degree identities in
Section~\ref{long1049:sec:degrees}. Appendix~\ref{long1049:app:index}
distinguishes these proofs from computer algebra.

The companion files contain exact certificates for the rational interval
bounds, $U_n,V_n$ for $1\le n\le4$, the sixteen coefficient-positivity
tests through rank eight, and the finite parameter box. The five complete
polynomial contents and seventy-six residue certificates are described in
Section~\ref{long1049:sec:coefficient-questions}. The historical
high-precision real remainder evaluations are separate numerical reports,
not reproduced by these exact computations. None of these finite
computations is needed for the proof of
Theorem~\ref{long1049:res:region}.
Theorem~\ref{long1049:res:region} is checked in Lean as
\href{https://github.com/wcook04/plectis-erdos/blob/f4e61ed8a6941571310049e22dc239baa1b04e12/lean/ErdosProblems/Erdos1049/PaperR17/SourceConsumers.lean#L113}{the rational-base region},
using the constructed polynomials and proved remainder estimates,
rather than assuming that suitable forms exist.

\paragraph{Attribution.}
Bundschuh and V\"a\"an\"anen proved the irrationality of $F(a/b)$ on the
region $\log b/\log a<1/2-1/\pi^{2}=0.3986788163576622\ldots$
\cite[Thm.~2, p.~177; hypotheses pp.~175--176]{bv1994}, and every base of
Theorem~\ref{long1049:res:region} with $b\le3$ is already theirs.  Their printed
hypothesis for $\alpha=-1$ is $\lambda<(1/2+1/\pi^{2})^{-1}$ with
$\lambda=\log h(q)/\log|q|$. Here the source uses $q$ for the base,
not its reciprocal, and
\[
 E_q(z)=\prod_{j\ge1}(1+zq^{-j}),\qquad
 L_q(z)=\frac{E_q'(z)}{E_q(z)}=\sum_{j\ge1}\frac1{q^j+z}.
\]
For $q=a/b>1$, $E_q(-1)>0$ and $L_q(-1)=F(a/b)$; in particular
$\alpha=-1$ is not one of the excluded zeros $-q^j$.
Since $a,b$ are coprime and $a>b>0$, the rational height is $h(q)=a$.
Thus $\lambda=\log a/\log(a/b)=1/(1-\log b/\log a)$, and the
printed hypothesis is equivalent to the displayed inequality.
Duverney proved another, strictly smaller, region for the same series
\cite[Th\'eor\`eme~2, p.~174]{duverney1996}.  Zudilin remarked that his
generalized $q$-logarithm results can be given at non-integer rational bases
under an assumption $\log|r|>c\log|s|$ for a computable $c>0$, without computing
a value \cite[Sec.~2, p.~4]{zudilin2016}.  Zudilin's 2004 paper supplies the forms and the exponent
$\mu$ under the standing hypothesis $p=1/q\in\Z\mathbin{\backslash}\{0,\pm1\}$, and it
states no rational-base result \cite[Sec.~2, p.~154; Thm.~1]{zudilin2004}.  The
contribution here is the rational specialisation of the 2004 forms, with the
denominator accounting and the limit passage carried out in full, which
identifies the printed $\mu$ as an admissible $c$ for $F$ itself, together with
the region that constant defines and its application to $31/4$.

\paragraph{Where the earlier criterion stops.}
Both regions are cut out by the same quantity.  The published cutoff
$1/2-1/\pi^{2}$ of \cite[Thm.~2, p.~177]{bv1994} is the reciprocal of
$\mu_{\mathrm{BV}}=2\pi^{2}/(\pi^{2}-2)$, the constant Van Assche later
recovered as an integer-base irrationality-exponent bound for $F(p)$
\cite[Thm.~1, p.~10]{vanassche2001}, and the cutoff $\theta^{*}$ of
Theorem~\ref{long1049:res:region} is the reciprocal of the smaller integer-base bound
$\mu=C_1/C_0$ printed at \cite[p.~162]{zudilin2004}.  For Zudilin's family this reciprocal relation uses both the exact degree
limit $d=C_1-C_0$ and the stronger evaluation bound
$|U_n(p)|\le p^{W_n}\exp(O_p(n))$ proved in
Corollary~\ref{long1049:cor:rational-base-measure}, together with decay
exponent $\sigma=C_0$. A general exponential coefficient-height bound
would contribute an additional term to the integer-base exponent
estimate; it cannot simply be discarded. The discussion following
Theorem~\ref{long1049:res:archcap} makes that distinction explicit.
The rational-base proof clears the denominators of the already
cancelled polynomials. It does not infer value irrationality merely by
taking the reciprocal of a published exponent bound.  The bases
gained are exactly the strip $s^{\mu}<r\le s^{\mu_{\mathrm{BV}}}$, which is empty
for $s=2$ and $s=3$ and is first occupied at $s=4$ by the single coprime numerator
$31$.  The integer comparisons and rational bracket in
Section~\ref{long1049:sec:regionbracket} establish membership in the region.
The irrationality conclusion in Theorem~\ref{long1049:res:31over4} also uses
the constructed forms and their proved analytic estimates. This application
does not improve either inherited integer-base exponent bound.

\paragraph{What a larger region would require.}
At $3/2$ the logarithmic parameter is
\[
 \frac{\log2}{\log3}=0.6309297535714574\ldots,
\]
which exceeds $\theta^*$ by $0.2252467\ldots$.
Theorem~\ref{long1049:res:region} gives no conclusion when
$\log b/\log a\ge\theta^*$. The same argument can apply to another family
once its polynomial integrality, degrees and fixed-base remainder
estimates are proved. A smaller ratio of the resulting constants
$C_1/C_0$ would enlarge the reciprocal-cutoff region. For example,
$C_1/C_0\le2.4234$ would give a cutoff greater than $0.4126$ and admit
$29/4$. An integer-base irrationality-exponent bound alone does not
supply the polynomial data needed for this specialisation.
The existing formal proofs establish irrationality and measure bounds
for the constructed forms; they do not supply a different family
whose divided remainder is small and nonzero at $3/2$.

The following condition is stronger than having estimates at $3/2$:
one polynomial family must work at every fixed $x>1$ with the same
leading constants. The error terms may depend on $x$. The 2004 family
above satisfies these hypotheses. Estimates for that family at only one
base, or for a different family at each base, do not verify them.
The conclusion concerns
these two-coordinate linear forms, not simultaneous forms in several
independent target values. Here height means the largest absolute
coefficient. The short paper uses the sum of the absolute coefficients;
for a polynomial of degree at most $d$,
\[
 \max_j|p_j|\le\sum_j|p_j|\le(d+1)\max_j|p_j|.
\]
For a nonzero polynomial of degree $d=O(n^2)$, the logarithms of
these norms differ by $O(\log n)$. They have the same leading
quadratic growth rate, but the literal zero-height conditions are
not identical: $1+X$ has maximum coefficient $1$ and coefficient sum
$2$. In particular, when $h=0$, the displayed bound
$hn^2(1+o(1))$ is zero, not an arbitrary $o(n^2)$ term.
The proof below also works with the additive bound $hn^2+o(n^2)$;
its evaluation estimate absorbs the $O(\log n)$ difference.
Alternatively, either norm convention satisfies the other paper's
hypothesis after replacing $h$ by any larger positive constant.
The degree conclusion is independent of that replacement.
In the rational-base conclusions, write $a/b$ with integers $a>b\ge1$.
The estimates hold without coprimality; a reduced representation gives
the smaller denominator-clearing factor.

\begin{theorem}[a degree restriction for estimates valid at every base]\label{long1049:res:archcap}
Let $(U_n,V_n)\in\Z[x]^{2}$ be a sequence such that, for constants
$\sigma,\delta>0$ and $h\ge0$ independent of $n$ and of the base,
\begin{enumerate}
\item $\Lambda_n(x):=U_n(x)F(x)-V_n(x)\ne0$ for every real $x>1$;
\item $\deg U_n,\deg V_n\le\delta n^{2}(1+o(1))$;
\item the coefficient heights satisfy
\[
 \log\max\bigl(H(U_n),H(V_n)\bigr)\le hn^{2}(1+o(1)),
\]
where $H(P)$ is the
largest absolute value of a coefficient of $P$;
\item the remainders satisfy
\[
 \log|\Lambda_n(x)|=-\sigma n^{2}\log x\,(1+o(1))
\]
for every real $x>1$.
\end{enumerate}
Put $d_n:=\max(\deg U_n,\deg V_n)$.  Then $\sigma\le\delta$, and for every fixed
rational base $a/b>1$,
\[
 \limsup_{n\to\infty}n^{-2}\log\bigl|b^{d_n}\Lambda_n(a/b)\bigr|
 \le\delta\log b-\sigma\log(a/b).
\]
Consequently the homogenised forms tend to zero whenever
$\log b/\log a<\sigma/(\sigma+\delta)$, a sufficient region whose cutoff is at
most $1/2$.  If the actual degrees satisfy $d_n/n^{2}\to d$, then $d\ge\sigma$
and the limit exists and equals $d\log b-\sigma\log(a/b)$; in that case the
forms tend to zero below $\log b/\log a=\sigma/(\sigma+d)$ and their absolute
values tend to infinity above it, so the exact-degree case has no decaying
homogenised forms at $3/2$.
\end{theorem}
\leannote{Lean: \lproof{ErdosProblems/Erdos1049/PaperLongCapR9.lean}{411}{long_record_archcap}.}

\begin{proof}
The coefficient-height bound controls evaluation at an integer base.
For every fixed $p>1$,
\[
 |U_n(p)|,|V_n(p)|
 \le (d_n+1)\max(H(U_n),H(V_n))p^{d_n}
 \le \exp\bigl((h+\delta\log p)n^2+o(n^2)\bigr).
\]
Here $\log(d_n+1)=o(n^2)$ by the degree bound. This is where the
coefficient-height hypothesis is needed.

Suppose $\sigma>\delta$ and choose an integer $p\ge2$ such that
$(\sigma-\delta)\log p>h$. Put $u_n=U_n(p)$, $v_n=V_n(p)$ and
$\ell_n=u_nF(p)-v_n$. The evaluation bound and hypothesis~(4) give
\[
 u_nv_{n+1}-u_{n+1}v_n=u_{n+1}\ell_n-u_n\ell_{n+1}=o(1).
\]
Indeed, each product on the right has absolute value at most
$\exp((h+(\delta-\sigma)\log p)n^2+o(n^2))$.
The expression on the left is an integer, so it is eventually zero.
Also $u_n\ne0$ eventually, since $u_n=0$ would make $\ell_n=-v_n$
a nonzero integer of absolute value less than $1$.
Thus $v_n/u_n$ is eventually a fixed rational number $r$.
If $F(p)=r$, then $\ell_n=0$; otherwise
$|\ell_n|=|u_n|\,|F(p)-r|\ge|F(p)-r|$.
Both contradict the nonzero remainders tending to zero. Hence
$\sigma\le\delta$.

For a fixed rational base $a/b>1$, taking logarithms gives
\[
 n^{-2}\log|b^{d_n}\Lambda_n(a/b)|
 =\frac{d_n}{n^2}\log b-\sigma\log(a/b)+o(1).
\]
Hypothesis~(2) yields the stated upper limit and sufficient region.
Since $\sigma\le\delta$, its cutoff is at most $1/2$.
If $d_n/n^2\to d$, apply the preceding integer-base argument with
$d+\varepsilon$ in place of $\delta$ for every $\varepsilon>0$ to obtain
$d\ge\sigma$. The displayed identity then gives the exact limit.
Its sign is that of $(\sigma+d)\log b-\sigma\log a$, which proves
both assertions away from equality. At $3/2$ this sign is positive
because $\log2/\log3>1/2\ge\sigma/(\sigma+d)$.
\end{proof}

\begin{corollary}[nondecay when $b<a<b^2$]
\label{long1049:cor:no-decay-below-square}
Under the hypotheses of Theorem~\ref{long1049:res:archcap}, for positive integers
$a,b$ with $b<a<b^2$, the undivided forms $b^{d_n}\Lambda_n(a/b)$ do not
tend to zero.  No limit of $d_n/n^2$ is assumed.
\end{corollary}
\leannote{Lean: \lproof{ErdosProblems/Erdos1049/PaperNoDecayR9.lean}{69}{cleared_below_square_not_tendsto_zero}.}
\begin{proof}
Write $x=a/b$ and suppose $C_n=b^{d_n}\Lambda_n(x)\to0$.
Eventual nonvanishing gives
$\log|C_n|=d_n\log b+\log|\Lambda_n(x)|$ for all sufficiently large $n$.
Also $|C_n|\le1$ eventually.  The lower side of the remainder asymptotic
therefore implies
\[
 d_n\log b\le-\log|\Lambda_n(x)|
       =\sigma\log x\,n^2+o(n^2).
\]
Set $c=\sigma\log x/\log b$.  Since $1<x<b$, we have $0<c<\sigma$,
and the displayed inequality gives $d_n\le(c+\varepsilon)n^2$ eventually
for each $\varepsilon>0$. Apply Theorem~\ref{long1049:res:archcap}
to the same family with degree upper rate $c$. Its height bound,
nonvanishing and remainder asymptotics at every real base greater than
one are unchanged. The theorem gives $\sigma\le c$, a contradiction.
\end{proof}
\noindent The no-decay conclusion is \href{https://github.com/wcook04/plectis-erdos/blob/3be82b1a7340284aea72e9a5c8493cb020843921/ErdosProblems/Erdos1049/PaperNoDecayR9.lean#L69}{kernel-checked in Lean} under the stated all-base hypotheses.


This strengthens the failure of a sufficient-cutoff test to an exclusion
of decay under the stated all-base hypotheses.  It does not supply an
actual-degree limit or assert divergence.  If $d_n/n^2\to d$, the separate
exact-degree result gives the stronger conclusion
$|b^{d_n}\Lambda_n(a/b)|\to\infty$ in the same strict region.
Equality $a=b^2$ remains unclassified.  The corollary concerns the undivided
forms; it does not exclude base-dependent content division or other
irrationality methods outside its hypotheses.


\noindent The coefficient-height hypothesis (3) supplies the evaluation bound
used in this proof. A degree bound alone gives no such estimate:
polynomials of degree zero can have arbitrarily large integer
coefficients. The constant $h$ is independent of the base, which lets
us choose one large integer $p$ with $(\sigma-\delta)\log p>h$
under the contradiction hypothesis $\sigma>\delta$.
This explains the use of (3); it does not prove that the hypothesis
cannot be weakened or omitted from the theorem.  At the boundary $\log b/\log a=\sigma/(\sigma+d)$ the normalised
logarithm is zero and these hypotheses decide neither behaviour.
Theorem~\ref{long1049:res:archcap} constrains families satisfying its hypotheses and
does not exclude every possible Pad\'e construction.

There is also a distinction between the degree cutoff and an
irrationality-exponent estimate at a fixed integer base $p\ge2$.
The hypotheses above give
\[
 |U_n(p)|\le\exp((h+\delta\log p)n^2+o(n^2)),\qquad
 |\Lambda_n(p)|=\exp(-\sigma\log p\,n^2+o(n^2)).
\]
Apply the integer-form argument in the proof of
Corollary~\ref{long1049:cor:rational-base-measure}, with coefficient
growth rate $h+\delta\log p$ and remainder decay rate $\sigma\log p$.
It gives
\[
 \mu_{\rm irr}(F(p))\le
 1+\frac{\delta}{\sigma}+\frac{h}{\sigma\log p}.
\]
Thus the reciprocal $\sigma/(\sigma+\delta)$ of the degree expression
$1+\delta/\sigma$ is not, under these hypotheses alone, the reciprocal
of the exponent bound furnished by that argument. To obtain the latter
identification it suffices to prove the stronger evaluated estimate
$|U_n(p)|\le\exp(\delta\log p\,n^2+o(n^2))$. This is separate from
bounding the polynomial's coefficients by $\exp(O(n^2))$.

For the family of Section~\ref{long1049:sec:source-forms} the fourth hypothesis is the
size estimate proved there and the second is~\eqref{long1049:eq:exact-degree}.  The third
is proved next, so Theorem~\ref{long1049:res:archcap} applies to that family.

\begin{lemma}[coefficient heights of the constructed polynomials]\label{long1049:res:sourceheight}
There is a constant $h$ with
$\log\max\bigl(H(U_n),H(V_n)\bigr)\le hn^{2}$ for every $n\ge1$, where $U_n$ and
$V_n$ are the polynomials of~\eqref{long1049:eq:integer-polynomial-pair}.
\end{lemma}
\leannote{Lean: \lproof{ErdosProblems/Erdos1049/PaperCompleteR21/SourceCoefficientHeights.lean}{49}{exists_quadratic_source_height_bound}, \lproof{ErdosProblems/Erdos1049/PaperCompleteR21/SourceCoefficientHeights.lean}{40}{maxPairHeight_source_eq}.}

\begin{proof}
For a polynomial or Laurent polynomial $P$, write $\|P\|$ for the sum
of the absolute values of its coefficients. Then
$\|PQ\|\le\|P\|\,\|Q\|$; for an ordinary polynomial,
$H(P)\le\|P\|$. This convention includes the Laurent monomials in
$c_{n,k}$.

By Lemma~\ref{long1049:res:omega-indicator} every $\nu_{n,l}$ is $0$ or $1$, so
$D_N/\Omega_n=\prod_{l\in S}\Phi_l$ over a subset $S\subseteq\{1,\dots,N\}$.
A monic polynomial with roots $\zeta_1,\ldots,\zeta_d$ on the unit circle
has coefficient norm at most $2^d$: expanding its linear factors bounds
the sum of the absolute coefficients by
$\prod_{j=1}^d(1+|\zeta_j|)=2^d$. Hence $\|\Phi_l\|\le2^{\varphi(l)}$ and
\[
 \Bigl\|\frac{D_N}{\Omega_n}\Bigr\|\le2^{\sum_{l\le N}\varphi(l)}\le2^{225n^{2}} .
\]
For $A_n$, the Gaussian binomial $\genfrac{[}{]}{0pt}{}{m}{r}_X$ has nonnegative
coefficients summing to $\binom{m}{r}$, so
$\|c_{n,k}\|\le\binom{k-1}{a_1-1}\binom{\beta_n-a_2-1}{\beta_n-k-1}\le2^{2\beta_n}$
by~\eqref{long1049:eq:c-explicit}, and summing the at most $\beta_n$ terms
of~\eqref{long1049:eq:A-explicit} gives $\|A_n\|\le\beta_n2^{2\beta_n}$.  Hence
$\|U_n\|\le\|D_N/\Omega_n\|\,\|A_n\|\le e^{O(n^{2})}$ and
$H(U_n)\le e^{O(n^{2})}$.

For $V_n$, bound it on the circle $|z|=2$ and use Cauchy's estimate
$H(V_n)\le\max_{|z|=2}|V_n(z)|$: each coefficient satisfies
$|v_i|\le2^{-i}\max_{|z|=2}|V_n(z)|$. The bound also holds for the
zero polynomial.  On that circle
$|z^{l}-1|\ge2^{l}-1\ge1$, so each of the at most $\beta_n+a_0$ inner terms
of~\eqref{long1049:eq:B-explicit} has modulus at most $1$; also
$|c_{n,k}(z)z^{a_0k}|\le\|c_{n,k}\|\,2^{K_n}$ and
$|D_N(z)/\Omega_n(z)|\le\|D_N/\Omega_n\|\,2^{225n^{2}}\le e^{O(n^{2})}$, while
$|z^{-M_n}|\le1$. There are $O(n)$ outer summands, each with $O(n)$
inner terms. After multiplication by the normalising factor, each
term is bounded by $e^{O(n^2)}$, since $K_n=O(n^2)$. Summing the
$O(n^2)$ terms preserves this bound. Hence
$\max_{|z|=2}|V_n(z)|\le e^{O(n^2)}$ and $H(V_n)\le e^{O(n^2)}$.
\end{proof}

\noindent With Lemma~\ref{long1049:res:sourceheight}, Zudilin's family satisfies all four
hypotheses of Theorem~\ref{long1049:res:archcap}, with exact degree limit $d=C_1-C_0$
by~\eqref{long1049:eq:degree-limits}, and decay exponent $\sigma=C_0$ because
$\log\Lambda_n(x)=-(K_n-W_n)\log x+o(n^{2})$ for each fixed real $x>1$ by the
size estimate of Section~\ref{long1049:sec:integer-forms}.  For it $\sigma/(\sigma+d)=\theta^{*}$ and
$(\sigma+d)/\sigma=\mu$. The stronger bound
$|U_n(p)|\le p^{W_n}\exp(O_p(n))$ from
Corollary~\ref{long1049:cor:rational-base-measure} removes the extra height
term at each integer base. Within this family the rational-base
threshold is therefore the reciprocal of the proved integer-base
irrationality-exponent bound.
Lemma~\ref{long1049:res:sourceheight} is kernel-checked in the Lean development.

\paragraph{Examples of rational bases.}
Among reduced $a/b$ with $1\le b<a\le60$, the region of
Theorem~\ref{long1049:res:region} has $137$ members, of which $78$ are
non-integral. The two largest admitted logarithmic ratios belong to
$53/5$ and $31/4$, respectively $0.000313$ and $0.001985$ below
$\theta^*$; the smallest excluded ratio is that of $52/5$, namely
$0.4073243836\ldots$, which exceeds the cutoff by $0.001641$.
The accompanying \texttt{rational\_base\_examples.json} certifies all
$1101$ reduced fractions in the stated range, these rounded margins
and the four fixed-denominator comparisons below, using rational
bounds for logarithms and the cutoff. The
bases in this region and outside the region of \cite{bv1994} are those in the strip
$s^{\mu}<r\le s^{\mu_{\mathrm{BV}}}$, which is empty for $s=2$ and $s=3$, is
$\{31\}$ for $s=4$, and is $\{53,54,56\}$ for $s=5$.  The strip is infinite:
its width $s^{\mu_{\mathrm{BV}}}-s^{\mu}$ eventually exceeds $s$, so for every
large enough denominator it contains an integer congruent to $1$ modulo $s$.
Hence $31/4$ is the member of the strip with least denominator and least numerator.

\paragraph{A certified finite parameter search.}
The parameter ratios $(14,12,14;27)$ are Zudilin's \cite[p.~162]{zudilin2004}.
They maximise $C_0/C_1$ in the following finite box: positive integer
parameters $(\alpha_0,\alpha_1,\alpha_2;\beta)$ with every entry at most
$30$, greatest common divisor $1$, and
\[
 \alpha_1\le\alpha_2,\qquad
 \alpha_1+\alpha_2<\beta\le\alpha_0+\alpha_2.
\]
These are the source's parameter inequalities
\cite[Sec.~5, p.~161]{zudilin2004}, with a bound imposed for enumeration.
The definitions of $C_0,C_1$ and their floor function are reproduced in
Section~\ref{long1049:sec:open}. There are exactly $37{,}533$ tuples.
Exactly two attain the maximum: $(14,12,14;27)$ and $(15,12,13;26)$.
Both give $\theta^*$, although they are distinct primitive directions.

The enumeration and comparisons are certified by rational intervals.
The fractions $j/c$ with $1\le c\le30$ partition $[0,1)$ into $278$
half-open intervals. Every floor function occurring in the search is
constant on each interval, and the first interval contributes zero. On each remaining interval $[u,v)$,
the integral is a sum of positive differences
$(k+u)^{-2}-(k+v)^{-2}$, weighted by $0$ or $1$. With $M=256$, the tail
lies between
\[
 (M+u)^{-1}-(M+v)^{-1}
 \quad\hbox{and}\quad
 (M+u)^{-1}-(M+v)^{-1}+(M+u)^{-2}-(M+v)^{-2}.
\]
This follows by comparing the positive decreasing summand with its
integral. Outward rounding of rational numbers and Machin's bounds for
$\pi$ give, for each of the two maximisers,
\[
 0.40568302137302<C_0/C_1<0.40568302139506.
\]
Every other tuple has $C_0/C_1<0.40563943278333$. The intervals therefore
separate these two candidates from all other tuples, without relying
on floating-point ordering.

Their exact equality is not inferred from overlapping intervals.
Write $\omega_A,\omega_B$ for their step functions in the order just
listed, and let $J_A,J_B$ be their integrals against $d(-\psi_1)$
on $[0,1]$. The floor formulas give
\[
 \omega_B(u)-\omega_A(u)
 =\lfloor13u\rfloor+\lfloor15u\rfloor-2\lfloor14u\rfloor.
\]
For every positive integer $c$, telescoping the $c$ intervals and
splitting the defining trigamma series into residue classes gives
\[
 \int_0^1\lfloor cu\rfloor\,d(-\psi_1(u))
 =\sum_{j=1}^{c-1}\psi_1(j/c)-(c-1)\psi_1(1)
 =\frac{c(c-1)\pi^2}{6}.
\]
Hence $J_B-J_A=\pi^2/3$. Both directions have $m=15$ and
$C_1=1091/2$, while the quadratic part of $C_0$ before the $3/\pi^2$ correction is
$266$ for $A$ and $265$ for $B$. Its decrease by $1$ exactly cancels
the increase $3(J_B-J_A)/\pi^2=1$. Thus $C_0$ also agrees.

The full tuple list, interval bounds and comparison certificates are
in \path{computations/parameter_box.json};
\path{computations/check_parameter_box.py} reproduces them.
A separate implementation, \path{computations/verify_parameter_box.py},
checks every tuple with $512$ tail terms and a different enumeration
order. This proves optimality only in the displayed finite box. No
bound on a direction outside that box follows, and this calculation
is not needed for the irrationality theorem.

The two-coordinate forms in $1$ and $F(p)$ should be distinguished
from simultaneous approximation to two target values. Postelmans and Van Assche prove the $\Q$-linear independence
of $1,\zeta_q(1),\zeta_q(2)$ for $q=1/p$ with integer $p\ge2$
\cite[Thm.~1.3, p.~3]{postelmansvanassche2007}.
Their confluent multiple little $q$-Jacobi construction has two
orthogonality conditions; its common integer normalisation and
nonvanishing argument are in Section~6, especially (6.1)--(6.4).
That theorem treats two target values and retains the integer
inverse-base restriction. It supplies neither the coefficient pairs
nor the all-base estimates assumed in Theorem~\ref{long1049:res:archcap};
it is not an application at $3/2$.

\subsection{Finite calculations}\label{long1049:sec:receipts}

This subsection separates reproduced exact checks from historical
numerical reports. The supplied reconstruction script and coefficient
lists cover $n=1,2,3,4$, and the preceding parameter search has an exact
certificate for every tuple in its finite box. The historical
high-precision evaluations were not rerun. None of these finite
calculations proves a statement for all indices, and none is a Lean proof. The
rank-eight tests and the rational interval certificate are described
separately.

The \href{https://github.com/wcook04/plectis-erdos/blob/39a0a2078c8f9fbe7fc24f912654815cc24685df/research/experiments/erdos1049/README.md}{public reproduction guide} gives the command and dependency pin for
\href{https://github.com/wcook04/plectis-erdos/blob/39a0a2078c8f9fbe7fc24f912654815cc24685df/research/experiments/erdos1049/direction_search.py}{the direction-search program}.
Run it with \texttt{--bound 30} and the output path named there; adding
\texttt{--check} compares the result without replacing the
\href{https://github.com/wcook04/plectis-erdos/blob/39a0a2078c8f9fbe7fc24f912654815cc24685df/research/experiments/erdos1049/receipts/direction-search-bound30.json}{stored search results}.
The stored record specifies mpmath~1.3.0 at thirty decimal digits.  Its exact enumeration
count is $37{,}533$, but its ordering by $C_0/C_1$ is numerical and has no
interval certificate in that historical record. The separate rational
certificate above verifies the maximisers in the stated finite box;
the archived program itself still uses numerical comparisons.

The program \texttt{computations/check\_source\_polynomials.py}
reconstructs the forms of Section~\ref{long1049:sec:source-forms}
directly in $\Z[X]$ from the displayed source formulas. It checks
division by $\Omega_n$ and $X^{M_n}$ with zero remainder and records
all coefficients of $U_n,V_n$ in \texttt{source\_n1.json} through
\texttt{source\_n4.json}. The exact degrees are
\[
\begin{array}{c|rrrr}
 n&1&2&3&4\\\hline
 K_n&587&2264&5032&8891\\
 W_n&333&1315&2944&5220\\
 \deg D_N&72&278&628&1102\\
 \deg\Omega_n&25&94&219&380.
\end{array}
\]
In each case $\deg V_n=W_n-1$ and both leading coefficients are
$(-1)^n$. Independently expanding $F(1/q)=\sum_{j\ge1}\tau(j)q^j$
checks the vanishing coefficients below order $K_n-W_n$ in
$U_n(1/q)F(1/q)-V_n(1/q)$ and its next twelve coefficients against
the positive source expression. Exact rational evaluation at
$31/4$, $3$ and $7/2$ also checks homogenisation; at the two noninteger
bases, the power $b^{W_n-1}$ fails to clear $U_n(a/b)$, as predicted
by the general leading-coefficient argument.

A separate historical numerical report evaluated the remainder identity
at those three bases to relative accuracy below $10^{-39}$ at $n=3$.
It reported positivity of $\widehat\Lambda_n$ and agreement with the
bounds on $H_n$ in Section~\ref{long1049:sec:integer-forms}. These
high-precision real evaluations were not reproduced by the exact
polynomial calculation just described.

The historical constant evaluation used $\omega$, the thirteen intervals
of Lemma~\ref{long1049:res:omega-indicator} and forty-digit trigamma values.
It reported
\[
 J=77.94318447500909\ldots,\qquad C_0=221.30008816500502\ldots,
\]
and $C_1/C_0=2.46497868357497\ldots$, agreeing with the source's printed
precision \cite[p.~162]{zudilin2004}. The rational interval estimate stated earlier, rather than these
rounded digits, supports the numerical inequality.

The main term $K_n\log b-(K_n-W_n)\log a$ at $31/4$ is negative
for every $1\le n\le400$, as verified by the supplied rational logarithm
enclosures and exact totient sums. Dividing by $n^2$ gives, to three
decimal places, $-58.478,-10.280,-4.297,-3.863$ at
$n=1,10,100,400$. This finite sign check does not assert negativity at
every index. The proved limit is
$C_1\log4-C_0\log31=-3.718\ldots$; the elementary convergence-rate bound
is $O(\log^2(n+2)/n)$.
Indeed, summing the $O(y\log(2y))$ endpoint errors over the floor blocks
with $k\le n$ gives $O(n\log^2(n+2))$; the remaining main-term tail is
$O(1)$ since its summands are $O(n^2/k^3)$.  The linear terms of $M_n$
add $O(n)$.  The stronger rate $O(1/n)$ does not follow from these estimates.  Only the limit enters Theorem~\ref{long1049:res:region}.

For orientation, the all-rank order and leading-coefficient formulas of
Theorem~\ref{long1049:res:zudilin-sharp-qorder} give the following first seven values:
\[
\begin{array}{c|rrrrrrr}
N&1&2&3&4&5&6&7\\\hline
\operatorname{ord}&0&1&5&14&30&55&91\\
\operatorname{lc}&1&6&108&4320&324000&40824000&8001504000.
\end{array}
\]


\section{Power comparisons and Hankel determinants}\label{long1049:sec:sharp}

The main Hankel calculation begins in Section~\ref{long1049:sec:hankel-order}
and uses none of the three numerical comparisons preceding it. Those
comparisons give bounds for the later scalar and residue-count tests. The power bracket gives
$\log3/\log2<65/41$ and shows that $65$ is the least
integer exponent $q$ for which $3^{41}<2^q$. Failure of the rank-$41$
selector count one row earlier uses the separate comparison
$2^{129}<3^{82}$. These comparisons settle the stated numerical
inequalities, not the existence of the approximation families to
which one might apply them.

\begin{theorem}[sharp power bracket]\label{long1049:res:powerbracket}
One has
\[
 2^{64}<3^{41}<2^{65}.
\]
Consequently
\[
 \frac{41}{65}<\frac{\log2}{\log3},
 \qquad
 \frac{\log3}{\log2}<\frac{65}{41}.
\]
\end{theorem}
\leannote{Lean: \lproof{ErdosProblems/Erdos1049/PaperFiniteAssembliesR7.lean}{24}{power_bracket}.}

\begin{proof}
Direct evaluation gives
\[
 \begin{aligned}
 2^{64}&=18446744073709551616,\\
 3^{41}&=36472996377170786403,\\
 2^{65}&=36893488147419103232.
 \end{aligned}
\]
Taking logarithms of the upper bound gives $41\log3<65\log2$.
Dividing this inequality by $65\log3$ and by $41\log2$, respectively,
gives the two stated logarithmic bounds; both divisors are positive.
\end{proof}

\noindent The integer sides are checked as
\lword{Erdos1049/AdelicHeightBridge.lean}{42}{threePow_fortyOne_lt_twoPow_sixtyFive}{the upper certificate}
and
\lword{Erdos1049/AdelicHeightBridge.lean}{45}{twoPow_sixtyFour_lt_threePow_fortyOne}{the sharp lower certificate}.

The upper power bound gives $\log2/\log3>41/65$, a sharper lower
bound than $81/200$. In particular,
$41/65-2/5=3/13$. Combining this with
$1/2-1/\pi^2<2/5$ gives the first gap below; the same logarithmic
lower bound is then compared with the rectangular expression.

For $\rho\ge0$ and $\sigma\ge1+\rho$, define
\[
 \Theta_{\mathrm{HP}}(\rho,\sigma)=
 \frac{(1+\rho^2)/2+\sigma-3\sigma^2/\pi^2}
 {(1+\rho)^2/2+\sigma(1+\rho)+(1+\rho^2)/2+\sigma}.
\]
The denominator is positive on this domain.

\begin{corollary}[gaps between the stated logarithmic thresholds]\label{long1049:res:sharpgaps}
For every $\rho,\sigma\in\R$ with $0\le\rho$ and $1+\rho\le\sigma$,
\[
 \frac3{13}<\frac{\log2}{\log3}
   -\left(\frac12-\frac1{\pi^2}\right),
 \qquad
 \frac3{13}<\frac{\log2}{\log3}-\Theta_{\mathrm{HP}}(\rho,\sigma),
\]
where $\Theta_{\mathrm{HP}}$ is the rectangular exponent threshold.  Moreover
\[
 \frac{\log3/\log2-1}{3}<\frac8{41}.
\]
\end{corollary}
\leannote{Lean: \lproof{ErdosProblems/Erdos1049/PaperFiniteAssembliesR7.lean}{34}{height_and_hankel_deficits}.}

\begin{proof}
The first inequality follows from
$41/65-2/5=3/13$, Theorem~\ref{long1049:res:powerbracket}, and
$1/2-1/\pi^2<2/5$. The polynomial calculation in
Section~\ref{long1049:sec:open} proves
$\Theta_{\mathrm{HP}}(\rho,\sigma)\le1/2-1/\pi^2$ on the stated
domain, which gives the second inequality.  The final inequality is a direct
rearrangement of $\log3/\log2<65/41$.
\end{proof}

\noindent Separate formal statements check the
\lword{Erdos1049/AdelicHeightBridge.lean}{86}{threeHalves_rectangular_hp_gap_gt_threeThirteenths}{bound for the rectangular expression}
and the
\lword{Erdos1049/AdelicHeightBridge.lean}{97}{threeHalves_hankelChargeThreshold_lt_eightFortyOne}{logarithmic bound $8/41$}.

The next inequalities compare two proposed savings in the denominator
exponent with the exponent $4N^3-3N^2$ before division. Here $E$ denotes
the exponent saved. The statement assumes the bounds $E\le N^3-N$ or
$E\le2N^3-N$; it does not derive them for a polynomial family. Under either
bound, the saving is less than $39/41$ of the original exponent. Thus
these bounds alone cannot justify the reduction required by this model.

\begin{theorem}[bounds for two proposed degree savings]\label{long1049:res:chargeceilings}
For every integer $N>0$,
\[
 41(N^3-N)<39(4N^3-3N^2).
\]
For every integer $N\ge2$,
\[
 41(2N^3-N)<39(4N^3-3N^2).
\]
Hence the same strict inequalities hold with the left side replaced by
$41E$ whenever, respectively, $E\le N^3-N$ or $E\le2N^3-N$.
\end{theorem}
\leannote{Lean: \lproof{ErdosProblems/Erdos1049/PaperFiniteAssembliesR7.lean}{44}{charge_ceilings}.}

\begin{proof}
After subtraction, the first inequality is
\[
 N(115N^2-117N+41)>0,
\]
whose quadratic factor is $115(N-1)^2+113(N-1)+39$, positive for
$N\ge1$. The second becomes
\[
 N(74N^2-117N+41)>0.
\]
Its quadratic factor is $74(N-2)^2+179(N-2)+103$, positive for
$N\ge2$. The assertions for $E$ follow by monotonicity.
\end{proof}

\noindent The inequalities under an upper bound on the saving are
\lword{Erdos1049/AdelicHeightBridge.lean}{123}{zudilinScalarContent_cannot_meet_required_charge}{the first degree bound}
and
\lword{Erdos1049/AdelicHeightBridge.lean}{153}{zudilinScalarPlusBorder_cannot_meet_required_charge}{the second degree bound}.

Lean checks the displayed integer and real inequalities in
Theorem~\ref{long1049:res:powerbracket},
Corollary~\ref{long1049:res:sharpgaps} and
Theorem~\ref{long1049:res:chargeceilings}. An application must separately
prove that its proposed divisor satisfies one of the stated bounds.
These inequalities do not rule out additional factors caused by
cancellation in a determinant, a different choice of integral forms,
or another proof of irrationality at $3/2$.

\subsection{The sharp \texorpdfstring{$q$}{q}-order of the normalised Hankel
determinant}\label{long1049:sec:hankel-order}

The preceding inequalities concern degrees used to clear denominators.
We now determine a different quantity: the first nonzero term, as a
formal power series in $q$, of Zudilin's normalised Hankel determinant
at $x=z=1$. Its order does not by itself give a divisor of an integer
evaluation. That requires an integral normalisation and a separate
divisibility proof.  Zudilin
builds the determinant by combining the Pad\'e-type approximations of Bundschuh
and Zudilin \cite{bundschuhzudilin2008} with B\'ezivin's method as developed
in \cite{krvz2009}; see \cite[Sec.~2, p.~3]{zudilin2016}.  His passage from the
$q$-order to the size of the determinant borrows the proofs of
\cite[Lemma~2, p.~12, and Prop.~3, p.~13]{krvz2009}, as he notes in
\cite[Sec.~4, p.~7]{zudilin2016}.

The source writes $\ell_p(x,z)=x\sum_{r\ge1}z^r/(p^r-x)$, so
$\ell_p(1,1)=F(p)$. Here $p=q^{-1}$ and $x,z$ are auxiliary
parameters, not the base. Write $N$ for the Hankel rank. With $q$ a
formal variable, the moments and determinant in question are
\[
 v_m^*=\sum_{t\ge0}q^{(m+1)t}
       \frac{(q;q)_m^3(q^{t+1};q)_m}{(q^{m+t+1};q)_{m+1}},
 \qquad V_N^*=\det(v_{i+j}^*)_{0\le i,j<N}.
\]
All product denominators have constant term $1$, so their inverses
exist in $\Z[[q]]$. The order of a nonzero series is the least
exponent of $q$ with a nonzero coefficient.

\begin{theorem}[the first nonzero term at every rank]\label{long1049:res:zudilin-sharp-qorder}
For every rank $N$, the normalised Hankel determinant $V_N^{*}$ of
\cite[Sec.~4, (6), p.~6]{zudilin2016}, evaluated at $x=z=1$, has
\[
 \operatorname{ord}_q V_N^{*}=\frac{N(N-1)(2N-1)}{6},
 \qquad
 \text{leading coefficient}\quad\frac{(N!)^{2}(N+1)!}{2^{N}} .
\]
\end{theorem}
\leannote{Lean: \lproof{ErdosProblems/Erdos1049/AllRow/Producer.lean}{173}{order_zudilinNormalizedHankelDet_all}, \lproof{ErdosProblems/Erdos1049/AllRow/Producer.lean}{199}{coeff_zudilinNormalizedHankelDet_all_rat}.}

\noindent The source proves the order inequality and also gives
$|V_N^*(q)|\le |q|^{N^3/3}\exp(O(N^2))$, referring to additional estimates
for the passage from formal order to analytic size
\cite[Sec.~4, p.~7]{zudilin2016}.
Equality in the order bound rules out a further initial power of $q$ at
$x=z=1$; it does not itself prove a fixed-$q$ estimate.
For $0<q<1$, the positive-measure argument below gives both positivity
and a two-sided comparison with the leading term, with logarithmic error
$O_q(N)$. It refines the subleading control, not the cubic exponent of $q$.

\paragraph{The row identity.}
We choose row operations whose first surviving coefficients can be
computed explicitly. The auxiliary lemma applies to more general
products than the moments above; its use here is to identify the
coefficient that remains after each row operation.

Work over $A=\Z[[q]]$. Let $H(X)=1+\sum_{s\ge1}a_s(q)X^s$.
For nonnegative integers $m,t$, put
\[
 W_m(t)=q^{(m+1)t}\prod_{r=1}^{m}H(q^{r}),
 \qquad
 D_j=\prod_{r=0}^{j-1}(I-q^{r}\mathcal N),
 \qquad
 E(m,j)=mj-\frac{j(j-1)}2,
\]
where $\mathcal N$ is the backward shift $(\mathcal Nf)_m=f_{m-1}$ in the index
$m$. Thus $D_j$ is Zudilin's backward-difference operator in
\cite[Sec.~4, (7), p.~6]{zudilin2016}.  Write $\bar H=H\bmod q$ and
$h_r=[X^{r}]\bar H(X)^{-1}$, with $h_r=0$ for $r<0$ and $h_0=1$.

The lemma allows arbitrary coefficients $a_s(q)\in\Z[[q]]$; the only
normalisation imposed on $H$ is its constant term $1$. It therefore
covers the product series used below, but not an unnormalised series
with a different constant term. All statements here are coefficientwise
identities of formal series, with no analytic convergence assumption.
For a general $H$, the coefficient in the lemma can vanish. The
application below computes it for the chosen products and proves the
required nonvanishing.

\begin{lemma}[a coefficient of each transformed row]\label{long1049:res:allrowinitial}
For $m\ge j\ge0$ one has $D_jW_m(t)\in q^{E(m,j)}A$ and
\[
 \bigl[q^{E(m,j)}\bigr]D_jW_m(t)=(-1)^{j}h_{j-t}.
\]
\end{lemma}
\leannote{Lean: \lproof{ErdosProblems/Erdos1049/PaperCompleteR21/AllRowInitialCoefficient.lean}{213}{all_row_initial}, \lproof{ErdosProblems/Erdos1049/PaperCompleteR21/AllRowInitialCoefficient.lean}{242}{all_row_initial_reciprocal}, \lproof{ErdosProblems/Erdos1049/PaperCompleteR21/AllRowInitialCoefficient.lean}{255}{all_row_initial_dvd}, \lproof{ErdosProblems/Erdos1049/PaperCompleteR21/AllRowInitialCoefficient.lean}{116}{paperE_eq_rowExponent}, and 3 further declarations in the \hyperref[long1049:sec:coverage]{coverage section}.}

\begin{proof}
Let $e_u$ denote a formal basis vector indexed by $u\ge0$, and suppress
the argument $q$ in $a_s(q)$. We use the transition rule
\[
 \Phi e_0=\sum_{s\ge1}a_se_{s-1},
 \qquad
 \Phi e_u=(q^{u}-1)e_{u-1}+q^{u}\sum_{s\ge1}a_se_{u+s-1}\quad(u>0).
\]
In $\Phi^je_t$, a coefficient means the sum of the weights of length-$j$
paths from $t$ to the specified endpoint. Each step lowers its index by
at most one. A path ending at $u$ therefore visits only indices at most
$\max(t,u+j)$, so each coefficient is a finite sum. No action on the
algebraic direct sum is assumed.

For $m\ge1$, a direct expansion gives
\[
 W_m(u)-W_{m-1}(u)=q^m\sum_v(\Phi e_u)_vW_{m-1}(v),
\]
since both sides equal
\[
 q^{mu}\prod_{r=1}^{m-1}H(q^r)\bigl(q^uH(q^m)-1\bigr).
\]
For the induction step take $m\ge j+1$ and note that
$E(m,j)-E(m-1,j)=j$. Applying $I-q^j\mathcal N$ to $q^{E(m,j)}W_{m-j}(u)$ therefore gives
$q^{E(m,j)}(W_{m-j}(u)-W_{m-j-1}(u))$.
The preceding identity contributes one further factor $q^{m-j}$, and
$E(m,j)+m-j=E(m,j+1)$. Induction, starting with $D_0=I$, gives
\begin{equation}\label{long1049:eq:state-expansion}
 D_jW_m(t)=q^{E(m,j)}\sum_u(\Phi^{j}e_t)_uW_{m-j}(u).
\end{equation}
The endpoint sum is also locally finite:
$\operatorname{ord}W_{m-j}(u)=(m-j+1)u$, so only finitely many endpoints
contribute to any fixed power of $q$.

Modulo $q$ the operator simplifies: $\bar\Phi e_u=-e_{u-1}$ for $u>0$, and
$\bar\Phi e_0=\sum_{s\ge1}a_s(0)e_{s-1}$.  Put $c_{j,t}=(\bar\Phi^{j}e_t)_0$.
Then $c_{j+1,t}=-c_{j,t-1}$ for $t>0$ and
$c_{j+1,0}=\sum_{s=1}^{j+1}a_s(0)c_{j,s-1}$, the sum terminating because a state
above $j$ cannot reach $0$ in $j$ steps.  Now $c_{0,t}=\delta_{t,0}=h_{-t}$, and
if $c_{j,t}=(-1)^{j}h_{j-t}$ for all $t$ then $c_{j+1,t}=(-1)^{j+1}h_{j+1-t}$ for
$t>0$ at once, while for $t=0$ the identity $\bar H\bar H^{-1}=1$ gives
$h_{j+1}=-\sum_{s\ge1}a_s(0)h_{j+1-s}$ and hence
$c_{j+1,0}=(-1)^{j}\sum_{s\ge1}a_s(0)h_{j+1-s}=(-1)^{j+1}h_{j+1}$.  So
$c_{j,t}=(-1)^{j}h_{j-t}$ for all $j,t$.  Since
$\operatorname{ord}W_{m-j}(u)=(m-j+1)u$ and $m\ge j$, in
\eqref{long1049:eq:state-expansion} only the state $u=0$ contributes at degree $E(m,j)$,
which gives both assertions.
\end{proof}

\noindent Apply the lemma to a summand of the normalised remainder:
\[
 T_{m,t}=q^{(m+1)t}\,(q;q)_m^{3}\,\frac{(q^{t+1};q)_m}{(q^{m+t+1};q)_{m+1}} .
\]
Set
\[
 H_t(X)=\frac{(1-X)^{3}(1-q^{t}X)^{2}}{(1-q^{t}X^{2})(1-q^{t+1}X^{2})}.
\]
Then $T_{m,t}=(1-q^{t+1})^{-1}W_m^{H_t}(t)$, by the telescoping identity
\[
 \prod_{r=1}^{m}H_t(q^r)=
 \frac{(q;q)_m^3(1-q^{t+1})(q^{t+1};q)_m}{(q^{m+t+1};q)_{m+1}}.
\]
The two denominator products contribute the consecutive factors
$1-q^{t+2},\ldots,1-q^{t+2m+1}$.
Reducing modulo $q$ gives $\bar H_t=(1-X)^{3}$ for $t>0$, so
$h^{(t)}_r=\binom{r+2}{2}$, and $\bar H_0=(1-X)^{4}/(1+X)$, so
$h^{(0)}_r=(r+1)(r+2)(2r+3)/6$. The scalar factor
$(1-q^{t+1})^{-1}$ has constant term $1$. The rows are $v_m^*=\sum_{t\ge0}T_{m,t}$. In $D_j$ only the shifts
$m,m-1,\ldots,m-j$ occur. Since $m\ge j$ and
$\operatorname{ord}T_{m-r,t}=(m-r+1)t\ge(m-j+1)t$, only finitely many
$t$ can affect any fixed coefficient after any of these shifts.
This justifies interchanging the sum and $D_j$, and then extracting its
coefficient at degree $E(m,j)$. Terms with $t>j$ contribute
$h^{(t)}_{j-t}=0$, so Lemma~\ref{long1049:res:allrowinitial} at
$m=j+\ell$, where $E(j+\ell,j)=j(j+1)/2+j\ell$, gives
\begin{equation}\label{long1049:eq:row-initial}
 D_jv_{j+\ell}^*=(-1)^{j}\frac{(j+1)^{2}(j+2)}2\,q^{\,j(j+1)/2+j\ell}
 +O\bigl(q^{\,j(j+1)/2+j\ell+1}\bigr)
 \qquad(j,\ell\ge0).
\end{equation}
The coefficient is obtained by summing
\[
 h^{(0)}_j+\sum_{t=1}^{j}h^{(t)}_{j-t}
 =\frac{(j+1)(j+2)(2j+3)}6+\binom{j+2}3
 =\frac{(j+1)^{2}(j+2)}2 .
\]

\begin{proof}[Proof of Theorem~\ref{long1049:res:zudilin-sharp-qorder}]
The operators $D_j$ act by lower unitriangular row operations, so they leave
$\det(v_{i+j}^*)_{0\le i,j<N}$ unchanged.  By~\eqref{long1049:eq:row-initial} the entry in
row $j$ and column $\ell$ has order $j(j+1)/2+j\ell$. For $N=2$,
the matrix of entry orders is
\[
 \begin{pmatrix}0&0\\1&2\end{pmatrix}.
\]
The off-diagonal product is the unique term of order one. Its
permutation sign cancels the negative leading coefficient of the second
row, giving $6q+O(q^2)$. At arbitrary rank the same argument selects
the reversed permutation. For a permutation $\varsigma$, the corresponding
Leibniz term has weight $\sum_j\bigl(j(j+1)/2+j\varsigma(j)\bigr)$.
By the rearrangement inequality, $\sum_jj\varsigma(j)$ is uniquely
minimised by the reversal
$\varsigma(j)=N-1-j$, the values $j$ being distinct. The minimum weight is
$\sum_{j<N}j^{2}=N(N-1)(2N-1)/6$, so exactly one Leibniz term attains it and no
cancellation is possible there. The sign of the reversal is
$(-1)^{N(N-1)/2}$, which cancels $\prod_{j<N}(-1)^{j}$, and the surviving
coefficient is
\[
 \prod_{j=0}^{N-1}\frac{(j+1)^{2}(j+2)}2=\frac{(N!)^{2}(N+1)!}{2^{N}} .
\]
\end{proof}

\paragraph{Formal order and analytic size are different questions.}
Theorem~\ref{long1049:res:zudilin-sharp-qorder} fixes the first nonzero power of $q$ and
its coefficient at each fixed rank.  It does not control the value at a fixed
rational $q$ as the rank grows.  For the rest of this subsection, write $B_N=\sum_{j<N}j^2$ and
$C_N=(N!)^2(N+1)!/2^N$ for the order and leading coefficient.
The integer polynomials $f_N(q)=C_Nq^{B_N}(1-q)^{N^3}$ have exactly
these same two quantities, whereas $f_N(2/3)=C_N(2/3)^{B_N}3^{-N^{3}}$ carries a further
cubic exponential factor that neither datum sees.  The following separate
positive-measure argument supplies the fixed-base estimate for $V_N^*$;
it is not inferred from the formal order.

\paragraph{A separate positive-measure estimate.}
We seek positive moment weights comparable to $(k+1)^2(k+2)/2$,
with constants depending only on $q$. In the determinant expansion
these constants give factors exponential in $N$, so they do not
change its cubic exponent of $q$.
For fixed $0<q<1$ write $P=(q;q)_\infty$, $Q=(\sqrt q;q)_\infty$,
$T=(-1;q)_\infty^2$, and
\[
 G_q(w)=\frac1{(w;q)_\infty^3}\sum_{t\ge0}\frac{w^t}{(q;q)_t}
          \frac{(q^tw^2;q)_\infty}{(q^tw;q)_\infty^2},
 \qquad \gamma_k=[w^k]G_q(w),\qquad c_k=\frac{(k+1)^2(k+2)}2.
\]
The choice $w=q^{m+1}$ collects all dependence on the moment index.
Indeed, the three finite products in the $t$th remainder summand become
\[
 (q;q)_m=\frac{P}{(w;q)_\infty},\quad
 (q^{t+1};q)_m=\frac{P}{(q;q)_t(q^tw;q)_\infty},\quad
 (q^{m+t+1};q)_{m+1}
 =\frac{(q^tw;q)_\infty}{(q^tw^2;q)_\infty}.
\]
Their substitution gives $v_m^*=P^4G_q(q^{m+1})$ term by term.
For fixed $q$ and $0<r<1$, the $t$th summand on $|w|\le r$ satisfies
\[
 \left|\frac{w^t(q^tw^2;q)_\infty}
 {(q;q)_t(w;q)_\infty^3(q^tw;q)_\infty^2}\right|
 \le \frac{(-r^2;q)_\infty}{P(r;q)_\infty^5}\,r^t.
\]
Indeed, $(q;q)_t\ge P$, each denominator product in $w$ has modulus
at least $(r;q)_\infty>0$, and the numerator product has modulus at
most $(-r^2;q)_\infty$. The bound is independent of $t$, so the sum
converges uniformly and absolutely on each such disk. The individual
products converge there as well, since their tails are bounded by
geometric series in $q$. Thus $G_q$ is holomorphic for $|w|<1$,
and its Taylor series may be evaluated at $w=q^{m+1}$ to obtain the
moment expansion. This analytic argument is separate from the formal
substitution in the short note.
The $q$-binomial theorem \cite[Eq.~17.2.37]{dlmf}
\[
 \frac{(Aw;q)_\infty}{(w;q)_\infty}
   =\sum_{j\ge0}\frac{(A;q)_j}{(q;q)_j}w^j\quad(0\le A\le1)
\]
follows by comparing coefficients in
$(1-w)R(w)=(1-Aw)R(qw)$ with $R(0)=1$; the series converges for $|w|<1$
since its coefficients are at most $P^{-1}$.  They are nonnegative, and
at least $(A;q)_\infty$ if $A<1$.  For $A=1$ the series equals $1$.
Factor the numerator of the $t$th term using
\[
 (q^tw^2;q)_\infty=(q^{t/2}w;q)_\infty(-q^{t/2}w;q)_\infty
 (q^{(t+1)/2}w;q)_\infty(-q^{(t+1)/2}w;q)_\infty.
\]
After pairing each positive-argument factor with one denominator,
the $t$th summand of $G_q$ becomes
\[
\begin{aligned}
 &\frac{w^t}{(q;q)_t}
  \frac{(q^{t/2}w;q)_\infty}{(w;q)_\infty}
  \frac{(q^{(t+1)/2}w;q)_\infty}{(w;q)_\infty}\\
 &\qquad{}\times
  \frac{(-q^{t/2}w;q)_\infty(-q^{(t+1)/2}w;q)_\infty}
       {(w;q)_\infty(q^tw;q)_\infty^2}.
\end{aligned}
\]
The $q$-binomial identity makes the two ratios nonnegative
coefficientwise; the remaining product also has nonnegative coefficients.
When $t=0$, the first ratio is $1$ and every coefficient of the second
is at least $Q$. The last fraction has coefficients at least those
of $(1-w)^{-3}$. Thus the $t=0$ term alone bounds $\gamma_k$ below by
$Q\binom{k+3}3$.

For the upper bound, if $R$ has nonnegative coefficients and
$R(1)\le C<\infty$, convolution with a nondecreasing sequence $d_k$
is bounded by $Cd_k$. At $t=0$, bound the second ratio coefficientwise
by $P^{-1}(1-w)^{-1}$. After extracting $(1-w)^{-3}$ from the last
fraction, its remaining factor has value at $w=1$ at most $TP^{-3}$.
This gives $TP^{-4}\binom{k+3}3$.
For $t\ge1$, both ratios are bounded by $P^{-1}(1-w)^{-1}$.
Extracting the one factor $(1-w)^{-1}$ from the last fraction leaves
a factor with value at $1$ at most $TP^{-3}$; also $(q;q)_t^{-1}\le P^{-1}$.
The $t$th summand is therefore bounded coefficientwise by
$TP^{-6}w^t(1-w)^{-3}$. Summing $t\ge1$ gives
$TP^{-6}\binom{k+2}3$ as the bound for its $k$th coefficient. Consequently
\[
 \frac Q3 c_k\le\gamma_k\le T(P^{-4}+P^{-6})c_k.
\]
Thus $\sum_{k\ge0}P^4\gamma_kq^k\delta_{q^k}$ is a finite positive measure
on $[0,1]$, with infinitely many distinct support points and moments
$v_m^*(q)$. In standard terminology, $(v_m^*(q))_{m\ge0}$ is a
Hausdorff moment sequence. This statement fixes $q$; the measure is
not a representing measure for the coefficient sequence $s_m(p)$
considered in the next subsection.
A nonzero polynomial of degree less than $N$ cannot vanish at all of its
first $N$ atoms; the associated Gram matrix is positive definite.
Truncate the measure to its first $K+1$ atoms. At fixed rank $N$, each
matrix entry converges as $K\to\infty$, and the determinant converges
because it is a polynomial in those entries. Finite Cauchy--Binet and
monotone convergence of the nonnegative tuple sums therefore give
Heine's expansion, whose integral form is reproved in
\cite[Sec.~2, (2)--(5), pp.~2--3]{zudilin2017det}:
\[
 V_N^*=\sum_{k_0<\cdots<k_{N-1}}
 \prod_i(P^4\gamma_{k_i}q^{k_i})\prod_{i<j}(q^{k_i}-q^{k_j})^2.
\]
Retaining the tuple $k_i=i$ and using
$\prod_{d=1}^{N-1}(1-q^d)^{2(N-d)}\ge P^{2N}$ gives the lower bound
below. For the upper bound put $k_i=i+\lambda_i$, where the
$\lambda_i$ are nonnegative and nondecreasing. Factoring the smaller
power from each Vandermonde difference gives the exact exponent
\[
 \sum_i k_i+2\sum_{i<j}k_i
 =B_N+\sum_{i=0}^{N-1}(2N-1-2i)\lambda_i.
\]
Every weight $2N-1-2i$ is at least $1$. Since $0<q<1$, the remaining
power of $q$ is at most $q^{\sum_i\lambda_i}$. Also
$c_{i+\lambda}/c_i\le(\lambda+1)^3$. Discard the remaining
Vandermonde factors, each bounded by $1$, and enlarge the nonnegative
sum by dropping the ordering restriction on the $\lambda_i$.
It now factors into $N$ copies of $\sum_{\lambda\ge0}(\lambda+1)^3q^\lambda$,
which gives
\[
 (P^6Q/3)^N C_Nq^{B_N}\le V_N^*(q)
 \le\left[P^4T(P^{-4}+P^{-6})
          \frac{1+4q+q^2}{(1-q)^4}\right]^N C_Nq^{B_N}.
\]
Both constants are positive and finite, so $V_N^*(q)>0$ and
$\log(V_N^*(q)/(C_Nq^{B_N}))=O_q(N)$, including $N=0$ under the empty
determinant convention.
Here $N\to\infty$ with $q$ fixed; the constants are not uniform as
$q\uparrow1$. The formal-order calculation instead fixes $N$ and expands
at $q=0$. No joint uniform limit is asserted. The measure depends on $q$,
not on the moment index or rank, and supplies no denominator factor for
the 2004 polynomial forms.
Section~\ref{long1049:sec:sharp-fixed-base} replaces the band $O_q(N)$ by an
exact asymptotic with an explicit positive constant, using the same weights.

\paragraph{Formal verification and computations.}
The positive-measure argument, including its generating-function
identification and infinite determinant expansion, is an ordinary proof,
not a Lean result. The separate formal-series argument has a checked
\lword{Erdos1049/AdelicHeightBridge.lean}{183}{hankelAssociatedCoeff_eq_reciprocal}{recurrence for the leading coefficients}.
The constructed determinant also has all-rank formal proofs of its
order and leading coefficient, recorded together in
\href{https://github.com/wcook04/plectis-erdos/blob/0b500c7cf8e8bb7ae343484378df02f277fb8194/lean/ErdosProblems/Erdos1049/AllRow/Producer.lean#L163}{the exact-order theorem},
and the two closed forms $6\operatorname{ord}=N(N-1)(2N-1)$ and
$2^{N}\mathrm{lc}=(N!)^{2}(N+1)!$ are
\href{https://github.com/wcook04/plectis-erdos/blob/0b500c7cf8e8bb7ae343484378df02f277fb8194/lean/ErdosProblems/Erdos1049/AllRow/Producer.lean#L173}{the order formula} and
\href{https://github.com/wcook04/plectis-erdos/blob/0b500c7cf8e8bb7ae343484378df02f277fb8194/lean/ErdosProblems/Erdos1049/AllRow/Producer.lean#L185}{the leading-coefficient formula}.
The Lean proof uses the first nonzero term of every transformed row, which
for row $1$ has order exactly $l+1$ and coefficient exactly $-6$ in every
column.  Independently of the proof, the order and the
leading coefficient were computed exactly for $1\le N\le7$, giving orders
$0,1,5,14,30,55,91$ and leading coefficients
$1,6,108,4320,324000,40824000,8001504000$, each matching the closed forms.

\paragraph{Earlier work and the unresolved value at $3/2$.}
The antecedent is the inequality of \cite[Sec.~4, p.~7]{zudilin2016}; the equality and
the leading coefficient are proved here.
Lemma~\ref{long1049:res:allrowinitial} and~\eqref{long1049:eq:row-initial}
give that row calculation, which is formalised in
\href{https://github.com/wcook04/plectis-erdos/blob/0b500c7cf8e8bb7ae343484378df02f277fb8194/lean/ErdosProblems/Erdos1049/AllRow/Producer.lean#L143}{the first nonzero term of each transformed row}.  Nothing in
this subsection decides the arithmetic nature of $F(3/2)$.


\subsection{The size of $V_N^*$ at a fixed base}
\label{long1049:sec:sharp-fixed-base}

The positive-measure argument traps $\log(V_N^*(q)/(C_Nq^{B_N}))$ between two
multiples of $N$. It leaves the exponential rate inside that band undetermined,
and says nothing about the correction below it. Both are determined here.
Throughout this subsection $0<q<1$ is fixed and
\[
 P=(q;q)_\infty,\qquad
 \mathcal M(q)=\prod_{d\ge1}(1-q^d)^{-d},\qquad
 L=F(1/q)=\sum_{r\ge1}\frac{q^r}{1-q^r},
\]
with $B_N=\sum_{j<N}j^2$ and $C_N=\prod_{k<N}c_k=(N!)^2(N+1)!/2^N$ as above.
Both products converge to positive finite values, and the abbreviation $L$ is
used only in this subsection; it is the value at $q$ of the divisor series
$\mathcal L$ of Section~\ref{long1049:sec:open}. Everything in this subsection
is an ordinary proof. None of it is formalised in Lean, and none of it changes
a denominator valuation or supplies a divisor for the 2004 polynomial forms.

The first step isolates what a Hankel determinant needs from its weights. The
hypotheses are two ratio conditions on the weights, so the theorem covers
other moment families of the same geometric shape as well.

\begin{theorem}[determinants of geometric moments]
\label{long1049:thm:geometric-universality}
Let $a_k>0$ satisfy, for fixed constants $C$ and $\kappa$,
\[
 \frac{a_{k+h}}{a_k}\longrightarrow1\quad(k\to\infty)
 \text{ for each fixed }h\ge0,
 \qquad
 \frac{a_{k+h}}{a_k}\le C(1+h)^\kappa\quad(k,h\ge0).
\]
Put $M_m=\sum_{k\ge0}a_kq^{(m+1)k}$ and $D_N=\det(M_{i+j})_{0\le i,j<N}$. Then
\[
 D_N\sim\mathcal M(q)^3\,q^{B_N}P^{2N}\prod_{k=0}^{N-1}a_k
 \qquad(N\to\infty).
\]
\end{theorem}

\begin{proof}
Taking $k=0$ in the ratio bound gives $a_h\le Ca_0(1+h)^\kappa$, so every moment
converges. Heine's expansion, as in
Section~\ref{long1049:sec:sharp}, writes
\[
 D_N=\sum_{k_0<\cdots<k_{N-1}}
 \prod_{i<N}\bigl(a_{k_i}q^{k_i}\bigr)
 \prod_{i<j<N}(q^{k_i}-q^{k_j})^2,
\]
a sum of positive terms. Write $k_i=i+\lambda_i$ with
$0\le\lambda_0\le\cdots\le\lambda_{N-1}$. The tuple $\lambda=0$ contributes
\[
 q^{B_N}\prod_{i<N}a_i\cdot\Delta_N,
 \qquad
 \Delta_N=\prod_{d=1}^{N-1}(1-q^d)^{2(N-d)}.
\]
Reverse the shifts by setting $\mu_j=\lambda_{N-j}$ for $1\le j\le N$, so that
$\mu_1\ge\mu_2\ge\cdots\ge\mu_N\ge0$ is a partition. Dividing the general
summand by the displayed contribution leaves
\[
 W_N(\mu)=
 q^{\sum_{j\le N}(2j-1)\mu_j}
 \prod_{j\le N}\frac{a_{N-j+\mu_j}}{a_{N-j}}
 \prod_{1\le j<k\le N}
 \left(\frac{1-q^{k-j+\mu_j-\mu_k}}{1-q^{k-j}}\right)^2,
\]
because the summand indexed by $(k_i)$ has $q$-exponent
$B_N+\sum_i(2N-1-2i)\lambda_i$ and the weight $2N-1-2i$ becomes $2j-1$ under
$i=N-j$. Thus $D_N/(q^{B_N}\Delta_N\prod_{i<N}a_i)=\sum_\mu W_N(\mu)$.

Let $\ell$ be the length of $\mu$, that is, the number of positive parts. Only
the indices $j\le\ell$ contribute a ratio different from $1$. For fixed
$j\le\ell$ the partition is decreasing, so $k-j+\mu_j-\mu_k\ge k-j$ for every
$k>j$; the numerators are at most $1$ and the denominators multiply to at least
$P$. Each of the $\ell$ squared factors is therefore at most $P^{-2}$, and the
weight ratios contribute at most $\prod_{j\le\ell}C(1+\mu_j)^\kappa$. Hence
\[
 W_N(\mu)\le(CP^{-2})^{\ell}
 \prod_{j=1}^{\ell}(1+\mu_j)^\kappa q^{(2j-1)\mu_j}
 \qquad\text{for every }N.
\]
Summing this majorant over all partitions of length exactly $\ell$, after
dropping the ordering constraint on the positive parts, factorises it into
\[
 (CP^{-2})^{\ell}\prod_{j=1}^{\ell}
 \sum_{u\ge1}(1+u)^\kappa q^{(2j-1)u}
 \le C_q^{\ell}q^{\ell^2},
 \qquad
 C_q=CP^{-2}\sum_{u\ge1}(1+u)^\kappa q^{u-1},
\]
using $q^{(2j-1)u}\le q^{2j-1}q^{u-1}$ and $\sum_{j\le\ell}(2j-1)=\ell^2$. The
resulting bound is summable in $\ell$, and it does not depend on $N$.

For each fixed partition $\mu$, the weight ratios tend to $1$ as $N\to\infty$
by the first hypothesis, and the Vandermonde quotient converges to the
corresponding infinite product. The limit of $W_N(\mu)$ is therefore the same
function of $\mu$ that the weights $a_k\equiv1$ produce, since those weights
give ratio $1$ identically and leave the Vandermonde quotient unchanged.
Dominated convergence over partitions now identifies
$\lim_N\sum_\mu W_N(\mu)$ with the corresponding limit in the unit-weight case.

For $a_k\equiv1$ the moments are $M_m=(1-q^{m+1})^{-1}$, and Cauchy's
determinant evaluates
\[
 D_N=\frac{q^{B_N}\Delta_N}{\prod_{0\le i,j<N}(1-q^{i+j+1})}.
\]
The exponent of $1-q^m$ in that denominator is $\min(m,N,2N-m)$, which
increases to $m$, so the denominator tends to $\mathcal M(q)^{-1}$ and the
quotient $\sum_\mu W_N(\mu)$ tends to $\mathcal M(q)$. Finally
\[
 \frac{\Delta_N}{P^{2N}}
 =\prod_{d<N}(1-q^d)^{-2d}\prod_{d\ge N}(1-q^d)^{-2N}
 \longrightarrow\mathcal M(q)^2,
\]
because the logarithm of the second product is $O_q(Nq^N)$. Multiplying the two
limits proves the theorem.
\end{proof}

The hypotheses hold for the actual weights $a_k=P^4\gamma_k$, and the reason is
an exact algebraic identity. It also sharpens the comparison
$(Q/3)c_k\le\gamma_k\le T(P^{-4}+P^{-6})c_k$ of the previous subsection to two
explicit powers of $P$.

\begin{proposition}[the moment weights as a product of two finite sums]
\label{long1049:prop:rogers-factorisation}
For $r\ge1$ and $k\ge0$ set
\[
 R_k^{(r)}(q)=\sum_{n_1+\cdots+n_r=k}
 \frac{(q;q)_k}{\prod_{j\le r}(q;q)_{n_j}},
\]
the sum being over compositions into $r$ nonnegative parts, so that
$R_k^{(r)}$ is the sum of all $q$-multinomial coefficients of degree $k$ and
length $r$, equivalently the multivariate Rogers--Szeg\H{o} polynomial at unit
arguments~\cite{vinroot2010}. These are polynomial identities in the
indeterminate $q$; at a prime power, $r=2$ gives the Galois number counting
the subspaces of a $k$-dimensional space over the field of that order, and
$r=3$ counts the flags $V_1\subseteq V_2$ in that space. The source is cited
for those identifications; the factorisation below is proved here. Then
\[
 \gamma_k(q)=\frac{R_k^{(2)}(q)R_k^{(3)}(q)}{(q;q)_k}.
\]
The product $R_k^{(2)}R_k^{(3)}$ lies in $\Z_{\ge0}[q]$, has degree
$\lfloor k^2/4\rfloor+\lfloor k^2/3\rfloor$, and has coefficient sum $6^k$.
Consequently $(q;q)_k\gamma_k\in\Z[q]$, and the weights
$a_k=P^4\gamma_k$ satisfy
\[
 P^5c_k\le a_k\le P^{-1}c_k
 \qquad(k\ge0),
\]
together with $a_{k+h}/a_k\le P^{-6}(1+h)^3$, so the hypotheses of
Theorem~\ref{long1049:thm:geometric-universality} hold with $\kappa=3$.
\end{proposition}

\begin{proof}
Write $\mathcal E(z)=(z;q)_\infty^{-1}$ and
$\partial_qf(z)=(f(z)-f(qz))/z$, so that $\partial_q\mathcal E=\mathcal E$ and
$\mathcal E(z)^r=\sum_{n\ge0}R_n^{(r)}z^n/(q;q)_n$ by Euler's expansion. The
$q$-Leibniz rule, which follows by induction from the product rule, is
\[
 \partial_q^n(fg)(z)=\sum_{j=0}^n\genfrac{[}{]}{0pt}{}{n}{j}_q
 (\partial_q^jf)(q^{n-j}z)\,(\partial_q^{n-j}g)(z).
\]
Write $h_k(z)=\sum_{j=0}^k\genfrac{[}{]}{0pt}{}{k}{j}_q(z;q)_j$. Applying the
rule to $\mathcal E\cdot\mathcal E$ and using
$\mathcal E(q^mz)=(z;q)_m\mathcal E(z)$ gives
$\partial_q^k\mathcal E(z)^2=\mathcal E(z)^2h_k(z)$; at $k=1$ both sides are
$\mathcal E(z)^2(2-z)$. Applying it once more, to
$\mathcal E\cdot\mathcal E^2$, gives
\[
 \partial_q^n\mathcal E(z)^3
 =\mathcal E(z)^3\sum_{k=0}^n\genfrac{[}{]}{0pt}{}{n}{k}_q(z;q)_kh_k(z),
\]
the power of $\mathcal E$ being $3$ and not $2$ here. Hence
\[
 \sum_{n\ge0}\frac{w^n}{(q;q)_n}\partial_q^n\mathcal E(z)^3
 =\mathcal E(w)\mathcal E(z)^3
  \sum_{k\ge0}\frac{w^k(z;q)_k}{(q;q)_k}h_k(z).
\]
Put $z=w$. On the left the coefficient of $w^n$ is
$R_n^{(3)}\sum_{j\le n}\bigl((q;q)_j(q;q)_{n-j}\bigr)^{-1}
=R_n^{(2)}R_n^{(3)}/(q;q)_n$. On the right, expand $h_k$ and set $k=j+t$; the
inner sum is evaluated by the $q$-binomial theorem as
\[
 \sum_{t\ge0}\frac{(wq^j;q)_t}{(q;q)_t}w^t
 =\frac{(w^2q^j;q)_\infty}{(w;q)_\infty},
\]
and the result is exactly $G_q(w)$. Every step holds coefficientwise, and
absolutely for $|w|$ small.

A $q$-multinomial coefficient has nonnegative integral coefficients, degree
$\sum_{i<j}n_in_j$, and value $1$ at $q=1$ after division by the corresponding
multinomial. Summing over compositions gives coefficient sum $r^k$ for
$R_k^{(r)}$ and hence $6^k$ for the product. The degree is maximised by
balanced parts, which gives $\lfloor k^2/4\rfloor$ for $r=2$ and
$\lfloor k^2/3\rfloor$ for $r=3$; no leading cancellation occurs because all
coefficients are nonnegative.

For the bounds, write $b_k^{(r)}=[z^k](z;q)_\infty^{-r}$, so that
$R_k^{(r)}=(q;q)_kb_k^{(r)}$ and
$\gamma_k=(q;q)_kb_k^{(2)}b_k^{(3)}$. Euler's expansion exhibits $b_k^{(r)}$ as
an $r$-fold convolution of the sequence $1/(q;q)_n$, whose terms lie in
$[1,P^{-1}]$; the number of compositions of $k$ into $r$ nonnegative parts is
$\binom{k+r-1}{r-1}$, so
\[
 \binom{k+r-1}{r-1}\le b_k^{(r)}\le P^{-r}\binom{k+r-1}{r-1}.
\]
Combining these at $r=2,3$ with $P\le(q;q)_k\le1$ and
$c_k=(k+1)\binom{k+2}2$ gives $P^5c_k\le a_k\le P^{-1}c_k$. The shift bound
follows from $c_{k+h}/c_k\le(1+h)^3$.
\end{proof}

Only one analytic input is still missing, namely the behaviour of $a_k/c_k$ to
second order. It is supplied by the two leading coefficients of $G_q$ at its
singularity $w=1$.

\begin{theorem}[the size of $V_N^*$ at a fixed base]
\label{long1049:thm:sharp-fixed-base}
Let $\gamma_{\!E}$ be Euler's constant. The product
\[
 \mathcal A(q)=e^{-8\gamma_{\!E}L}
 \prod_{k\ge0}\left(\frac{a_k}{c_k}\,e^{8L/(k+1)}\right)
\]
converges to a positive finite value. With $K(q)=\mathcal A(q)\mathcal M(q)^3$,
\[
 V_N^*(q)\sim K(q)\,C_Nq^{B_N}P^{2N}N^{-8L},
\]
equivalently
\[
 \log V_N^*(q)=B_N\log q+\log C_N+2N\log P-8L\log N+\log K(q)+o(1).
\]
\end{theorem}

\begin{proof}
Write $\Pi(w)=(qw;q)_\infty$, so that $(w;q)_\infty=(1-w)\Pi(w)$, and for
$t\ge1$ put
\[
 f_t(w)=\frac{(q^tw^2;q)_\infty}{(q;q)_t(q^tw;q)_\infty^2},
 \qquad
 \mathcal S(w)=\sum_{t\ge1}w^t\bigl(f_t(w)-P^{-1}\bigr).
\]
On each compact subset of $|w|<q^{-1}$ the bracket is $O_q(q^t)$ uniformly, so
$\mathcal S$ is analytic there. Separating the summand $t=0$, whose
denominator carries the factor $(1-w)^2$, and subtracting the constant part of
the remaining summands gives
\[
 \mathcal D(w):=P^4(1-w)^4G_q(w)
 =\frac{P^4}{\Pi(w)^3}
 \left((1+w)\frac{(qw^2;q)_\infty}{\Pi(w)^2}+\frac wP
       +(1-w)\mathcal S(w)\right).
\]
Hence $\mathcal D$ is analytic on a disc of radius greater than $1$.

At $w=1$ one has $\Pi(1)=P$ and $\Pi'(1)/\Pi(1)=-L$, since
$\Pi'/\Pi=-\sum_{j\ge1}q^j/(1-q^jw)$. Also $(qw^2;q)_\infty$ equals $P$ at
$w=1$, and the logarithmic derivative of $(qw^2;q)_\infty/\Pi(w)^2$ at $w=1$ is
$-2L+2L=0$. Therefore the bracket has value $3/P$ and derivative $(2-L)/P$
at $w=1$; here $\mathcal S(1)=L/P$, because
$f_t(1)=\bigl(P(1-q^t)\bigr)^{-1}$ and
$\sum_{t\ge1}\bigl((1-q^t)^{-1}-1\bigr)=L$. Combining with the factor
$P^4\Pi^{-3}$, whose logarithmic derivative at $1$ is $3L$, gives
\[
 \mathcal D(1)=3,\qquad
 \mathcal D'(1)=3\left(3L+\frac{2-L}3\right)=2+8L.
\]
Subtracting the principal part of $\mathcal D(w)(1-w)^{-4}$ at $w=1$ leaves a
function analytic on a larger disc, so
\[
 a_k=3\binom{k+3}3-(2+8L)\binom{k+2}2+O_q(k),
 \qquad
 \frac{a_k}{c_k}=1-\frac{8L}{k+1}+O_q\bigl((k+1)^{-2}\bigr),
\]
the second line because $3\binom{k+3}3=c_k(k+3)/(k+1)$ and
$\binom{k+2}2=c_k/(k+1)$.

Consequently $\log(a_k/c_k)+8L/(k+1)=O_q((k+1)^{-2})$, which is summable, so
$\mathcal A(q)$ converges; it is positive because every factor is. Since
$\sum_{k<N}(k+1)^{-1}=\log N+\gamma_{\!E}+o(1)$,
\[
 \prod_{k<N}a_k\sim\mathcal A(q)\,C_NN^{-8L}.
\]
Proposition~\ref{long1049:prop:rogers-factorisation} supplies the uniform shift
bound and the last expansion supplies the fixed-shift limit, so
Theorem~\ref{long1049:thm:geometric-universality} applies to $(a_k)$ and gives
$V_N^*\sim\mathcal M(q)^3q^{B_N}P^{2N}\prod_{k<N}a_k$. Substituting the product
asymptotic proves both displays.
\end{proof}

\paragraph{What the sharp estimate settles, and what it leaves.}
The positive-measure argument of the previous subsection had recorded this
shape as a candidate without its constant. The proof, the limiting Gram factor
$\mathcal M(q)^3$ and the convergent expression for the constant are supplied
here. The estimate replaces the band $O_q(N)$ by the explicit terms
$2N\log P$, $-8F(1/q)\log N$ and the constant $\log K(q)$, with an $o(1)$
error, and the weights now have an exact finite description rather than an
infinite-product representation alone.

The cubic exponent of $q$ is untouched, and so is every prime-power valuation.
At $3/2$ the deficit is unchanged and is exactly
\[
 \frac{39}{41}\bigl(4N^3-3N^2\bigr)-\bigl(2N^3-N\bigr)
 =\frac{N(74N^2-117N+41)}{41}>0\qquad(N\ge2),
\]
the positivity being immediate from
$74N^2-117N+41=74(N-2)^2+179(N-2)+103$. Its leading size is $(74/41)N^3$, so no
estimate that leaves the cubic coefficient alone can close it. The finite
factorisation describes the same weights in a new way; it does not license
substituting their denominator for the denominator of a cleared linear form.


\subsection{Coefficient moments and cyclotomic content}
\label{long1049:sec:coefficient-questions}
The positive measure just constructed belongs to the remainders $v_m^*$.
It is not a measure for their coefficients in $F(p)$. To make the latter
question precise, fix $p>1$, use Gaussian binomial polynomials, and put
\[
 \begin{aligned}
 R_m(p)&=\sum_{k=0}^m(-1)^{m+k}p^{k(k+1)/2}
          \genfrac{[}{]}{0pt}{}{m}{k}_p\genfrac{[}{]}{0pt}{}{m+k}{k}_p,\\
 s_m(p)&=([m]_p!)^3R_m(p),
 \end{aligned}
\]
where $[m]_p!=\prod_{j=1}^m(1+p+\cdots+p^{j-1})$.
Each $R_m$ is the signed value at $x=p^{m+1}$ of the little
$q$-Legendre polynomial of degree $m$, with $q=p^{-1}$; this
identification is discussed in the related-work section.
The polynomial normalisation of~\cite[Sec.~3, p.~5]{zudilin2016}, at
$x=z=1$, gives
\[
 \begin{aligned}
 \alpha_m&=p[p(p-1)^3]^m s_m(p),\\
 \beta_m&=\left[\alpha_m\sum_{j\ge1}\tau(j)p^{-j}\right]_+-1,\\
 v_m^*&=\alpha_mF(p)-\beta_m.
 \end{aligned}
\]
The brackets mean the polynomial part at infinity. Thus
$\alpha_0=p$ and $\beta_0=0$. This fixes the normalisation before any
content or positivity test.

\paragraph{A positive expansion is not a moment representation.}
Substituting $x=p^{m+1}$ into Van Assche's alternative expansion
\cite[(16), p.~4]{vanassche2001} gives the identity
\[
 R_m(p)=\sum_{k=0}^m\genfrac{[}{]}{0pt}{}{m}{k}_p\genfrac{[}{]}{0pt}{}{m+k}{k}_p
 p^{(m-k)(m-k+1)/2}\prod_{j=m-k+1}^{m}(p^j-1).
\]
Indeed $(qx;q)_k=(-1)^k\prod_{j=m-k+1}^m(p^j-1)$, so the two signs cancel.
This proves $R_m(p)>0$ for $p>1$, and nonnegative coefficients in the
variable $t=p-1$; it does not prove Hankel positivity. In fact,
\[
 \det(R_{i+j}(p))_{i,j<3}=-36(p-1)^2+O((p-1)^3)
 \quad(p\downarrow1).
\]
A geometric factor $c\rho^m$ with $c,\rho>0$ would change a Hankel
matrix only by a positive scalar and an invertible diagonal congruence,
so it would preserve positive definiteness. The factor $([m]_p!)^3$
is not geometric: its values at $m=0,1,2$ are $1,1,(1+p)^3$.
It cannot be discarded in testing the moment condition. Nor does its
positivity alone prove that multiplication by it repairs the failed
condition for $R_m$. At $p=1$, $s_m(1)=(m!)^3$ is the moment sequence of a
product of three independent unit exponential variables, since their
$m$th moments multiply. This is only the polynomial endpoint; $F(1)$
diverges. Berg's Theorem~5.1 states that $(m!)^c$ is Stieltjes indeterminate
for $c>2$ \cite{berg2007}. Thus a measure exists at the endpoint but is not
unique. Neither uniqueness nor a canonical deformation is a premise of
the question at $p>1$.

\paragraph{Finite certificates and an all-rank degree check.}
Put $D_{N,h}(p)=\det(s_{i+j+h}(p))_{0\le i,j<N}$, with $D_{0,h}=1$.
Exact polynomial calculations give strictly positive coefficients in
$t=p-1$ for $D_{N,h}(1+t)$ when $h=0,1$ and $1\le N\le8$.
Thus sixteen determinant polynomials are positive for every real $p\ge1$.
Their degrees, in increasing rank, are
\[
\begin{array}{c|rrrrrrrr}
h=0&0&10&54&156&340&630&1050&1624\\
h=1&2&26&96&236&470&822&1316&1976.
\end{array}
\]
The calculation uses exact arithmetic in $\Z[p]$. Starting from
$D_{0,h}=1$ and $D_{1,h}=s_h$, the Desnanot--Jacobi identity gives
\[
 D_{N,h}=
 \frac{D_{N-1,h}D_{N-1,h+2}-D_{N-1,h+1}^{\,2}}{D_{N-2,h+2}}
 \qquad(N\ge2).
\]
The degree argument below shows that each denominator is nonzero.
The computation checks that every division has zero remainder, then
substitutes $p=1+t$ and tests every coefficient. All $8824$ coefficients
in the sixteen polynomials are strictly positive. The full lists are in
\path{computations/finite_coefficients.json}, computed by
\path{computations/finite_coefficients_generate.py}.
An independent program, \path{computations/check_finite_coefficients.py},
evaluates each stored polynomial at more integer bases than its degree and
compares the values with fraction-free determinants of the integer Hankel
matrices, which identifies the stored list with $D_{N,h}(1+t)$ in $\Z[t]$.
As a separate check of the normalisation, direct integer determinants
at $p=1,2,3,5,11$ agree with evaluations of all sixteen polynomials.
Those evaluations are checks, not the proof of polynomial positivity.
This is a finite computer-algebra calculation and gives no all-rank result.

For the eight unshifted determinants, positivity is also a
\href{\laterepobase/lean/ErdosProblems/Erdos1049/PaperR20/HankelKroneckerCertificate.lean\#L458}{Lean theorem},
proved by a different certificate. Put $T=2^{2929}$. An explicit bound
shows that every coefficient of $D_{N,0}(1+t)$, $N\le8$, has absolute
value less than $T/2$. An integer polynomial with this property is
determined by its value at $t=T$, and its coefficients are the base-$T$
digits of that value whenever every digit is less than $T/2$. The Lean
kernel computes the eight integers $D_{N,0}(1+T)$ exactly and checks this
digit condition and that the lowest digit is nonzero. Every coefficient
of $D_{N,0}(1+t)$ is therefore nonnegative and the constant term is
positive, so $D_{N,0}(p)>0$ for every real $p\ge1$. This is the positivity
used in Proposition~\ref{long1049:res:finite-pencil}. The strict
positivity of every coefficient, and the shifted determinants $D_{N,1}$,
rest on the computation above.

The degrees admit a separate all-rank proof. The Gaussian binomial
$\genfrac{[}{]}{0pt}{}{u}{v}_p$ is monic of degree $v(u-v)$. In the signed sum for $R_m$,
the degree increases by $2m-k$ from the $k$th to the $(k+1)$st term.
The unique maximal term is $k=m$, with positive leading coefficient.
Hence $R_m$ is monic of degree $(3m^2+m)/2$, and $s_m$ is monic of
degree $3m^2-m$. In the determinant, the only permutation-dependent
part of the degree is $6\sum_i i\sigma(i)$; strict rearrangement makes
the identity its unique maximum. Therefore
\[
 \deg D_{N,h}=\sum_{i=0}^{N-1}\bigl(3(2i+h)^2-(2i+h)\bigr),
 \qquad \operatorname{lc}D_{N,h}=1.
\]
More generally the same argument works for a fixed minor with distinct
increasing row and column indices. Each such minor is positive for all
sufficiently large $p$, with a threshold that may depend on the minor.
At $p=1$, both leading Hankel families are positive definite by the
infinite-support product measure, so each fixed rank is also positive in
some neighbourhood of $1$. Neither argument provides a neighbourhood or
large-base threshold uniform over every rank.

A Stieltjes moment representation requires positive semidefiniteness of
both Hankel families at all ranks. The positive-definite formulation
for nondegenerate sequences is stated in
\cite[Lemma~2.1, p.~4]{wangzhu2016}; finite support requires allowing
semidefinite matrices. The necessity of both conditions is visible
from their quadratic forms:
if a positive measure $\mu$ on $[0,\infty)$ satisfies
$s_m=\int_0^\infty x^m\,d\mu(x)$, and
$P(x)=\sum_{i=0}^{N-1}c_ix^i$ has real coefficients, then
\[
 \sum_{i,j=0}^{N-1}c_ic_j s_{i+j+h}
 =\int_0^\infty x^hP(x)^2\,d\mu(x)\ge0,
 \qquad h=0,1.
\]
The shifted test records nonnegative support, not merely positivity
of the measure. For example, the positive point mass at $-1$ has
moments $(-1)^m$. Its unshifted quadratic form is $P(-1)^2$, whereas
its shifted form is $-P(-1)^2$. Thus a Hamburger moment sequence,
which permits support anywhere on $\R$, need not be a Stieltjes
moment sequence. The two tests above do not assert a measure for the
entire coefficient sequence considered here.
Coefficientwise total positivity would be stronger than the sixteen
certificates above. Positive production matrices and path constructions
can supply such a mechanism in other families
\cite[Sec.~2]{lrz2017}\cite[Thms.~9.7--9.8]{psz2023}; no such matrix for
this moving-degree sequence has been established here.

\paragraph{A finite spectral consequence.}
Write $A_N=(\alpha_{i+j})_{0\le i,j<N}$ and
$B_N=(\beta_{i+j})_{0\le i,j<N}$. We study the matrix pencil
$YA_N-B_N$, where $Y$ is a scalar variable, through the roots of its
determinant.
Let $D=\operatorname{diag}(1,p(p-1)^3,\ldots,[p(p-1)^3]^{N-1})$.
Then $A_N=pD(s_{i+j})D$. For $p>1$ this is an invertible positive diagonal
congruence. The remainder representation already proves
$F(p)A_N-B_N=(v^*_{i+j})>0$ at every rank.

\begin{proposition}[finite coefficient positivity and pencil roots]
\label{long1049:res:finite-pencil}
For every real $p>1$ and $1\le N\le8$, $A_N$ is positive definite and
all roots of $\det(YA_N-B_N)$ are real and strictly less than $F(p)$.
The roots at consecutive ranks $N,N+1\le8$ interlace non-strictly.
\end{proposition}
\leannote{Not formalised: the coefficient positivity of the eight unshifted determinants is established by an exact integer computation outside the Lean kernel, and the passage through the eigenvalues and the interlacing of the pencil roots are not formalised. See the \hyperref[long1049:sec:coverage]{coverage section}.}
\begin{proof}
The leading principal minors of $(s_{i+j})_{i,j<N}$ are exactly
$D_{k,0}(p)$ for $1\le k\le N$; their coefficient positivity proved above, together with
Sylvester's criterion, gives positive definiteness. The diagonal congruence gives
$A_N>0$. The real symmetric matrix $A_N^{-1/2}B_NA_N^{-1/2}$ has these
pencil roots as eigenvalues, and $F(p)I-A_N^{-1/2}B_NA_N^{-1/2}>0$ puts
them strictly below $F(p)$. The minimum--maximum principle for
$x^{\mathsf T}B_Nx/(x^{\mathsf T}A_Nx)$ on nested coordinate subspaces
gives non-strict interlacing. These are subspaces for the original
pencils; the separately conjugated symmetric matrices need not be
principal submatrices of one another.
\end{proof}
Proposition~\ref{long1049:res:finite-pencil} uses only the eight
unshifted certificates $D_{k,0}$, not an all-rank coefficient measure.
The shifted certificates $D_{k,1}$ are used below for a truncated
Stieltjes moment representation and the first fifteen continued-fraction
coefficients, not for this application of Sylvester's criterion.
An infinite-support coefficient measure would extend the pencil argument
to all ranks. At the certified ranks, the largest root is a nondecreasing
lower bound for $F(p)$; convergence to $F(p)$, simple roots, strict
interlacing and denominator control are not established. At $p=1$ the
congruence degenerates and $F$ diverges, so the proposition excludes that
endpoint.

The pencil is different from the Jacobi matrix of the coefficient
moments. This is already visible at rank two. The defining formulas give
\[
 \beta_0=0,\qquad \beta_1=p^3(p-1)^2(p+2)>0\quad(p>1).
\]
Consequently $\det B_2=-\beta_1^2<0$. Since $A_2>0$, the product of the
two real pencil roots is $\det B_2/\det A_2<0$: one root is negative
and the other positive. In contrast, the finite Stieltjes construction
below has only positive nodes. Positivity of the coefficient Hankel
matrix therefore does not identify these two spectral constructions.

\paragraph{Formal continued fractions and finite quadrature.}
Write $D_n=D_{n,0}$, $E_n=D_{n,1}$, with $D_0=E_0=1$.
All these determinant polynomials are nonzero by the all-rank degree
calculation above. The classical Stieltjes continued fraction is thus
defined over the rational-function field $\Q(p)$, with the convention
\[
 \sum_{m\ge0}s_m(p)z^m
 =\cfrac{1}{1-\cfrac{\lambda_1z}{1-\cfrac{\lambda_2z}{1-\cdots}}},
\]
\[
 \lambda_{2n-1}=\frac{E_nD_{n-1}}{D_nE_{n-1}},\quad
 \lambda_{2n}=\frac{D_{n+1}E_{n-1}}{E_nD_n}\quad(n\ge1).
\]
Here $s_0=1$, and the identity is in $\Q(p)[[z]]$. The classical
moment correspondence is recalled in \cite[pp.~1--2]{sw2024}.
Specialising a ratio at a fixed real $p$ requires its displayed
denominators to be nonzero. The sixteen certificates ensure this and
$\lambda_1,\ldots,\lambda_{15}>0$ for every $p\ge1$; they do not ensure
it at all indices. Seeking positive formulas for all these ratios is
therefore a specific version of the coefficient-moment question.
The generating series has radius of convergence zero for every
$p\ge1$, so the displayed identity must be read formally. At $p=1$
this follows from $s_m(1)=(m!)^3$. For $p>1$, the $k=0$ term of the
positive expansion for $R_m$ and the inequality
$[m]_p!\ge p^{m(m-1)/2}$ give
\[
 s_m(p)\ge p^{2m^2-m}.
\]
Thus $s_m(p)^{1/m}\to\infty$. This concerns the moment power series;
it does not decide convergence of the continued fraction at an
individual nonzero value of $z$.

For each fixed $p\ge1$, the two positive $8\times8$ moment matrices give
a positive measure with eight atoms representing $s_0,\ldots,s_{15}$. This can be proved
without assuming an infinite representing measure. Set
\[
 H_0=(s_{i+j})_{i,j<8},\qquad H_1=(s_{i+j+1})_{i,j<8},\qquad
 T=H_0^{-1}H_1,
\]
and give $\R^8$ the inner product
$\langle u,v\rangle=u^{\mathsf T}H_0v$.
The operator $T$ is self-adjoint and positive definite because
$H_0T=H_1$ is symmetric and positive definite. If $e_0,\ldots,e_7$
are the coordinate vectors, the Hankel identities give
$Te_j=e_{j+1}$ for $0\le j<7$. Hence $e_0$ is cyclic: its first
eight iterates form a basis. The spectral theorem now gives eight
distinct positive eigenvalues and a positive weight at each, namely
the squared norm of the corresponding projection of $e_0$.
Cyclicity ensures that no projection vanishes and that no eigenspace
has dimension greater than one.

Let $\nu$ be the resulting atomic measure. For $0\le m\le14$, choose
$i,j\le7$ with $i+j=m$. Self-adjointness gives
\[
 \int x^m\,d\nu(x)=\langle T^m e_0,e_0\rangle
 =\langle e_i,e_j\rangle=s_m.
\]
For the last moment,
\[
 \int x^{15}\,d\nu(x)=\langle T^7e_0,T T^7e_0\rangle
 =e_7^{\mathsf T}H_1e_7=s_{15}.
\]
The weights sum to $s_0=1$. Hence the measure agrees with the moment
functional on every polynomial of degree at most $15$; this is the
degree range certified for the eight-node quadrature obtained here.
In an orthonormal polynomial basis it is the classical Jacobi construction;
Golub and Welsch describe the computation of its nodes and weights
\cite[Sec.~2, (2.6); Sec.~4]{golubwelsch1969}.
Nothing here identifies the subsequent moments of $\nu$ with
$s_{16},s_{17},\ldots$. This finite construction must also be distinguished
from the coefficient pencil $(A_N,B_N)$.
For rational modifications of an already known moment functional, see
Krattenthaler \cite[Thm.~1; Sec.~6]{krattenthaler2021}; such a modification
producing the entire sequence $s_m$ has not been identified.

\paragraph{Testing for additional cyclotomic factors.}
Independently of the moment question, let $v_{\Phi_d}$ denote the
multiplicity of the cyclotomic factor $\Phi_d$ in a nonzero polynomial
in $\Q[p]$, and set $v_{\Phi_d}(0)=\infty$.
For a $2\times2$ matrix with entry valuations $0,2,2,5$, the two
determinant products have valuations $5$ and $4$. The lower one is
unique, so the determinant has valuation $4$. A tie is different:
all four entries of
$\left(\begin{smallmatrix}1&1\\1&1+\Phi_d\end{smallmatrix}\right)$
have valuation zero, but its determinant is $\Phi_d$.
At general rank, we must find the least total valuation and then test
whether the terms attaining it cancel. For $N,d\ge1$, define
\[
 \begin{aligned}
 \mathcal H_N(Y;p)&=\det(\alpha_{i+j}Y-\beta_{i+j})_{i,j<N},\\
 t_{m,d}&=\min(v_{\Phi_d}(\alpha_m),v_{\Phi_d}(\beta_m)),\\
 e_{N,d}&=\min_{\sigma\in S_N}\sum_{i=0}^{N-1}t_{i+\sigma(i),d}.
 \end{aligned}
\]
The last expression is a minimum-cost assignment: row $i$ is matched
to column $\sigma(i)$ at cost $t_{i+\sigma(i),d}$.
The valuation of a polynomial in $Y$ means the minimum valuation of
its coefficients in $\Q[p]$. Each determinant term is
divisible by $\Phi_d^{e_{N,d}}$. Set
\[
 \Gamma_{N,d}(Y)=
 \sum_{\substack{\sigma\in S_N\\\sum_i t_{i+\sigma(i),d}=e_{N,d}}}
 \operatorname{sgn}(\sigma)
 \prod_i\overline{\Phi_d^{-t_{i+\sigma(i),d}}
        (\alpha_{i+\sigma(i)}Y-\beta_{i+\sigma(i)})},
\]
where the bar is reduction in $(\Q[p]/(\Phi_d))[Y]$.
Then $v_{\Phi_d}\mathcal H_N=e_{N,d}$ exactly when
$\Gamma_{N,d}\ne0$. This isolates the cancellation that a genuine
additional factor would require.
This test also has a determinant form. The integer costs are finite,
since every $\alpha_m$ is nonzero. Minimum-cost assignment duality gives
integer row and column potentials with
\[
 u_i+v_j\le t_{i+j,d},\qquad
 \sum_i u_i+\sum_jv_j=e_{N,d}.
\]
Here is a direct construction, so their existence is not an additional
hypothesis. Choose a minimum-cost permutation $\sigma$. On the column
indices put a directed edge from $\sigma(i)$ to $j$ of length
$t_{i+j,d}-t_{i+\sigma(i),d}$. A negative directed cycle would improve
$\sigma$ by reassigning the corresponding rows along that cycle;
there is therefore no such cycle. Add a new vertex with an edge of
length zero to every column. The shortest-path distances $v_j$ are
finite integers and satisfy
\[
 v_j-v_{\sigma(i)}\le t_{i+j,d}-t_{i+\sigma(i),d}.
\]
Putting $u_i=t_{i+\sigma(i),d}-v_{\sigma(i)}$ gives the inequalities
above, with equality on the chosen assignment and hence equality of
totals.

Divide row $i$ by $\Phi_d^{u_i}$ and column $j$ by $\Phi_d^{v_j}$
over $\Q(p)$. The inequalities ensure that every resulting entry is
in $\Q[p,Y]$, even if some potentials are negative. Reduce these
entries modulo $\Phi_d$. Only those with
$u_i+v_j=t_{i+j,d}$ survive. The determinant of the reduced matrix is
$\Gamma_{N,d}$, because every non-minimal permutation vanishes on
reduction. Thus the valuation exceeds $e_{N,d}$ exactly when the
reduced matrix is singular over $(\Q[p]/(\Phi_d))(Y)$.
A kernel relation valid at just one value of $Y$ is insufficient:
the determinant must vanish as a polynomial in $Y$. Conversely, one
nonzero value $\Gamma_{N,d}(Y_0)$ proves equality with the assignment
bound. Neither assignment duality nor the existence of the potentials
forces the reduced determinant to vanish.


The supplied exact calculation gives the following monic contents.
Here $\operatorname{cont}_Y$ means the monic gcd in $\Q[p]$ of the
coefficients in $Y$:
\[
\begin{array}{c|l}
N&\operatorname{cont}_Y\mathcal H_N\\\hline
1&p\\
2&p^5(p-1)^4\\
3&p^{14}(p-1)^{15}(p+1)^4\\
4&p^{30}(p-1)^{32}(p+1)^8(p^2+p+1)^4\\
5&p^{55}(p-1)^{55}(p+1)^{19}(p^2+1)^4(p^2+p+1)^8.
\end{array}
\]
These factorisations attain the assignment bound for every $d$ at
these five ranks. The residue certificates verify equality for
$d\le8$; the displayed complete factorisations contain no other
cyclotomic factors, so the nonnegative assignment bound is zero at
all remaining $d$. The power of $p$ is not cyclotomic content.

For the table, the script computes the full polynomials
$\mathcal H_N(Y_0;p)$ at $Y_0=0,\ldots,N$ by fraction-free
elimination, checking every polynomial division. Their monic gcd in
$\Q[p]$ is $\operatorname{cont}_Y\mathcal H_N$: evaluation gives one
divisibility direction, and interpolation of the degree-at-most-$N$
polynomial in $Y$ gives the other. The calculation includes every
factor, not just a prescribed list of cyclotomic candidates.

The same script produces a nonzero residue witness for every pair
\[
 1\le N\le8,\qquad 1\le d\le\max(8,2N-2).
\]
There are $76$ such pairs. For each, the certificate gives an optimal
assignment, integer dual potentials, an integer $Y_0\in\{0,\ldots,N\}$
and the nonzero residue $\Gamma_{N,d}(Y_0)$. The signed
minimum-cost subset recurrence and the determinant of the reduced
matrix give the same residue. A separate check computes the full
polynomial $\mathcal H_N(Y_0;p)$ and verifies directly that division
by $\Phi_d^{e_{N,d}}$ leaves exactly this nonzero residue.
At $N=2,d=1$, for example, the entry valuations are
$\left(\begin{smallmatrix}0&2\\2&5\end{smallmatrix}\right)$,
$e_{2,1}=4$, and $\Gamma_{2,1}(0)=-9$.

The source coefficients, determinant-value polynomials and all
witnesses are in \path{computations/cyclotomic_content.json}.
The scripts \path{computations/check_cyclotomic_content.py} and
\path{computations/verify_cyclotomic_content.py} reproduce the
calculation and its direct determinant checks. At ranks six to eight
these certificates cover only the stated cyclotomic-index window;
complete coefficient contents are computed only through rank five.

These computations distinguish systematic common factors from
cancellation beyond entry valuations, a distinction also important in
fraction-free matrix decompositions \cite[Sec.~2; Thm.~6]{middeke2021}.
In particular the gcd argument at the points $0,\ldots,N$ is over
$\Q[p]$: the rational Vandermonde inverse is not generally integral, so
it must not be used to assert the corresponding coefficient gcd in $\Z[p]$.


Proposition~4 of Krattenthaler--Rochev--V\"a\"an\"anen--Zudilin
\cite[pp.~14--15; (4.10), p.~17]{krvz2009} obtains cyclotomic factors
through root-of-unity annihilation for a different first-order tail
recurrence. No corresponding recurrence has been established for this
moving-diagonal pencil. The finite witnesses prove equality with the assignment bound at
these $76$ pairs. They do not exclude excess at larger ranks or, at
ranks six to eight, cyclotomic indices outside the stated window. Whether minimum-valuation residues ever cancel
systematically, and whether $s_m$ is a Stieltjes sequence, remain separate
questions. The remainder measure answers neither.

\section{Rescaling integer rows}\label{long1049:sec:primitive}

Fix a rational base $a/b>1$. An irrationality argument constructs
integer-coefficient forms in $1$ and $F(a/b)$ that are nonzero and
tend to zero. Bounds on coefficient height help construct those forms or bound an
irrationality exponent, but height times error need not tend to zero for
the one-form irrationality criterion.
An integer row $(U,V)$ is primitive when $\gcd(|U|,|V|)=1$.
Dividing a nonzero row by this gcd divides its remainder by the same
integer. For an integer coefficient pair \((U,V)\) and a real target
\(S\), put
\[
 L_S(U,V)=US-V,
 \qquad
 \Delta\bigl((U_n,V_n),(U_m,V_m)\bigr)=U_nV_m-U_mV_n.
\]
We call \(L_S(U,V)\) the \emph{error} of the row at \(S\); a good approximation
is one that makes it small.  The second expression is their $2\times2$ determinant. It eliminates
\(S\) exactly:
\[
\Delta=U_mL_S(U_n,V_n)-U_nL_S(U_m,V_m).
\]

If this integer determinant does not vanish, then
\[
 1\le|\Delta|
  \le|U_m|\,\bigl|L_S(U_n,V_n)\bigr|+|U_n|\,\bigl|L_S(U_m,V_m)\bigr|,
\]
and the two errors cannot both be smaller than
\(1/(|U_n|+|U_m|)\).  The next theorem compares the divisor introduced by rescaling with
the resulting change in the determinant's absolute value.

\begin{theorem}[rescaling rows and their determinant]\label{long1049:res:content}
Let \(S\) be real, let \((U_n,V_n)\) and \((U_m,V_m)\) be pairs of integers, and
let \(c_n,c_m\) be integers.  Then
\[
 L_S(c_nU_n,c_nV_n)=c_nL_S(U_n,V_n),
\]
\[
 \Delta\bigl(c_n(U_n,V_n),c_m(U_m,V_m)\bigr)
 =c_nc_m\Delta\bigl((U_n,V_n),(U_m,V_m)\bigr),
\]
and consequently
\[
 \left|\Delta\bigl(c_n(U_n,V_n),c_m(U_m,V_m)\bigr)\right|
 =|c_n|\,|c_m|\,
   \left|\Delta\bigl((U_n,V_n),(U_m,V_m)\bigr)\right|.
\]
In particular \(c_nc_m\) divides the scaled determinant. Multiplication
by these scalars therefore introduces a divisor whose absolute value is
exactly the factor multiplying the determinant's absolute value.
\end{theorem}
\leannote{Lean: \lproof{ErdosProblems/Erdos1049/PaperFiniteAssembliesR7.lean}{59}{integer_scalar_content}.}

\begin{proof}
Expand the linear form and use the bilinearity of the determinant.
The original integer determinant is the quotient witnessing divisibility;
it need not be primitive.
\end{proof}

\begin{example}\label{long1049:ex:content}
Take \((U_n,V_n)=(1,2)\) and \((U_m,V_m)=(3,5)\), so that
\(\Delta=1\cdot5-3\cdot2=-1\).  Multiplying the first row by \(c_n=6\) and the
second by \(c_m=10\) gives the rows \((6,12)\) and \((30,50)\), whose
determinant is \(6\cdot50-30\cdot12=-60\).  That determinant is now divisible
by \(60\), which looks like a local gain of \(60\); and its absolute value has
risen from \(1\) to \(60\), which is a cost of exactly the same size.
\end{example}

\noindent Lean checks the error identity in
\lword{Erdos1049/RationalPadeArithmetic.lean}{84}{rationalPadeError_mul}{error
scaling}, the determinant identity in
\lword{Erdos1049/RationalPadeArithmetic.lean}{94}{rationalPadeExteriorDet_mul_contents}{content
factorisation}, the exact absolute-height identity in
\lword{Erdos1049/RationalPadeArithmetic.lean}{104}{natAbs_rationalPadeExteriorDet_mul_contents}{absolute
determinant scaling}, and the divisor statement in
\lword{Erdos1049/RationalPadeArithmetic.lean}{116}{contentProduct_dvd_rationalPadeExteriorDet}{content-product
divisibility}.  The elimination identity is the checked
\lword{Erdos1049/RationalPadeArithmetic.lean}{124}{rationalPadeExteriorDet_cast_eq}{exterior
determinant identity}.

The identities allow zero scalars, but cancelling the scalar factors
requires them to be nonzero. Multiplication alone gives no improvement
in the comparison between a divisor and the determinant's absolute
value. The theorem neither asserts that the determinant is nonzero nor
estimates the original approximation error. Cancelling polynomial
factors before evaluation, finding factors common to several minors,
and adding different rows remain separate operations.

\section{Congruences after evaluation at \texorpdfstring{$3/2$}{3/2}}\label{long1049:sec:endpoints}

The polynomials in Zudilin's construction give integer rows after
evaluation and denominator clearing~\cite[Secs.~3--5]{zudilin2004}.
The next results apply to arbitrary polynomials in $\Z[X]$, not just
those polynomials. We first compute the evaluated integers modulo $2$
and $3$, then count their possible residues modulo powers of these
primes. The tools are elementary congruences, the pigeonhole principle
and B\'ezout's identity. The final scalar inequality,
Theorem~\ref{long1049:res:scalar}, concerns the parameters of Zudilin's
construction rather than its coefficient polynomials.

Substituting \(X=3/2\) into an integer polynomial produces a rational
number whose denominator is cleared by a sufficiently large power of
\(2\). To use the same operation when adding polynomials, first fix a
truncation index \(W\ge0\). For \(P(X)=\sum_i p_iX^i\in\Z[X]\), put
\[
 H_W(P)=\sum_{i=0}^{W}p_i\,3^i2^{W-i}.
\]
This is the
\lword{Erdos1049/ZudilinConeArithmetic.lean}{98}{homEvalThreeTwo}{denominator-cleared evaluation at $3/2$}.  Each term \(p_iX^i\) with \(i\le W\) contributes
\(p_i3^i2^{W-i}\). Keep the index \(W\) fixed when polynomials are
added. When \(W\ge\deg P\), this is exactly the cleared numerator, since
\[
 \sum_{i=0}^{W}p_i\,3^i2^{W-i}=2^{W}\sum_{i=0}^{W}p_i\left(\tfrac32\right)^{i}
 =2^{W}P\!\left(\tfrac32\right).
\]

Here a unit coefficient in $\Z$ means $1$ or $-1$. Thus a monic
polynomial of degree $W$ with constant term $\pm1$ satisfies both unit
conditions below. Those conditions are sufficient, not necessary:
the congruences themselves only require that the relevant coefficient
not be divisible by the relevant prime.

\begin{theorem}[endpoint residues]\label{long1049:res:endpoints}
Let \(P=\sum_ip_iX^i\in\Z[X]\) and let \(W\ge0\).  Then
\[
 H_W(P)\equiv p_0\,2^W\pmod 3,\qquad
 H_W(P)\equiv p_W\,3^W\pmod 2.
\]
Consequently a unit constant coefficient prevents divisibility by \(3\), and
a unit coefficient at index \(W\) prevents divisibility by \(2\).
\end{theorem}
\leannote{Lean: \lproof{ErdosProblems/Erdos1049/PaperFiniteAssembliesR7.lean}{74}{endpoint_residues}.}

\begin{proof}
Modulo \(3\), every summand with \(i>0\) vanishes; modulo \(2\), every summand
with \(i<W\) vanishes.  The remaining powers are units in the corresponding
residue fields.
\end{proof}

Since \(2^{W}\) is invertible modulo \(3\) and \(3^{W}\) is invertible modulo
\(2\), the two congruences say more than the stated consequence: divisibility
of \(H_W(P)\) by \(3\) is decided by \(p_0\) alone, and divisibility by \(2\) by
\(p_W\) alone.  The rest of the coefficient vector is invisible to both primes.
The coefficient \(p_W\) is the top coefficient of \(P\) exactly when
\(\deg P=W\), and is zero when \(\deg P<W\). The identity
\(H_W(P)=2^{W}P(3/2)\) is guaranteed when \(\deg P\le W\).
Without a degree bound, increasing the truncation index gives
\[
 H_{W+1}(P)=2H_W(P)+p_{W+1}3^{W+1},
\]
so the truncated sum doubles precisely when the extra coefficient is zero.

\begin{example}\label{long1049:ex:endpoints}
Take \(W=2\).  The three polynomials below differ only at an endpoint.
\begin{center}
\begin{tabular}{@{}llcc@{}}
\toprule
\papertablehead
\(P\) & \(H_2(P)\) & \(3\mid H_2(P)\) & \(2\mid H_2(P)\)\\
\midrule
\(X^2+1\)  & \(1\cdot4+0\cdot6+1\cdot9=13\) & no  & no \\
\(X^2+3\)  & \(3\cdot4+0\cdot6+1\cdot9=21\) & yes & no \\
\(2X^2+1\) & \(1\cdot4+0\cdot6+2\cdot9=22\) & no  & yes \\
\bottomrule
\end{tabular}
\end{center}
The first has both endpoints equal to \(1\) and its evaluation, \(13\), is
divisible by neither prime; as a check, \(2^{2}\bigl((3/2)^2+1\bigr)=13\).  The
second and third show that each hypothesis is used: spoiling the constant
endpoint admits \(3\), and spoiling the top endpoint admits \(2\).
\end{example}

\noindent Formal proofs cover the congruences
\lword{Erdos1049/ZudilinConeArithmetic.lean}{203}{homEvalThreeTwo_mod_three}{modulo $3$} and
\lword{Erdos1049/ZudilinConeArithmetic.lean}{320}{homEvalThreeTwo_mod_two}{modulo $2$}, together with the consequences for a
\lword{Erdos1049/ZudilinConeArithmetic.lean}{338}{three_not_dvd_homEvalThreeTwo_of_const_unit}{unit constant coefficient} and a
\lword{Erdos1049/ZudilinConeArithmetic.lean}{356}{two_not_dvd_homEvalThreeTwo_of_top_unit}{unit coefficient of $X^W$}.

In the following proposition, a unit top endpoint means that the
coefficient at index $W$ is $1$ or $-1$, with no degree bound imposed.
The two conditions apply to different entries of the polynomial pair.

\begin{proposition}[common divisor]\label{long1049:res:commonmult}
Let \(U,V\in\Z[X]\) and let \(W\ge0\).  If \(U\) has unit top endpoint, \(V\)
has unit constant endpoint, and an integer
\(c\) divides both \(H_W(U)\) and \(H_W(V)\), then
\[
 2\nmid c\qquad\text{and}\qquad 3\nmid c.
\]
\end{proposition}
\leannote{Lean: \lproof{ErdosProblems/Erdos1049/ZudilinConeArithmetic.lean}{398}{commonMultiplier_not_two_not_three_of_endpoint_units}.}
\begin{proof}
If $2\mid c$, then $2\mid H_W(U)$, contrary to the unit coefficient
at index $W$ and Theorem~\ref{long1049:res:endpoints}. Similarly,
$3\mid c$ would imply $3\mid H_W(V)$, contrary to the unit constant
coefficient.
\end{proof}

\noindent This is the checked
\lword{Erdos1049/ZudilinConeArithmetic.lean}{398}{commonMultiplier_not_two_not_three_of_endpoint_units}{common-divisor
exclusion}.  The proposition does not say that the two evaluations are coprime;
it says only that every common divisor is coprime to $6$. It gives
no restriction on the other prime factors of that divisor.

\begin{example}\label{long1049:ex:commonmult}
Take \(W=2\), \(U=X^2+3\) and \(V=5X^2+1\).  The top endpoint of \(U\) and the
constant endpoint of \(V\) are both \(1\), and
\[
 H_2(U)=3\cdot4+1\cdot9=21,\qquad H_2(V)=1\cdot4+5\cdot9=49 .
\]
Here $\gcd(21,49)=7$. The endpoint assumptions therefore allow a
nontrivial common divisor, but that divisor is coprime to $6$.
\end{example}

\begin{corollary}[limits of rescaling and common-divisor cancellation at $3/2$]\label{long1049:res:nomult}
Under the endpoint hypotheses of Proposition~\ref{long1049:res:commonmult},
every common divisor of the unscaled evaluations $H_W(U)$ and $H_W(V)$
is coprime to $6$. Multiplying two integer rows by nonzero integers
$c_n,c_m$ multiplies their determinant by $c_nc_m$ and its absolute
value by $|c_nc_m|$. Cancelling this introduced scalar factor therefore
leaves the original comparison between divisor and determinant size
unchanged.
\end{corollary}
\leannote{Lean: \lproof{ErdosProblems/Erdos1049/PaperFiniteAssembliesR7.lean}{220}{endpoint_scalar_content_exclusion}.}

\begin{proof}
The common-divisor assertion is Proposition~\ref{long1049:res:commonmult}.
By Theorem~\ref{long1049:res:content}, rescaling the two rows by
$c_n,c_m$ multiplies their errors by $c_n,c_m$, respectively, and
their $2\times2$ determinant by $c_nc_m$. For nonzero scalars, the
introduced divisor has the same absolute value as the introduced
factor in that determinant. Dividing out that factor restores the
original determinant, including its original divisor-to-size comparison.
\end{proof}

One further consequence of Theorem~\ref{long1049:res:endpoints} is worth stating,
because it bears on the most natural way one might hope to import an existing
denominator reduction.  Write $\Phi_m$ for the $m$th cyclotomic polynomial and,
for coprime $a>b\ge1$, put $\Phi_m(a,b)=b^{\varphi(m)}\Phi_m(a/b)$ for its
homogenisation, with exponent $\varphi(m)=\deg\Phi_m$.

\begin{proposition}[coprimality of homogenised cyclotomic values]
\label{long1049:res:cyclounit}
Let $a>b\ge1$ with $\gcd(a,b)=1$ and let $m\ge1$.  Then
$\gcd(\Phi_m(a,b),ab)=1$.  In particular $\gcd(\Phi_m(3,2),6)=1$ for every $m$.
\end{proposition}
\leannote{Lean: \lproof{ErdosProblems/Erdos1049/ZudilinConeArithmetic.lean}{302}{cyclotomicHomEval_isCoprime_mul}.}

\begin{proof}
The polynomial $\Phi_m$ is monic and $\Phi_m(0)=\pm1$. The endpoint
argument in Theorem~\ref{long1049:res:endpoints} also works at $(a,b)$:
modulo a prime dividing $b$, only $a^{\varphi(m)}$ survives, and
modulo a prime dividing $a$, only $\Phi_m(0)b^{\varphi(m)}$ survives.
Coprimality of $a$ and $b$ makes each surviving term nonzero.
\end{proof}

The kernel-checked declaration
\lword{Erdos1049/ZudilinConeArithmetic.lean}{302}{cyclotomicHomEval_isCoprime_mul}{coprimality of homogeneous cyclotomic values}
proves Proposition~\ref{long1049:res:cyclounit} in the same homogeneous-evaluation
representation, under the displayed coprimality assumptions. The
later analytic deductions require their own proofs.

Rhin and Viola's factorial-coset quotients~\cite{rhinviola1996}
reduce denominators of rational forms in $\zeta(2)$. Zudilin's
$q$-analogue has an order-twelve symmetry group for the normalised
series; its denominator cancellation uses the order-six subgroup
preserving the required sign condition
\cite[Sec.~4, pp.~159--161]{zudilin2004}.
Proposition~\ref{long1049:res:cyclounit} applies to these cyclotomic
factors: each homogenised value at $(3,2)$, and hence any product of
them, is divisible by neither $2$ nor $3$. It makes no such assertion
about an arbitrary integer or factorial factor, which can contain
both primes. Cancelling a common cyclotomic factor can still reduce
the real absolute values. That reduction is distinct from producing
the powers of $2$ or $3$ sought in Section~\ref{long1049:sec:open}.

\noindent The rescaling and common-divisor statements used in
Corollary~\ref{long1049:res:nomult} have the Lean proofs cited above;
their combination is an ordinary deduction, not a separately formalised
result. Scalar multiplication can introduce powers of $2$ and $3$,
but does not improve this comparison. Under the endpoint hypotheses,
the unscaled pair has no common factor at either prime. Neither
observation shows that a gain at these primes is necessary for a proof
by linear forms at $3/2$.

The operation studied next is different: take an integer combination
of several rows and ask that the combination be divisible where the
individual rows are not. The proposed divisor must come from
cancellation between rows, not from multiplying a row by that divisor.
The endpoint congruences are the first case of a divisibility condition that can be imposed
to any depth, and it is that condition, read additively, which is counted
below.

We first raise the two congruences to prime powers.  Fix depths \(R,S\ge0\).
For $P\in\Z[X]$, let $J_{3,R}(P)$ and $J_{2,S}(P)$ be the residues of
$H_W(P)$ modulo $3^R$ and $2^S$, respectively.
Theorem~\ref{long1049:res:endpoints} computes them when $R=S=1$.  Their
vanishing is exactly the requested divisibility:
\[
 J_{3,R}(P)=0\iff 3^R\mid H_W(P),\qquad
 J_{2,S}(P)=0\iff 2^S\mid H_W(P).
\]
These are the checked
\lword{Erdos1049/ZudilinConeArithmetic.lean}{191}{bottomJet3_eq_zero_iff_dvd}{criterion for divisibility by $3^R$} and
\lword{Erdos1049/ZudilinConeArithmetic.lean}{197}{topJet2_eq_zero_iff_dvd}{criterion for divisibility by $2^S$}.

The \emph{vector of four residues} of a coefficient pair
\((U,V)\) is then the quadruple
\[
 \bigl(J_{3,R}(U),J_{3,R}(V),J_{2,S}(U),J_{2,S}(V)\bigr)
 \in(\Z/3^R\Z)^2\times(\Z/2^S\Z)^2 ,
\]
two residues for each of the two primes, one from each entry of the pair.  By
the displayed criteria it vanishes exactly when \(3^R\) divides both specialised
entries and \(2^S\) divides both.

Instead of requiring each input row to have a common divisor, we seek
a small integer combination whose two entries are both divisible by
$3^R2^S$. Since $H_W$ is additive, the four residues of the combination
are the corresponding sums of the input residues. This allows a
pigeonhole argument on subset sums.

\begin{theorem}[equal residues for two subset sums]\label{long1049:res:jetkernel}
Fix a truncation index \(W\) and depths \(R,S\), and let \((U_j,V_j)_{j<M}\) be any \(M\)
pairs of integral polynomials. Represent each subset of
\(\{0,\dots,M-1\}\) by its indicator vector in \(\{0,1\}^M\).
If the \(2^M\) subsets outnumber the possible residue vectors in
\[
 (\Z/3^R\Z)^2\times(\Z/2^S\Z)^2,
\]
then two distinct subsets have the same residue vector.  Subtracting their
indicator vectors gives a nonzero coefficient vector in
\(\{-1,0,1\}^M\) cancelling all four residues.  The target has exact cardinality
\[
 (3^R)^2(2^S)^2.
\]
In particular, if \(R>0\) and \(4R+2S\le M\), such a collision exists.
\end{theorem}
\leannote{Lean: \lproof{ErdosProblems/Erdos1049/PaperFiniteAssembliesR7.lean}{248}{fourJet_paper_statement}.}

\begin{proof}
Send each subset to the sum of the residue vectors of its members.
The pigeonhole principle gives two distinct subsets with the same
residue vector. The number of possible vectors is the product of
the four moduli.  For \(R>0\),
\[
 (3^R)^2(2^S)^2<(4^R)^2(2^S)^2=2^{4R+2S}\le2^M,
\]
which proves the stated sufficient threshold.
\end{proof}

The power bracket improves the generic coefficient $4R$ when the depth
$R$ is a positive integer multiple of $41$. All depths and row counts
in the following corollary, including $T$, are integers.

\begin{corollary}[the exact count at depth $41$]\label{long1049:res:rankfortyone}
Let $T>0$.  At modulus $3^{R}$ with $R=41T$, any family of
$M\ge130T+2S$ integral polynomial pairs has two distinct binary selectors
with the same residue vector.  For $T=1$ the coefficient $130$ is exact for
this counting argument:
\[
 2^{129+2S}<\bigl| (\Z/3^{41}\Z)^2\times
                       (\Z/2^S\Z)^2\bigr|.
\]
No exact-optimality assertion is made here for $T>1$.
\end{corollary}
\leannote{Lean: \lproof{ErdosProblems/Erdos1049/PaperFiniteAssembliesR7.lean}{149}{rank_fortyone}.}

\begin{proof}
The upper power inequality gives
\[
 (3^{41T})^2(2^S)^2
 <(2^{65})^{2T}(2^S)^2=2^{130T+2S}\le2^M,
\]
so Theorem~\ref{long1049:res:jetkernel} applies.  For $T=1$, direct integer
evaluation gives $2^{129}<3^{82}$; multiplying by $(2^S)^2$ gives the
displayed reverse count at $129+2S$.
\end{proof}

\noindent The formal statements check the
\lword{Erdos1049/AdelicHeightBridge.lean}{279}{exists_distinct_binary_selectors_same_fourJet_of_rank_41}{sufficient count $130T+2S$}
and, when $T=1$, the
\lword{Erdos1049/AdelicHeightBridge.lean}{292}{fourJet_card_gt_two_pow_of_rank_41}{failure of the count at $129+2S$}.

For all depths at once the ambient-cardinality inequality $2^{M}>3^{2R}2^{2S}$
holds exactly when
\[
 M\ \ge\ \bigl\lfloor2R\log_23+2S\bigr\rfloor+1,
\]
by the definition of the floor function; irrationality of $\log_2 3$
is not needed for this strict-inequality reformulation.
At $R=41$ this is $130+2S$, and
at $R=41\cdot31$ it is $4029+2S$, one below the uniform bound $4030+2S$ of
Corollary~\ref{long1049:res:rankfortyone}; the corollary trades exactness for a
certificate that is a single integer comparison.  Here $41$ is the exponent of $3$ in the modulus, not the number of
input rows; at that depth the counting bound is $130+2S$ rows.

Counting alone does not ensure that the two selectors produce different
analytic remainders. A bound on the number of selectors giving each
real remainder is one way to obtain that additional conclusion.

\begin{theorem}[equal residues with different values]\label{long1049:res:boundedfibre}
Let $A$ and $B$ be finite sets, let $f:A\to B$, and let $g:A\to C$ be any
map into a set $C$.  Suppose every fibre of $g$ has at most $k$ elements.  If
\[
 |B|k<|A|,
\]
then there exist distinct $x,y\in A$ such that
\[
 f(x)=f(y)\qquad\hbox{and}\qquad g(x)\ne g(y).
\]
Thus, when $f$ records the four residues and $g$ records the real
remainder, a bound on the multiplicities of equal remainders guarantees
a pair with equal residues and different remainders.
\end{theorem}
\leannote{Lean: \lproof{ErdosProblems/Erdos1049/AdelicHeightBridge.lean}{1832}{exists_ne_map_eq_map_ne_of_card_mul_lt}.}

\begin{proof}
If every pair in a common $f$-fibre also had the same $g$-value, each
$f$-fibre would lie in one $g$-fibre and hence have size at most $k$.
Summing over the at most $|B|$ fibres of $f$ would give
$|A|\le |B|k$, contrary to the hypothesis.
\end{proof}

\noindent The \lword{Erdos1049/AdelicHeightBridge.lean}{311}{exists_ne_map_eq_map_ne_of_card_mul_lt}{finite-set counting statement} is formalised.

If $g$ is injective, the hypothesis holds with $k=1$. It need not hold
with a useful small $k$ for subset sums: repeated input rows produce
many equal sums, even when every row is primitive. The theorem assumes
a bound on every fibre of $g$, not merely on the fibres of $(f,g)$.
It gives a nonzero difference but no upper bound on its real size.
The short paper states a different, quantitative version using both
residues and intervals of real values.


The next improvement needs a much stronger hypothesis than
primitivity: adjacent determinants must vanish modulo the chosen
modulus. Unit multiples of one unimodular row satisfy it. Arbitrary
primitive rows do not; for example, $(1,0)$ and $(0,1)$ have determinant
$1$. Under this hypothesis the possible sums lie on a single line,
so only one residue coordinate must be counted.

\begin{theorem}[vanishing minors and a residue count]
\label{long1049:res:plucker-collapse}
Let $R_0$ be a commutative ring and let $w_n=(A_n,B_n)\in R_0^2$.  Suppose
that each row is unimodular, meaning that $u_nA_n+v_nB_n=1$ for some
$u_n,v_n\in R_0$, and that every adjacent minor vanishes:
\[
 A_nB_{n+1}-B_nA_{n+1}=0\qquad(n\ge0).
\]
Then every pairwise minor $A_iB_j-B_iA_j$ vanishes.  In particular, take
$R_0=\Z/(2^S3^R)\Z$ with $R>0$.  If $S+2R\le k$, there are two distinct
binary selectors $s,t\in\{0,1\}^{k}$ such that
\[
 \sum_{i<k}s_iw_i=\sum_{i<k}t_iw_i.
\]
Thus $S+2R$ rows suffice. The ambient two-coordinate argument gives
the sufficient bound $2S+4R$.
\end{theorem}
\leannote{Lean: \lproof{ErdosProblems/Erdos1049/PaperFiniteAssembliesR7.lean}{269}{plucker_paper_statement}.}

\begin{proof}
If $(a,b)$ is unimodular, say $ua+vb=1$, and $ay-bx=0$, then
\[
 (x,y)=(ux+vy)(a,b).
\]
Indeed, the first coordinate follows by replacing $ay$ with $bx$, and the
second by the reverse substitution.  Apply this identity to consecutive rows:
each next row is a scalar multiple of the current one.  Induction places the
entire tail on the line through $w_0$ and proves the pairwise-minor assertion.
Right multiplication by
\[
 \begin{pmatrix}u_0&-B_0\\v_0&A_0\end{pmatrix}
\]
has determinant $u_0A_0+v_0B_0=1$ and sends $w_0$ to $(1,0)$.
All transformed selector sums therefore have second coordinate zero
and occupy at most $2^S3^R$ values.  Finally
\[
 2^S3^R<2^S4^R=2^{S+2R}\le2^k,
\]
and pigeonhole gives the two selectors.
\end{proof}

\noindent No particular coordinate needs to be invertible: modulo six,
$(2,3)$ is unimodular because $-2+3=1$, although neither entry is a unit.
Some nondegeneracy is essential. The rows $(1,0),(0,0),(0,1)$ have
zero adjacent minors but outer minor one.
For the counting conclusion, vanishing is required only in the quotient
ring. The integer rows $(1,0),(1,6)$ used in Section~\ref{long1049:sec:open}
are independent over $\Q$, with determinant $6$, but coincide modulo $6$.
Thus dependence modulo the modulus neither follows from primitivity nor
implies dependence of the integer rows.

\noindent Lean checks the unit-coordinate form inside the supported root: the
\href{https://github.com/wcook04/plectis-erdos/blob/7148c5ae4a4a2f3802cccdeed62c394f64593340/lean/ErdosProblems/Erdos1049/BezoutPluckerJets.lean\#L191}{vanishing of all pairwise minors},
the resulting
\href{https://github.com/wcook04/plectis-erdos/blob/7148c5ae4a4a2f3802cccdeed62c394f64593340/lean/ErdosProblems/Erdos1049/BezoutPluckerJets.lean\#L230}{equal residues for distinct subsets},
and the
\href{https://github.com/wcook04/plectis-erdos/blob/7148c5ae4a4a2f3802cccdeed62c394f64593340/lean/ErdosProblems/Erdos1049/BezoutPluckerJets.lean\#L249}{explicit $S+2R$ threshold}.
Each of the three assumes that the second coordinate of every row is a unit.
The unimodular-row statement proved above is the stronger one.
This conditional theorem is stronger than the ambient count of possible residue vectors only
after its minor-vanishing hypothesis has been established.  No such all-tail
hypothesis is proved here for an actual $q$-Ap\'ery or Zudilin family, and the
theorem says nothing about whether the resulting selector difference has
nonzero analytic remainder.

The count in Corollary~\ref{long1049:res:rankfortyone} is sharp at
$T=1$ for comparison with the full residue space. To obtain different
remainders, Theorem~\ref{long1049:res:boundedfibre} additionally needs
a bound on how often the same real value occurs. Such a bound is not
proved here for the $q$-Ap\'ery or Zudilin remainder family.

\noindent The target count is the checked
\lword{Erdos1049/ZudilinConeArithmetic.lean}{131}{fourJetSignature_card}{cardinality of the residue space}; the abstract collision is the checked
\lword{Erdos1049/ZudilinConeArithmetic.lean}{142}{exists_distinct_binary_selectors_same_fourJet}{pigeonhole argument for the residue map}, and the linear sufficient condition is the checked
\lword{Erdos1049/ZudilinConeArithmetic.lean}{161}{exists_distinct_binary_selectors_same_fourJet_of_rank}{sufficient row count for equal residues}. The pigeonhole argument gives
divisibility, not nonvanishing.  Pigeonhole cancellation itself requires no
independence.  Additional information about the input family is needed to
ensure that the resulting nonzero selector difference has a nonzero combined
polynomial pair and analytic remainder.  None of the statements proved here
supplies such a family or proves either nonvanishing conclusion.

\begin{example}\label{long1049:ex:jetcount}
At depths \(R=S=1\) the space of four residue coordinates is
\((\Z/3\Z)^2\times(\Z/2\Z)^2\), of cardinality \(9\cdot4=36\), and the threshold
reads \(M\ge4\cdot1+2\cdot1=6\).  With six pairs there are \(2^{6}=64\) binary
selectors against \(36\) targets, so two of them collide and their difference
is a vector in \(\{-1,0,1\}^{6}\), not identically zero, killing all four residues.
\end{example}

The following elementary comparison concerns only the two scalar
exponents, not the coefficient polynomials.

\begin{theorem}[a restriction on the two scalar exponents]\label{long1049:res:scalar}
Let \(C_1>0\).  If \(C_0\le0\) or \(2C_0\le C_1\), then
\[
 C_0\log 3-C_1\log 2<0.
\]
In the positive branch $C_0>0$ and $2C_0\le C_1$, the stronger estimate is
\[
 C_0\log3-C_1\log2< -\frac{17}{41}C_0\log2.
\]
\end{theorem}
\leannote{Lean: \lproof{ErdosProblems/Erdos1049/PaperFiniteAssembliesR7.lean}{177}{scalar_margin}.}

\begin{proof}
If $C_0\le0$, then $C_0\log3-C_1\log2\le-C_1\log2<0$.
If $C_0>0$ and $2C_0\le C_1$, use $3^{41}<2^{65}$ to obtain
\[
 C_0\log3-C_1\log2
 \le C_0(\log3-2\log2)
 <-\frac{17}{41}C_0\log2<0.
\]
\end{proof}

\noindent Written multiplicatively, the conclusion is \(3^{C_0}<2^{C_1}\).
The inequality is the checked
\lword{Erdos1049/ZudilinConeArithmetic.lean}{428}{three_two_scalar_margin_neg}{three-halves
scalar margin}.  The positive-branch deficit is the checked
\lword{Erdos1049/AdelicHeightBridge.lean}{242}{three_two_scalar_margin_lt_explicit}{$17/41$ margin}.
The scope of this elementary inequality matters.  The primary Zudilin
theorem~\cite[Theorem~1, p.~154; Sec.~5, pp.~161--162]{zudilin2004} supplies an integer-base
irrationality-exponent estimate on its parameter cone, and the elementary
inequality \(\mu\ge2\) then forces \(2C_0\le C_1\) whenever \(C_0>0\); Lean
checks that implication separately.  The primary 2004 theorem is stated for an integer inverse base.
The rational specialisation proved earlier in this record has separate
proofs for the constructed forms. The scalar statement in this paragraph is
only the displayed inequality; it neither constructs a new family at
$3/2$ nor formalises a universal exclusion of linear-form methods.

\section{Failure of the stated clearing conditions at \texorpdfstring{$3/2$}{3/2}}\label{long1049:sec:corridor}

\paragraph{Erd\H{o}s's integer-base argument, in outline.}
For an integer $b\ge2$, Erd\H{o}s rewrites
$F(b)=\sum_{n\ge1}\tau(n)b^{-n}$.  A Chinese-remainder construction forces
arbitrarily long blocks in which the divisor coefficients have the powers of
$b$ needed to make the corresponding base-$b$ digits zero. Tail estimates
control carries, and the argument also establishes that the expansion
does not terminate. These are separate requirements: arbitrarily long
zero blocks without eventual termination exclude eventual periodicity,
and hence rationality~\cite[pp.~63--66]{erdos1948}.
Positivity of the original summands alone would not prove nontermination;
for example, $\sum_{n\ge1}(b-1)b^{-n}=1$.

\paragraph{The direct cut-and-clear attempt studied here.}
A more naive attempt is to suppose $F(b)=p/q$, cut at $N$, and multiply the
remaining identity by $qb^N$.  This does not by itself trap a positive integer
below $1$: the first uncleared tail contribution is
$q\tau(N+1)/b$.  The clearing conditions below describe one bounded-window version
of this attempt. They fail at $3/2$.

Write $\beta=r/s$ for the base and $c(n)$ for the coefficient of $\beta^{-n}$,
so that $c=\tau$ in the case at hand.  The term $c(n)\beta^{-n}$ is
$c(n)s^{n}/r^{n}$.  Clearing the
power of $r$ leaves the numerator factor $s^{n}$ in place.  That factor is
invisible when $s=1$ and grows geometrically when $s\ge2$.  At $3/2$ it is
$2^{n}$.


We return to the attempt to clear successive terms of the series.
The multiplier must satisfy a divisibility condition and must also be
small enough to leave a remainder less than $1$. These demands give
incompatible inequalities at $3/2$. The argument below uses only six
natural-number parameters; it is not an exclusion of other ways to
approximate $F(3/2)$.

\begin{definition}[clearing conditions]\label{long1049:def:corridor}
For natural numbers $a,b,N,K,Q,D$, the clearing conditions are
\[
 a>0,\qquad Q>0,\qquad D>0,\qquad D\le N+K,\qquad
 a^{K}\mid QD,\qquad Q\,b^{\,N+K+1}<a^{\,K+1}.
\]
\end{definition}

\noindent The reading is: $a$ and $b$ are the numerator and denominator of the
base, playing the roles of $r$ and $s$ in the partial-sum calculation
above, so that
$(a,b)=(3,2)$ is the case of interest;
$N$ is the shift, $K$ is the width of the cleared window, $Q$ is the
accumulated clearing factor, and $D$ is the final coefficient being cleared.
The bound $D\le N+K$ is the only property of the coefficient used; for the
divisor-counting coefficient it holds because $\tau(n)\le n$.  The divisibility $a^{K}\mid QD$ tests the final coordinate of the cleared
window; it does not by itself clear every earlier coordinate. For a
reduced base $a/b$, the denominator of
$Qa^N D(b/a)^{N+K}$ is cleared exactly when $a^K\mid QD$, since
$a$ and $b$ are coprime. For example, at $a=2$, $b=1$,
$N=21$, $K=3$ and $Q=1$, the final term clears because
$2^3\mid\tau(24)=8$, but the preceding scaled term is
$\tau(23)/2^2=1/2$ and is not integral.

The final inequality is a necessary smallness test, not a sufficient
bound for the entire tail. For $a>b\ge1$ and $Q>0$, positivity and
$\tau(m)\ge1$ give
\[
 Q a^N\sum_{m\ge N+K+1}\tau(m)(b/a)^m
 >\frac{Qb^{N+K+1}}{a^{K+1}}.
\]
A tail smaller than $1$ must therefore satisfy the final inequality.
The divisibility gives $a^K\le Q(N+K)$, while the smallness test requires
$Qb^{N+K+1}<a^{K+1}$. The same $Q$ must satisfy both. These
conditions are defined in Lean by the
\lword{Erdos1049/RationalBaseLambert.lean}{113}{CoordinatewiseCorridor}{six clearing conditions}.

\begin{theorem}[a necessary inequality for clearing]\label{long1049:res:corridorbound}
If $(a,b,N,K,Q,D)$ satisfies the clearing conditions, then
\[
 b^{\,N+K+1}<a\,(N+K).
\]
\end{theorem}
\leannote{Lean: \lproof{ErdosProblems/Erdos1049/RationalBaseLambert.lean}{121}{coordinatewiseCorridor_implies_pow_lt_linear}.}

\begin{proof}
Since $QD>0$ and $a^{K}\mid QD$, we have
$a^{K}\le QD$, and $D\le N+K$ gives
$a^{K}\le Q(N+K)$.  Hence
\[
 Q\,b^{\,N+K+1}<a^{\,K+1}=a^{K}\cdot a\le Q(N+K)\cdot a,
\]
and cancelling the positive factor $Q$ gives the claim.
\end{proof}

\noindent Formalised as the
\lword{Erdos1049/RationalBaseLambert.lean}{121}{coordinatewiseCorridor_implies_pow_lt_linear}{power-versus-linear
consequence}.

The necessary inequality in Theorem~\ref{long1049:res:corridorbound}
shows the difference between these clearing tests at integer and
noninteger rational bases.  At $b=1$ it reads $1<a(N+K)$, which holds for every
$a\ge2$ and every nonempty window; this necessary inequality imposes no
obstruction.  The other clearing conditions still have to be checked.
At $b\ge2$ the left side is exponential in $N+K$ and the right side is linear,
so the conditions can hold only for bounded $N+K$ at each fixed base.  At $(a,b)=(3,2)$ the
necessary inequality already fails when $N=K=1$.

\begin{example}\label{long1049:ex:corridor}
Take the smallest window with $N,K\ge1$, namely $N=K=1$, and numerator $a=3$.  At $b=1$ the tuple
$(3,1,1,1,3,1)$ satisfies the conditions: $D=1\le2$, the divisibility reads $3\mid3$, and
the last inequality reads $3\cdot1^{3}=3<3^{2}=9$. At $b=2$ no choice works.
The last inequality becomes $Q\cdot2^{3}<3^{2}$, which forces $Q=1$; the divisibility
then reads $3\mid D$, and the only candidates are $D=1$ and $D=2$, neither
divisible by $3$.  This is the proof of
Theorem~\ref{long1049:res:corridorbound} in miniature: it derives $a^{K}\le Q(N+K)$,
here $3\le2$.
\end{example}

Feasibility of these necessary conditions is not a successful tail
estimate. For $(a,b,N,K,Q,D)=(2,1,1,1,1,2)$ all six conditions hold,
and $D=\tau(N+K)$. Nevertheless,
\[
 2\sum_{m\ge3}\frac{\tau(m)}{2^m}
 >2\left(\frac2{2^3}+\frac3{2^4}+\frac2{2^5}\right)=1.
\]
The failure at $3/2$ is useful because it already occurs at the weaker,
necessary tests; their feasibility at an integer base proves no bound
for the full remainder.

\begin{proposition}\label{long1049:res:exp}
For every natural number $x\ge2$ we have $3x<2^{\,x+1}$.
\end{proposition}
\leannote{Lean: \lproof{ErdosProblems/Erdos1049/RationalBaseLambert.lean}{142}{three_mul_lt_two_pow_succ}.}

\begin{proof}
Induct on $x$, starting from $6<8$ at $x=2$.  For the step, $2^{x+1}\ge2^{2}>3$ when
$x\ge1$, so $3(x+1)=3x+3<2^{x+1}+2^{x+1}=2^{x+2}$.
\end{proof}

\noindent Formalised as the
\lword{Erdos1049/RationalBaseLambert.lean}{142}{three_mul_lt_two_pow_succ}{exponential
comparison}.

\begin{theorem}[failure of the stated clearing conditions at $3/2$]\label{long1049:res:nocorridor}
For all $N\ge1$ and $K\ge1$ and all natural $Q,D$, the tuple
$(3,2,N,K,Q,D)$ does not satisfy the clearing conditions.
\end{theorem}
\leannote{Lean: \lproof{ErdosProblems/Erdos1049/RationalBaseLambert.lean}{155}{threeHalves_no_coordinatewiseCorridor}.}

\begin{proof}
The clearing conditions would give $2^{\,N+K+1}<3(N+K)$ by
Theorem~\ref{long1049:res:corridorbound}, contradicting Proposition~\ref{long1049:res:exp}
applied to $x=N+K\ge2$.
\end{proof}

\noindent Formalised as the
\lword{Erdos1049/RationalBaseLambert.lean}{155}{threeHalves_no_coordinatewiseCorridor}{failure of the conditions at $3/2$}.

Theorem~\ref{long1049:res:nocorridor} excludes the divisibility pattern and
necessary tail test in Definition~\ref{long1049:def:corridor} at $3/2$.
It gives no conclusion about the rationality of $F(3/2)$ or about
clearing schemes with different conditions.

\section{Successive scaled remainders}\label{long1049:sec:tail}

The preceding clearing test concerns a chosen partial sum. The
recurrence below instead compares successive scaled differences from
an arbitrary rational number. It is an algebraic identity; to regard
those differences as tails one must also identify that number with
the sum of a convergent series.

Let $r,s,B,\xi$ be rationals with $r\ne0$ and let $c:\N\to\Q$ be arbitrary.
Define the prefix and the scaled remainder by
\[
 P_N=\sum_{m=0}^{N-1}\frac{c(m+1)\,s^{\,m+1}}{r^{\,m+1}},
 \qquad
 U_N=B\,r^{N}\bigl(\xi-P_N\bigr).
\tag{$\ast$}\label{long1049:eq:tailstate}
\]
Thus $P_N$ is the partial sum through index $N$, and $U_N$ is its
scaled difference from $\xi$. These are the
\lword{Erdos1049/RationalBaseLambert.lean}{168}{rationalBasePrefixQ}{rational-base partial sum} and the
\lword{Erdos1049/RationalBaseLambert.lean}{181}{rationalBaseClearedTailQ}{scaled remainder}.

\begin{theorem}[recurrence for the scaled remainder]\label{long1049:res:tailrec}
Let $r,s,B,\xi\in\Q$ with $r\ne0$, let $c:\N\to\Q$, and let $P_N$ and $U_N$ be as
in~\eqref{long1049:eq:tailstate}.  Then for every $N$,
\[
 U_{N+1}=r\,U_N-B\,c(N+1)\,s^{\,N+1}.
\]
\end{theorem}
\leannote{Lean: \lproof{ErdosProblems/Erdos1049/RationalBaseLambert.lean}{187}{rationalBaseClearedTailQ_succ}.}

\begin{proof}
Expanding
$P_{N+1}=P_N+c(N+1)s^{N+1}/r^{N+1}$ and $r^{N+1}=r^{N}\cdot r$ in the
definition of $U_{N+1}$ and clearing the denominator $r^{N+1}$, which is
nonzero, gives the identity.
\end{proof}

\noindent Formalised as the
\lword{Erdos1049/RationalBaseLambert.lean}{187}{rationalBaseClearedTailQ_succ}{recurrence for successive scaled remainders}.

For a reduced positive base $r/s>1$, the forcing term
$Bc(N+1)s^{N+1}$ contains the denominator power absent at integer
bases. Its size gives a useful bound on two consecutive remainders,
although it need not bound each remainder separately.

\begin{theorem}[the forcing term]\label{long1049:res:forcing}
Let $s,B$ be natural numbers and $c:\N\to\N$.
\begin{enumerate}
\item If $s\ge2$, $B\ge1$ and $c(N+1)\ge1$, then
$2^{\,N+1}\le B\,c(N+1)\,s^{\,N+1}$.
\item If $s=1$, then $B\,c(N+1)\,s^{\,N+1}=B\,c(N+1)$.
\end{enumerate}
\end{theorem}
\leannote{Lean: \lproof{ErdosProblems/Erdos1049/PaperFiniteAssembliesR7.lean}{186}{forcing_term}.}

\begin{proof}
For the first part,
$2^{N+1}\le s^{N+1}=1\cdot s^{N+1}\le Bc(N+1)s^{N+1}$,
using $B\,c(N+1)\ge1$.  The second part is the definition with $s=1$.
\end{proof}

\noindent Formalised as the
\lword{Erdos1049/RationalBaseLambert.lean}{204}{twoPow_le_rationalBaseForcingNat}{exponential
lower bound} and the
\lword{Erdos1049/RationalBaseLambert.lean}{218}{rationalBaseForcingNat_one}{integer-base special case}.

At $s=1$ the forcing term is $Bc(N+1)$, so it grows only as fast as the
coefficient; for the divisor function this is $O(N^{\varepsilon})$ for every
$\varepsilon>0$. This comparison does not itself construct a bounded
sequence of remainders.  At $s\ge2$ the same
term is at least $2^{N+1}$ whenever the coefficient is nonzero.

\begin{example}\label{long1049:ex:forcing}
Take $B=1$, $c=\tau$ and $N=9$, so that the coefficient is $\tau(10)=4$.  At
$s=1$ the forcing term is $4$.  At $s=2$ it is $4\cdot2^{10}=4096$, and part~(1)
of Theorem~\ref{long1049:res:forcing} already guarantees at least $2^{10}=1024$ without
knowing the coefficient at all.
\end{example}

The recurrence gives an adjacent-pair lower bound. For a natural index
$N$, with $B,s$ positive integers, $s\ge2$ and $c(N+1)\ge1$, the
triangle inequality gives
\[
 2^{N+1}\le Bc(N+1)s^{N+1}
 =|rU_N-U_{N+1}|
 \le(1+|r|)\max\{|U_N|,|U_{N+1}|\}.
\]
If these assumptions hold at every index, the full sequence is unbounded.
They do not imply $|U_N|\to\infty$ under the stated algebraic hypotheses. For example, take $r=-2$, $s=2$,
$B=1$, $c(n)=1$ and $\xi=0$. Directly from the defining partial sums,
\[
 U_{2j}=0,\qquad U_{2j+1}=-2^{2j+1}\qquad(j\ge0).
\]
The forcing term grows exponentially while every even remainder is
zero. This is an algebraic example, not a convergent Lambert series: it
shows that the adjacent-pair bound does not control every subsequence.
It concerns the displayed normalisation only and supplies no small
nonzero integer forms.

\section{The height criterion at \texorpdfstring{$7/2$}{7/2}}\label{long1049:sec:sevenhalves}

In 1994 Bundschuh and V\"a\"an\"anen proved an irrationality criterion for a
family of rational bases cut out by a height condition
\cite[Theorem~2, p.~177; hypotheses pp.~175--176]{bv1994}.  We keep their
notation: $q$ is the base, and $\alpha$
and $\lambda$ are the parameters of the criterion.  In its special case
$\alpha=-1$, the printed hypothesis is
\[
 \lambda<\left(\frac12+\frac1{\pi^2}\right)^{-1}.
\]
At $q=7/2$ the Archimedean parameter is
$\lambda=\log 7/\log(7/2)$.  The criterion therefore applies once the following
strict inequality is checked.

\begin{theorem}[the $7/2$ height condition]\label{long1049:res:sevenhalves}
\[
 \frac{\log 7}{\log(7/2)}
 <
 \left(\frac12+\frac1{\pi^2}\right)^{-1}.
\]
\end{theorem}
\leannote{Lean: \lproof{ErdosProblems/Erdos1049/RationalBaseLambert.lean}{83}{sevenHalves_archimedean_height_condition}.}
\begin{proof}
Since $2^{18}=262144<823543=7^7$, we have
$\log2/\log7<7/18$. Also $\pi>3$ gives
$1/2+1/\pi^2<11/18$. Hence
\[
 \frac{\log7}{\log(7/2)}
 =\frac1{1-\log2/\log7}<\frac{18}{11}
 <\left(\frac12+\frac1{\pi^2}\right)^{-1}.
\]
Every denominator is positive, so the reciprocal inequalities have
the displayed directions.
\end{proof}

\noindent Numerically the two sides are $1.5533\ldots$ and $1.6630\ldots$, so
the condition holds with a margin of about $0.11$.  The Lean proof factors the
estimate through the explicit
\lword{Erdos1049/RationalBaseLambert.lean}{44}{BundschuhVaananenHeightRegion}{height-region
predicate} and the
\lword{Erdos1049/RationalBaseLambert.lean}{32}{sevenHalves_power_certificate}{integer
certificate} $2^{18}<7^7$, the resulting
\lword{Erdos1049/RationalBaseLambert.lean}{48}{sevenHalves_log_ratio_lt}{logarithmic
ratio bound} $\log2/\log7<7/18$, the
\lword{Erdos1049/RationalBaseLambert.lean}{64}{one_div_pi_sq_lt_one_nine}{$\pi$-bound}
$1/\pi^2<1/9$, and the
\lword{Erdos1049/RationalBaseLambert.lean}{73}{sevenHalves_bundschuhVaananen_margin}{strict
margin}
\[
 \frac{\log2}{\log7}<\frac12-\frac1{\pi^2}.
\]
The
\lword{Erdos1049/RationalBaseLambert.lean}{83}{sevenHalves_archimedean_height_condition}{final
height inequality} then rewrites
$\log(7/2)=\log7-\log2$ and closes the displayed condition.

This formalises the complete elementary parameter check at $q=7/2$; it does
not formalise Bundschuh and V\"a\"an\"anen's analytic irrationality theorem,
whose proof occupies pp.~189--193 of the source.  The conclusion that
$F(7/2)$ is irrational is consequently cited from that theorem, not claimed as
a Lean theorem here.

The criterion does not cover $3/2$, and the results of
Sections~\ref{long1049:sec:corridor} and~\ref{long1049:sec:tail} give no evidence either way
about whether $F(3/2)$ is irrational.

\subsection{The \texorpdfstring{$81/200$}{81/200} logarithmic region}

Consider the rational-height region
\[
 \frac{\log b}{\log a}<\frac{81}{200}
 \qquad(a>b>0).
\]
We use $81/200$ as a rational lower bound on $\theta^{*}$, not as a
separately derived analytic threshold. The bracket
$81/200<\theta^{*}<1/2$ in Section~\ref{long1049:sec:regionbracket}
puts this entire smaller region inside
Theorem~\ref{long1049:res:region}. Membership of the smaller region is
an integer comparison.  The Lean module of
this library treats the displayed inequality as a definition and proves
elementary memberships and exclusions.  In particular,
\[
 \frac25
 <\frac{\log4}{\log31}
 <\frac{81}{200}
 <\frac{\log2}{\log3}.
\]
The first two comparisons come respectively from $31^2<4^5$ and
$4^{200}<31^{81}$; the last comes from $3^{81}<2^{200}$.  The lower bound is
the checked theorem
\lword{Erdos1049/ZudilinHeightRegion.lean}{65}{twoFifths_lt_thirtyoneFour_log_ratio}{two-fifths
lower bound}.  Consequently $31/4$, and every positive power $(31/4)^r$, lies
in the enlarged region
(\lword{Erdos1049/ZudilinHeightRegion.lean}{45}{thirtyoneFour_mem_zudilinHeightRegion}{$31/4$ satisfies the inequality},
\lword{Erdos1049/ZudilinHeightRegion.lean}{59}{thirtyoneFour_power_mem_zudilinHeightRegion}{power family}).
The same base lies strictly outside the earlier
Bundschuh--V\"a\"an\"anen region
(\lword{Erdos1049/ZudilinHeightRegion.lean}{91}{thirtyoneFour_outside_bundschuhVaananenHeightRegion}{$31/4$ is outside the earlier region}),
so the $81/200$ inequality defines a strict set-theoretic enlargement of the
earlier recorded logarithmic region.  Combined with the bracket
$81/200<\theta^{*}$, these memberships are exactly the finite input to
Theorem~\ref{long1049:res:31over4}, which is where the irrationality of $F(31/4)$ and of
every $F\bigl((31/4)^{r}\bigr)$ is proved.

It still does not approach $3/2$.  The exact comparison
\[
 3^{81}<2^{200}
 \quad\Longrightarrow\quad
 \frac{81}{200}<\frac{\log2}{\log3}
\]
is Lean-checked, as is the conclusion that $3/2$ belongs to neither height
region
(\lword{Erdos1049/ZudilinHeightRegion.lean}{105}{eightyOneTwoHundredths_lt_threeHalves_log_ratio}{comparison with $\log2/\log3$},
\lword{Erdos1049/ZudilinHeightRegion.lean}{122}{threeHalves_outside_zudilinHeightRegion}{$3/2$ fails the condition}).
Thus $81/200$ is the larger of the two explicitly defined cutoffs used below,
while a cutoff that includes $3/2$ must be strictly larger than
$\log2/\log3\approx0.6309$.

\section{A denominator-exponent model for Pad\'e-type
forms}\label{long1049:sec:pade}

To clear the rational coefficients in a Pad\'e approximation, a
proposed common denominator must dominate every summand's denominator.
The calculation below checks two explicit exponent expressions against
one proposed bound. No approximation family is specified for these
expressions, so the result is an algebraic comparison only. Applying
it would require a coefficient formula, an integrality proof, and a
nonzero remainder estimate for that formula.

The proposed denominator exponent is $(3n^2-n)/2$. We compare twice
each exponent, so every displayed identity is over $\Z$. The doubled
exponent is $\widetilde E_n=3n^2-n$.

\begin{proposition}[exponent model: summand bound and exact gap]\label{long1049:res:pade}
Let $\widetilde{E}_n=3n^{2}-n$ and put
\[
 \widetilde{P}(n,k)=2\bigl(k(n-k)+nk\bigr)+k(k-1),
\]
\[
 \widetilde{Q}(n,m)=2(n^{2}-n)+j^{2}+2jm+j-m^{2}+3m,
 \qquad j=n-m-1 .
\]
Then, for integers $n,k,m$:
\begin{enumerate}
\item if $0\le k\le n$, then $\widetilde{P}(n,k)\le\widetilde{E}_n$, and the
gap factors as
$\widetilde{E}_n-\widetilde{P}(n,k)=(n-k)(3n-k-1)$;
\item $\widetilde{E}_n-\widetilde{Q}(n,m)=2\bigl(n+m(m-1)\bigr)$ identically;
\item if $n\ge0$ and $m\ge1$, then $\widetilde{Q}(n,m)\le\widetilde{E}_n$.
\end{enumerate}
\end{proposition}
\leannote{Lean: \lproof{ErdosProblems/Erdos1049/PaperFiniteAssembliesR7.lean}{201}{pade_summand_bound_and_gap}.}

\noindent Part~(1) is the
\lword{Erdos1049/RationalPadeArithmetic.lean}{30}{rationalPadePSummandDenExpTwice_le}{summand
exponent bound}, part~(2) the
\lword{Erdos1049/RationalPadeArithmetic.lean}{52}{rationalPadeQMaxDenExpTwice_gap}{exact
gap identity}, and part~(3) the
\lword{Erdos1049/RationalPadeArithmetic.lean}{60}{rationalPadeQMaxDenExpTwice_le}{maximal
exponent bound}.  The gap in~(2) is an identity in $\Z[n,m]$, so~(3) follows
from $m(m-1)\ge0$ and $n\ge0$. The range $m\ge1$ is sufficient,
not necessary for this polynomial inequality: for integer $m$, one
also has $m(m-1)\ge0$ when $m\le0$. No claim about a summand of an
unspecified approximation is needed.

\begin{example}\label{long1049:ex:pade}
At $n=2$ the proposed doubled exponent is $\widetilde{E}_2=10$, and
$\widetilde{P}(2,k)$ takes the values $0,6,10$ at $k=0,1,2$.  The three gaps are
$10,4,0$, matching the factorisation $(2-k)(5-k)$ of part~(1); the summand at
$k=2$ is the one that saturates the proposed denominator.  For part~(2), at
$m=1$ we have $j=0$ and $\widetilde{Q}(2,1)=6$, with gap
$4=2\bigl(2+1\cdot0\bigr)$.
\end{example}

These are inequalities between polynomials in the exponents. They
establish that the proposed exponent $\widetilde{E}_n$ dominates the two
displayed summand exponent expressions, and nothing further.  Positivity of the remainder, its rate of
decay, and the comparison of that rate against the denominator height are the
analytic obligations, and none of them is treated here, so nothing in this
section is an irrationality measure.

\section{Complements and further questions}\label{long1049:sec:open}

The next question asks for integer combinations whose two evaluated
coefficients are divisible by large powers of $2$ and $3$, and whose
remainder is nonzero and small after division. With unrestricted input
polynomials, it is equivalent to irrationality at $3/2$, as we prove
below. The quadratic bounds in its statement do not themselves isolate
a method of proving irrationality. A question about a specified
approximation family must impose that restriction separately.

\begin{problem}[a divided linear form with small nonzero remainder]\label{long1049:prob:kernel}
Exhibit an integer constant \(C\ge1\) and, for every sufficiently large
positive integer \(n\), positive integers \(W_n,R_n,S_n,M_n\) such that
\[
 n^2\le W_n,R_n,S_n\le Cn^2,
 \qquad 4R_n+2S_n\le M_n\le Cn^2,
\]
together with polynomial pairs
\((U_{n,j},V_{n,j})\in\Z[X]^2\) for \(0\le j<M_n\), each of degree at most
the common degree bound \(W_n\), whose specialised integer rows are
primitive:
\[
 \gcd\!\bigl(H_{W_n}(U_{n,j}),H_{W_n}(V_{n,j})\bigr)=1.
\]
Find a nonzero vector
\(\lambda^{(n)}\in\{-1,0,1\}^{M_n}\) for which, on putting
\[
 U_n=\sum_{j<M_n}\lambda^{(n)}_jU_{n,j},
 \qquad V_n=\sum_{j<M_n}\lambda^{(n)}_jV_{n,j},
\]
the pair \((U_n,V_n)\) is not \((0,0)\), all four residues vanish,
\[
 J_{3,R_n}(U_n)=J_{3,R_n}(V_n)=0,
 \qquad J_{2,S_n}(U_n)=J_{2,S_n}(V_n)=0,
\]
where every residue in this display is formed using the common degree bound \(W_n\),
and the resulting divided integer linear form
\[
 A_n=\frac{H_{W_n}(U_n)}{3^{R_n}2^{S_n}},
 \qquad
 B_n=\frac{H_{W_n}(V_n)}{3^{R_n}2^{S_n}},
 \qquad
 \rho_n=A_nF(3/2)-B_n
\]
satisfies the explicit analytic condition
\[
 0<|\rho_n|<\frac1n.
\]
\end{problem}

The residue equations make \(A_n,B_n\) integers.  A solution would prove
irrationality: if \(F(3/2)=a/b\) in lowest terms, every nonzero \(\rho_n\) has
absolute value at least \(1/b\), contradicting the displayed bound for
\(n>b\).  Theorem~\ref{long1049:res:jetkernel} supplies only a nonzero signed relation
with the four residue equations once the pairs and size inequality are present; it
does not supply primitive input rows, a nonzero combined polynomial pair, or
the nonvanishing and decay of \(\rho_n\).

\paragraph{Why the choice of family matters.}
The converse uses only the pigeonhole principle and B\'ezout's identity.
Let $\xi$ be any irrational real number. Comparing the fractional parts
of $0,\xi,\ldots,(n+1)\xi$ in $n+1$ equal half-open intervals gives
integers $A\ge1$ and $B$ with
\[
 A\le n+1,\qquad 0<|A\xi-B|<\frac1{n+1}<\frac1n.
\]
Set $w=n^2$, $W_n=R_n=S_n=w$, $M_n=6w$, and $D=6^w$.
The three integer rows
\[
 (DA,1),\qquad (1,DB-1),\qquad (-1,0)
\]
are primitive and sum to $(DA,DB)$. Pad them to $M_n$ rows with
$(1,0)$, and use the selector $(1,1,1,0,\ldots,0)$.

To realise every coordinate as a polynomial evaluation of degree at
most $w$, choose integers $u,v$ such that $u2^w+v3^w=1$. For any
integer $h$, the polynomial $h(u+vX^w)$ satisfies
$H_w(h(u+vX^w))=h$. Lift all row coordinates this way.
The combined polynomial pair is nonzero because its first evaluation
is $DA\ne0$. Its two evaluations are divisible by $3^w2^w$, and its
divided remainder is $A\xi-B$. All the conditions of
Problem~\ref{long1049:prob:kernel}, with $F(3/2)$ replaced by $\xi$,
hold with $C=6$.

Even a coefficient-height bound of the form $\exp(O_\xi(n^2))$ would
not exclude these lifts. Choose $0\le v<2^w$, which gives $|u|<3^w$.
Since $|B|\le(n+1)|\xi|+1$, every lifted coefficient has absolute value
at most $(|\xi|+2)(n+1)18^w$.
Together with the preceding rationality contradiction, this proves that
the unrestricted construction problem for a real target $\xi$ is
equivalent to the irrationality of $\xi$. It does not establish either
assertion for $F(3/2)$.

The short paper records the same unrestricted construction requirements.
To pose a more specific question, one must replace the free choice of
polynomial pairs by an actual set given by coefficient formulas,
admissible parameters and permitted normalisations. A family name or
an unspecified exponential height bound is not such a restriction;
the lifts just constructed also have explicit formulas and a fixed
height bound. For a fixed indexed family, the distinction between all
sufficiently large indices and a sparse subsequence can matter.
Arbitrary padding and reindexing remove it in the unrestricted problem.


Corollary~\ref{long1049:res:nomult} gives no improvement from integer
rescaling and, under its endpoint hypotheses, no common factor $2$ or
$3$ in the unscaled evaluations. It makes no general assertion about
polynomial cancellation before specialisation, combinations of rows,
determinant-specific divisibility or other approximation families.

There are two separate tasks: find equal residues without obtaining a
zero polynomial pair or a zero remainder, then prove that the divided
remainder tends to zero. A coefficient-height estimate is needed when
the particular approximation argument calls for it; it is not an extra
hypothesis of the elementary irrationality criterion above.

\subsection{A congruence relation with nonzero remainder}

Fix, for each $n$, a common degree bound $W_n$ and a specified family
\[
 (U_{n,j},V_{n,j},\mathcal R_{n,j})_{j<M_n},
\]
where $U_{n,j},V_{n,j}\in\Z[X]$ have degree at most $W_n$ and
\[
 \mathcal R_{n,j}(t)=U_{n,j}(t)F(t)-V_{n,j}(t)\qquad(t>1).
\]
This identity uses the same normalisation as the evaluated rows and
all subsequent residue equations. Fix all divisions before forming the
residue map. Require each
polynomial pair to have coefficient gcd $1$, by dividing its common
integer coefficient factor if necessary. This is different from making
the two evaluated integers coprime, as required in
Problem~\ref{long1049:prob:kernel}. The two operations need not preserve
the same polynomial family. Coefficient content $1$ alone gives no
coprimality with $6$ after evaluation: at $W=1$, the pairs $(X,3)$
and $(2X-1,2)$ have coefficient content $1$ but give the integer rows
$(3,6)$ and $(4,4)$, respectively.

The gcd
\[
 \gcd\!\bigl(H_{W_n}(U_{n,j}),H_{W_n}(V_{n,j})\bigr)
\]
of one evaluated row can differ from the gcd of the final sum. Neither
is divided out without also dividing the corresponding remainder and
measuring the height after that division. If, in addition, the
coefficient of $X^{W_n}$ in $U_{n,j}$ and the constant coefficient of
$V_{n,j}$ are both units, Proposition~\ref{long1049:res:commonmult}
makes this row gcd coprime to $6$. These endpoint assumptions are
extra conditions, not consequences of coefficient primitivity.
Dividing by such a gcd preserves divisibility by every power of $2$
and $3$ for that individual row. It need not preserve a relation
with a fixed selector when different rows are divided by different
contents. For example, at $W=1$ the pairs $(X-1,X-1)$ and $(X+1,X+1)$
give rows $(1,1)$ and $(5,5)$. Their sum is zero modulo $6$, whereas
the sum of their primitive normalisations is $(2,2)$, not zero modulo
$6$. Both pairs satisfy the two unit conditions just stated. The
residue map must therefore be formed again after row-by-row normalisation.

Dividing a specialised row by its content need not preserve the chosen
polynomial family, and lifting the divided row back is a constrained problem. For fixed $W$, the map $P\mapsto H_W(P)$ on integer polynomials of
degree at most $W$ is surjective onto $\Z$, since $3^{W}$ and $2^{W}$ are
coprime, so a lift always exists. The issue is not an arbitrary exponential
height bound, as the construction above shows. It is whether the lift
belongs to the chosen approximation family and satisfies that family's
remainder identity and quantitative height estimate. With $W=1$,
for instance, $(X+1,X+1)$ specialises to $(5,5)$, whose primitive
normalisation $(1,1)$ lifts to $(X-1,X-1)$, but not to an integer scalar
multiple of the original pair.  Any lift used below is therefore supplied together with its
degree bound, its height bound and its exact remainder.

For target depths $R_n,S_n$, let $J_n(\lambda)$ be the vector of four residues of
the pair $\sum_j\lambda_j(U_{n,j},V_{n,j})$, and define
\[
 \mathcal C_n=
 \{\lambda\in\{-1,0,1\}^{M_n}\mathbin{\backslash}\{0\}:J_n(\lambda)=0\},
\]
\[
\begin{aligned}
 K_n^{\mathrm{poly}}
   &=\left\{\lambda\in\{-1,0,1\}^{M_n}:
       \sum_j\lambda_jU_{n,j}=0,\ \sum_j\lambda_jV_{n,j}=0\right\},\\
 K_n^{\mathrm{rem}}
   &=\left\{\lambda\in\{-1,0,1\}^{M_n}:
       \sum_j\lambda_j\mathcal R_{n,j}(3/2)=0\right\}.
\end{aligned}
\]
The second set consists of signed relations whose remainder vanishes
at $3/2$. It also contains every relation with an identically zero
remainder function. These are sets of restricted coefficient vectors,
not assertions that the signed cube is a vector space.

The exact pigeonhole condition is
\begin{equation}
 2^{M_n}>3^{2R_n}2^{2S_n},
 \qquad\text{equivalently}\qquad
 M_n>2R_n\log_2 3+2S_n.
\label{long1049:eq:exact-jet-threshold}
\end{equation}
Thus the least integer number of rows satisfying this counting test is
\[
 \left\lfloor2R\log_2 3+2S\right\rfloor+1.
\]
The checked condition $M\ge4R+2S$ for $R>0$ is a convenient sufficient
corollary, not the exact threshold. Applying Cauchy--Schwarz to the
sizes of the residue classes shows that there are at least
\begin{equation}
 \frac12\left(\frac{2^{2M_n}}{3^{2R_n}2^{2S_n}}-2^{M_n}\right)
\label{long1049:eq:collision-count}
\end{equation}
unordered pairs of distinct subsets with equal residues. It would
therefore suffice to prove that fewer than this many pairs give a zero
polynomial combination or a zero remainder at $3/2$.

Theorem~\ref{long1049:res:jetkernel} supplies only
$\mathcal C_n\ne\varnothing$; it does not rule out zero polynomial pairs
or zero real remainders. For a useful approximation family, bounding
the multiplicities in~\eqref{long1049:eq:collision-count} remains a
possible way to obtain nonvanishing. Without restrictions on the
family, however, nonvanishing alone is elementary.

\begin{example}[nonvanishing without decay]\label{long1049:prob:escape}
For $n\ge1$, set $W_n=R_n=S_n=n^2$ and $M_n=6n^2$. Take
\[
\begin{aligned}
 (U_{n,0},V_{n,0})&=((1+2^{W_n})X^{W_n},1),\\
 (U_{n,j},V_{n,j})&=(X^{W_n},1)\qquad(1\le j<M_n),
\end{aligned}
\]
and define $\mathcal R_{n,j}(t)=U_{n,j}(t)F(t)-V_{n,j}(t)$ for $t>1$.
All pairs have coefficient gcd $1$ and primitive evaluated rows. The
selector $(1,-1,0,\ldots,0)$ belongs to
$\mathcal C_n\mathbin{\backslash}
(K_n^{\mathrm{poly}}\cup K_n^{\mathrm{rem}})$, but the divided
integer form equals $F(3/2)$ for every $n$.
\end{example}

Indeed, the evaluated rows are
$(3^{W_n}+6^{W_n},2^{W_n})$ and $(3^{W_n},2^{W_n})$; their first
coordinates are odd, so both rows are primitive. Their difference is
$(6^{W_n},0)$, while the polynomial difference is
$(2^{W_n}X^{W_n},0)$. Thus all four residues vanish. The unscaled remainder difference
is $3^{W_n}F(3/2)>0$. Clearing the evaluation denominators multiplies
it by $2^{W_n}$, giving $6^{W_n}F(3/2)$. Dividing this cleared
integer form by $3^{R_n}2^{S_n}=6^{W_n}$ leaves the same positive
constant $F(3/2)$. The counting condition
\eqref{long1049:eq:exact-jet-threshold} holds because $\log_2 3<2$.

Example~\ref{long1049:prob:escape} meets the two nonvanishing
requirements in Problem~\ref{long1049:prob:kernel} without any
irrationality assumption. What it fails is the smallness condition.
The useful question is therefore to obtain nonvanishing and decay in
the same prescribed approximation family, not merely to exhibit some
family with a nonzero congruence relation.

\paragraph{Selector spans and multiplicities.}
The real-bin argument can be sharpened fibre by fibre. Fix a positive
integer $n$. For $M$ integer rows $(A_j,B_j)$ and a modulus $D\ge1$, set
$e_j=A_jF(3/2)-B_j$. For each attained residue vector $b$, let $T_b$
be the span of the selector remainders $\sum_j\varepsilon_je_j$
with $\varepsilon_j\in\{0,1\}$ and
$\sum_j\varepsilon_j(A_j,B_j)\equiv b\pmod D$.
Let the integer $k_b$ bound the number of selectors attaining any one exact real
value in that fibre. Then
\[
 2^M>\sum_b k_b\left(\left\lfloor\frac{nT_b}{D}\right\rfloor+1\right)
\]
produces two selectors with equal residues and distinct remainders
less than $D/n$ apart. Within each residue fibre, subtract the least
remainder, multiply by $n/D$, and take floors. The bin indices range
from $0$ to $\lfloor nT_b/D\rfloor$, including a separate final index
when the span is a positive integral multiple of $D/n$. Equal indices give a
remainder difference strictly less than $D/n$. If every bin contained
only one real value, its selector count would be at most $k_b$,
contradicting the displayed inequality. Subtracting the two selectors and dividing by $D$ gives
an integer form with nonzero absolute value less than $1/n$.
A common span $T=\sum_j|e_j|$ and multiplicity bound $k$ yield the
coarser count $Qk(\lfloor nT/D\rfloor+1)$, where $Q$ is the number of
attained residue vectors. All inputs must use the same row
normalisation. Primitive input rows can still have repeated subset
sums, and positivity need not survive subtraction.

\paragraph{The lattice after evaluation.}
For example, $(1,0),(1,6)$ generate $\Z\oplus6\Z$. They give two
residues modulo $2$. Within that lattice, the vectors whose two
coordinates are even form $2\Z\oplus6\Z$; dividing them by $2$ gives
$\Z\oplus3\Z$, of index $3$. Smith normal form separates the residue
count from this remaining index for any rank-two lattice.

Let primitive integer rows $u_j\in\Z^2$ span a rank-two lattice $L$, and
let $g$ be the gcd of their $2\times2$ minors. Since at least one row is
primitive, the Smith invariants are $1,g$
\cite[Thms.~2.3--2.4, p.~3]{stanley2016}. Hence, for $D\ge1$,
\[
 |\operatorname{im}(L\longrightarrow(\Z/D\Z)^2)|
 =\frac{D^2}{\gcd(g,D)},\qquad
 [\Z^2:(L\cap D\Z^2)/D]=\frac{g}{\gcd(g,D)}.
\]
Indeed an integral unimodular change of coordinates takes $L$ to
$\Z\times g\Z$; reduction modulo $D$ and intersection with $D\Z^2$
give the two formulas. This is Smith normal form over $\Z$, not over
$\Z[p]$. It measures the actual image, rather than the ambient residue space. If two independent divided rows have coordinate height at most
$H$ and remainders of absolute value at most $\varepsilon$, their nonzero
integer determinant gives the necessary inequality
$2H\varepsilon\ge g/\gcd(g,D)$. That is not a necessary condition for
producing a single nonzero form.
If the specialised rows are not primitive, let $d_1\mid d_2$ be the
positive Smith invariants instead. For $D\ge1$ the corresponding formulas are
\[
 |\operatorname{im}(L\bmod D)|=
 \frac{D^2}{\gcd(d_1,D)\gcd(d_2,D)},\qquad
 [\Z^2:(L\cap D\Z^2)/D]=
 \frac{d_1d_2}{\gcd(d_1,D)\gcd(d_2,D)}.
\]
They follow by applying the one-dimensional calculation to each summand
$d_i\Z$. Polynomial coefficient content alone does not determine either
integer Smith invariant after specialisation.

\paragraph{The lattice spanned by the actual tails.}
The formulas above apply to any rank-two lattice. For the rows that the
approximation problem actually supplies, the invariant can be computed
outright, and it is the smallest one compatible with the base. Let $a>b\ge1$ be
coprime, write the partial sum $S_m(b/a)=\sum_{r\le m}(b/a)^r/(1-(b/a)^r)$ in
lowest terms as $P_m/Q_m$ with $Q_m>0$, and set $P_0=0$, $Q_0=1$. The primitive
integer row of the $m$th tail is $(Q_m,P_m)$, which represents
$Q_mF(a/b)-P_m$.

\begin{proposition}[the prefix lattice of the tails]
\label{long1049:res:tail-lattice}
Every $Q_m$ is coprime to $ab$, and $b$ divides every $P_m$. For every prefix
containing $m=0$ and $m=1$,
\[
 \operatorname{span}_{\Z}\{(Q_m,P_m):0\le m<M\}=\Z\times b\Z
 \qquad(M\ge2),
\]
so the Smith invariants are $1,b$, the gcd of the $2\times2$ minors is $b$, and
the image modulo $D$ has cardinality $D^2/\gcd(b,D)$. Multiplying an individual
tail by a nonzero rational weight leaves its primitive row unchanged up to
sign.
\end{proposition}

\begin{proof}
A common denominator of $S_m(b/a)$ is a product of the integers $a^r-b^r$ for
$r\le m$, each coprime to $ab$; over that denominator every summand has
numerator divisible by $b^r$. Reduction to lowest terms removes no prime factor
of $b$, which gives the first two assertions. The rows at $m=0,1$ are $(1,0)$
and $(a-b,b)$, and subtracting $a-b$ copies of the first from the second gives
$(0,b)$; every later row already lies in $\Z\times b\Z$ by the first
assertions. The Smith and image statements are the specialisation of the
formulas above at $g=b$. Finally the rational line through a primitive pair
meets $\Z^2$ in exactly the multiples of that pair, so clearing denominators in
a nonzero rational multiple and dividing by the coordinate gcd returns the same
pair up to sign.
\end{proof}

At $3/2$ the first three partial sums are $0$, $2$ and $14/5$, with rows
$(1,0)$, $(1,2)$ and $(5,14)$. Their span is already $\Z\times2\Z$. There is no
three-adic collapse, and the only local saving is the single factor of two that
the image formula displays. For a window beginning at $m=s$ the corresponding
statement is that every pairwise minor is divisible by $b^{s+1}$, with the
valuation of the gcd equal to $(s+1)v_p(b)$ at each prime $p\mid b$ and equal
to zero at each prime $p\mid a$. The adjacent-row identity
\[
 Q_sP_{s+1}-P_sQ_{s+1}
 =\frac{Q_sQ_{s+1}b^{s+1}}{a^{s+1}-b^{s+1}}
\]
attains both, since its denominator is a unit at every prime dividing $ab$.

The consequence for the construction problem is a restriction on one
operation, and only on that one. Rescaling the individual tails and then taking
integer combinations stays inside $\Z\times b\Z$, so it leaves the minor gcd
and the endpoint valuations exactly as above. Rational combinations are a
different matter and do leave that lattice: at $3/2$ the half-sum of the rows
$(1,0)$ and $(1,2)$ is the primitive row $(1,1)$, whose minor against $(1,0)$
is $1$. What such a combination costs is a denominator in the coefficients,
which the proposition does not measure, so escaping the lattice this way is not
by itself a saving. A genuinely improved system needs primitive rows of a
different provenance, whose minors carry extra valuation and whose divided
remainders are still nonzero and small.



\subsection{Estimating the divided remainder}

The next displayed margin is an additional research target for a specified
comparison of height and remainder, not a necessary irrationality criterion.
First fix such a construction and define its height $H_n\ge1$, undivided
nonzero remainder $L_n$ and exactly once-counted certified divisor
$D_n=3^{R_n}2^{S_n}$. If another divisor is used, replace the two logarithmic
terms below by $\log D_n$. A scalar-form height and an exterior-determinant
height cannot be interchanged: the comparison with a nonzero integer must
be derived for the actual objects selected. For a scalar construction,
integral coefficients and $0<|L_n|/D_n\to0$ already suffice, with no
extra height factor. Conversely, the proposed negative margin would
imply this scalar decay because $H_n\ge1$. Allowing an arbitrarily
small positive $H_n$ would lose that implication.

\begin{problem}[an additional height and remainder estimate]
Prove the explicit estimate
\begin{equation}
 \limsup_{n\to\infty}
 \frac{\log H_n+\log|L_n|-R_n\log3-S_n\log2}{n^2}<0.
\label{long1049:eq:negative-margin}
\end{equation}
Every denominator, row content and final-combination content must already be
included in $H_n$ and $L_n$.
\end{problem}

Nonvanishing alone does not address~\eqref{long1049:eq:negative-margin}.  Conversely,
a formal decay estimate cannot supply a nonzero form if every
combination with the required residues has zero remainder.  A proposed divisor saving must be compared with the height and remainder
of the same explicitly normalised objects.

\subsection{A restriction on Mahler functional equations}

Mahler's method requires suitable functional equations. For the divisor
generating series $\mathcal L(z)=\sum_{n\ge1}\tau(n)z^n$, Bell and
Smertnig's classification rules out a $k$-Mahler equation for every $k\ge2$
\cite[Thm.~1.3 and the consequences on p.~3]{bellsmertnig2026}.  The
proposition below proves the simultaneous $2$/$3$ case using the theorem
of Adamczewski and Bell and the functional nonrationality already
proved in Section~\ref{long1049:sec:source-forms}.
The 2026 works are cited as the identified preprints, not as journal
publications. A function-level obstruction to these functional equations
does not determine the arithmetic nature of any one rational-base value.

To apply the classification, note that $\tau$ is multiplicative:
for coprime integers $m,n$, each divisor of $mn$ has a unique
factorisation into a divisor of $m$ and a divisor of $n$.
Suppose that $\mathcal L$ were $k$-Mahler. The classification would give a prime $p$, an integer
$r\ge0$ and an eventually periodic function $\chi$ with
$\tau(m)=m^r\chi(m)$ whenever $p\nmid m$. For a prime $\ell\ne p$,
this forces
\[
 \chi(\ell^j)=\frac{j+1}{\ell^{jr}}\qquad(j\ge0).
\]
If $r=0$, these values are unbounded. If $r>0$, they are nonzero
and tend to zero. Both alternatives contradict the finite range of
an eventually periodic function.


Here the ambient space is $\Q((z))$, the field of formal Laurent
series, viewed as a vector space over $\Q(z)$. Stability means that
substituting $z^k$ for $z$ sends each member of the subspace back into
that subspace; this substitution is not a $\Q(z)$-linear map.

\begin{proposition}[no finite simultaneous $2/3$-system]\label{long1049:res:nomahler}
Let
\[
 \mathcal L(z)=\sum_{n\ge1}\frac{z^n}{1-z^n}.
\]
There is no finite-dimensional $\Q(z)$-vector space that contains $\mathcal L$
and is stable under both $z\mapsto z^2$ and $z\mapsto z^3$.
\end{proposition}
\leannote{Lean: \lproof{ErdosProblems/Erdos1049/PaperCompleteR21/SimultaneousMahlerSystem.lean}{1140}{no_finite_simultaneous_two_three_system}, \lproof{ErdosProblems/Erdos1049/PaperCompleteR21/SimultaneousMahlerSystem.lean}{228}{mahler_of_stable}, \lproof{ErdosProblems/Erdos1049/PaperCompleteR21/SimultaneousMahlerSystem.lean}{203}{isMahler_subs}, \lproof{ErdosProblems/Erdos1049/PaperCompleteR21/SimultaneousMahlerSystem.lean}{909}{divisorLambert_subs_not_rational}, and 8 further declarations. Conditional on the rationality theorem of Adamczewski and Bell; see the \hyperref[long1049:sec:coverage]{coverage section}.}

\begin{proof}
Suppose $V$ were such a space, of dimension $d$.  Stability under $z\mapsto z^2$
places the $d+1$ elements $\mathcal L(z),\mathcal L(z^2),\dots,\mathcal
L(z^{2^{d}})$ in $V$, so they are linearly dependent over $\Q(z)$; clearing
denominators gives polynomials $P_0,\dots,P_d$, not all zero, with
$\sum_{i=0}^{d}P_i(z)\mathcal L(z^{2^i})=0$.

A Mahler equation requires a nonzero coefficient of the unshifted
function. To obtain one here, let $j$ be the least index with
$P_j\ne0$. Write each polynomial uniquely as
$P_i(z)=\sum_{r=0}^{2^j-1}z^rQ_{i,r}(z^{2^j})$.
Every series $\mathcal L(z^{2^i})$ with $i\ge j$ has exponents
divisible by $2^j$, so the relation splits by exponent residues modulo
$2^j$. Choose $r$ with $Q_{j,r}\ne0$ and put $u=z^{2^j}$. The
corresponding relation is
\[
 \sum_{i=j}^{d}Q_{i,r}(u)\mathcal L(u^{2^{i-j}})=0,
\]
with a nonzero coefficient of $\mathcal L(u)$. Thus $\mathcal L$ is
$2$-Mahler. The same argument with $3$ in place of $2$ makes it
$3$-Mahler. Since $2$ and
$3$ are multiplicatively independent, a theorem of Adamczewski and Bell
\cite[Thm.~1.1, p.~6]{adamczewskibell2013} then forces $\mathcal L$ to be a rational
function. Since $\mathcal L(z)=F(1/z)$ for $0<z<1$, this contradicts
the nonrationality of $F$ proved in Section~\ref{long1049:sec:source-forms}.
Rivin's periodic-coefficient corollary
\cite[Cor.~6.4, p.~9]{rivin2026} gives the same nonrationality conclusion.
\end{proof}

Proposition~\ref{long1049:res:nomahler} uses no property of the point
$2/3$: the obstruction is functional and appears before regularity at a particular point is considered.
For this scalar function, Bell and Smertnig's single-base classification
already implies the proposition. The proof above instead derives it
from simultaneous closure and elementary functional nonrationality,
using the Adamczewski--Bell theorem.  A
construction using additional functions or functional relations must specify
those functions and its closure conditions; the single-base statement is not
an obstruction to every approximation method. The classification and the
Adamczewski--Bell theorem are cited, not proved here; the needed
nonrationality has the elementary proof given earlier.

The Lean proof of Proposition~\ref{long1049:res:nomahler}
(\href{\laterepobase/lean/ErdosProblems/Erdos1049/PaperCompleteR21/SimultaneousMahlerSystemUnconditional.lean\#L298}{no finite simultaneous $2/3$-system})
uses neither cited theorem, and it uses only stability under
$z\mapsto z^2$. As in the first step of the proof above, that stability
gives polynomials $P_0,\dots,P_d$, not all zero, with
$\sum_{i=0}^{d}P_i(z)\mathcal L(z^{2^i})=0$, and this relation also holds
between the values of these functions on the open unit disc. If $\zeta$ is
a primitive $n$th root of unity, then
$(1-r)\mathcal L(r\zeta)/\log(1/(1-r))\to1/n$ as $r\uparrow1$.
These limits, taken at roots of unity of suitable orders, force every
$P_i$ to vanish. The same argument excludes, for every $k\ge2$, a
finite-dimensional space containing $\mathcal L$ and stable under
$z\mapsto z^k$ alone
(\href{\laterepobase/lean/ErdosProblems/Erdos1049/PaperCompleteR21/SimultaneousMahlerSystemUnconditional.lean\#L263}{single-base statement}).

\subsection{A limitation of the rectangular exponent model}

Consider the two-parameter expression $\Theta_{\mathrm{HP}}$ defined in
Section~\ref{long1049:sec:sharp}. On the domain $\rho\ge0$,
$\sigma\ge1+\rho$, its denominator
\[
 \frac{(1+\rho)^2}{2}+\sigma(1+\rho)+\frac{1+\rho^2}{2}+\sigma
\]
is positive. Put $u=\sigma-1-\rho\ge0$. Multiplying the difference
$\Theta_{\mathrm{HP}}(\rho,\sigma)-(1/2-1/\pi^2)$ by
$2\pi^2$ times this denominator gives exactly
\[
 -\pi^2\rho^2-\pi^2\rho u-2\pi^2\rho-2\rho^2-10\rho u
 -4\rho-6u^2-8u.
\]
Every term is nonpositive. The original difference has the same sign, and
it vanishes exactly when $\rho=0$ and $\sigma=1$
(\lword{Erdos1049/HermitePadeNoGo.lean}{48}{hpClearedGap_expansion}{exact expansion},
\lword{Erdos1049/HermitePadeNoGo.lean}{58}{hpClearedGap_nonpos}{nonpositivity},
\lword{Erdos1049/HermitePadeNoGo.lean}{75}{hpClearedGap_eq_zero_iff}{equality case}).
Equivalently, within that model the displayed threshold never exceeds the
classical one-function margin, with equality only at the classical endpoint
(\lword{Erdos1049/HermitePadeNoGo.lean}{103}{rectangular_hp_threshold_le_classical}{bound},
\lword{Erdos1049/HermitePadeNoGo.lean}{126}{rectangular_hp_threshold_eq_classical_iff}{equality}).
This bounds only the displayed exponent model. Applying it to an
approximation family would require the polynomial construction,
integrality and asymptotic estimates connecting that family to
$\Theta_{\mathrm{HP}}$.
The rational cutoff $81/200$ discussed above satisfies
\[
 \frac{81}{200}<\frac{\log2}{\log3}.
\]
For a sufficient criterion of the form $\log b/\log a<T$, including
$3/2$ requires $T>\log2/\log3$, not merely an improvement on
$81/200$. This numerical target applies to that form of criterion,
not to every irrationality argument.

A concrete optimisation question is available within the 2004
construction itself. Write a source parameter direction as
$\alpha=(\alpha_0,\alpha_1,\alpha_2;\beta)$, with positive integer entries satisfying
\[
 \alpha_1\le\alpha_2,\qquad
 \alpha_1+\alpha_2<\beta\le\alpha_0+\alpha_2,\qquad
 \gcd(\alpha_0,\alpha_1,\alpha_2,\beta)=1.
\]
The source parameters are $a_j=\alpha_jn+1$ and $b=\beta n+2$;
here $b$ is not a rational-base denominator. The gcd condition is a
normalisation of the parameter direction, not a primitivity assertion
about evaluated rows. Its invariance under dilation is checked below.
For these directions set
\[
 \begin{gathered}
 c_{00}=\alpha_0+\alpha_1+\alpha_2-\beta,\quad
 c_{01}=\alpha_0,\quad c_{11}=\alpha_1,\quad c_{21}=\alpha_2,\\
 c_{12}=\beta-\alpha_1,\quad c_{22}=\beta-\alpha_2,\qquad
 m=\max(c_{00},c_{01},c_{11},c_{21},c_{12},c_{22}),
 \end{gathered}
\]
and define the periodic step function
\[
 \omega_\alpha(u)=\max\left\{\begin{aligned}
 &0,\\[-2pt]
 &\lfloor c_{21}u\rfloor+\lfloor c_{22}u\rfloor
       -\lfloor c_{11}u\rfloor-\lfloor c_{12}u\rfloor,\\[-2pt]
 &\lfloor c_{01}u\rfloor+\lfloor c_{21}u\rfloor
       -\lfloor c_{00}u\rfloor-\lfloor c_{12}u\rfloor
 \end{aligned}\right\}.
\]
Zudilin's constants in (25) and (26), already used in the finite
direction scan, are
\[
 \begin{aligned}
 C_1(\alpha)&=(\alpha_0+\alpha_1+\alpha_2)\beta
               -\frac{\alpha_1^2+\alpha_2^2+\beta^2}{2},\\
 C_0(\alpha)&=\frac{\alpha_1^2}{2}+\alpha_0\alpha_1
            +(\beta-\alpha_2)(\alpha_2-\alpha_1)\\
 &\hspace{8mm}-\frac3{\pi^2}
       \left(m^2-\int_0^1\omega_\alpha(u)\,d(-\psi_1(u))\right).
 \end{aligned}
\]
Here $\psi_1$ is the trigamma function used earlier. The step function
is bounded, nonnegative, periodic with period $1$, and zero near zero,
so the integral is finite.

The gcd normalisation does not change the ratio being optimised.
Indeed, periodicity and
$-\psi_1'(u)=2\sum_{j\ge0}(u+j)^{-3}$ give, by Tonelli's theorem,
\[
 \int_0^1\omega_\alpha(u)\,d(-\psi_1(u))
   =\int_0^\infty\frac{2\omega_\alpha(u)}{u^3}\,du.
\]
For any positive integer $k$, the floor formulas give
$\omega_{k\alpha}(u)=\omega_\alpha(ku)$. Substitution in the last integral
therefore multiplies it by $k^2$. All other terms in $C_0$ and $C_1$
are quadratic in the direction, including $m^2$. Hence
$C_i(k\alpha)=k^2C_i(\alpha)$ for $i=0,1$, so passing to a primitive
direction preserves both positivity and $C_0/C_1$.
Let $\mathcal A$ be exactly the displayed primitive directions with
$C_0(\alpha)>0$ and $C_1(\alpha)>0$.

\begin{problem}[optimising the published parameter directions]
Determine
\[
 \sup_{\alpha\in\mathcal A}\frac{C_0(\alpha)}{C_1(\alpha)},
\]
or improve its bounds. In particular, does an admissible direction give
a ratio strictly greater than the value $\theta^*$ in
Theorem~\ref{long1049:res:region}, attained at $(14,12,14;27)$?
A finite numerical scan does not establish optimality over $\mathcal A$.
\end{problem}

The immediate bounds are $\theta^*\le\sup_{\mathcal A}C_0/C_1\le1/2$.
Indeed, at any integer base $p\ge2$ the source gives
$\mu_{\rm irr}(F(p))\le C_1(\alpha)/C_0(\alpha)$, whereas pigeonhole
approximation gives $\mu_{\rm irr}(F(p))\ge2$.
Thus even the optimal reciprocal criterion in this class cannot exceed
$1/2$, which is below $\log2/\log3$. Optimising these directions cannot
include $3/2$ by that criterion. For other noninteger bases, a new
direction would still require the polynomial degree and coefficient-height
estimates used in the rational-base transfer; the integer-base source
alone does not supply that application.
The unrestricted question of which rational bases give irrational values
is not settled here and is not reduced to any one of these problems.

\section*{Related work and comparison of constructions}
The following comparisons place the two constructions in their
literature. The distinction between an integer base, a rational base and
a statement about a function is essential in each comparison.

\paragraph{Functional nonrationality is not value irrationality.}
At the level of functions, Rivin proves that if a sequence \(\gamma\) and its
divisor-sum sequence are both eventually linearly recurrent, then \(\gamma\)
is finitely supported
\cite[Theorem~1.1, p.~2; proof pp.~6--7]{rivin2026}.
His periodic-coefficient corollary therefore shows that
\[
  \sum_{n\ge1}\frac{z^n}{1-z^n}
\]
is not a rational function~\cite[Corollary~6.4, p.~9]{rivin2026}.
This is an exact structural statement about the function underlying
Problem~\#1049, but it gives no irrationality statement for a special value
at \(z=1/t\): a nonrational function may take rational values at particular
rational points.

\paragraph{Pad\'e and orthogonal-polynomial constructions.}
In 1991
Borwein proved the irrationality of shifted series
$\sum_{n\ge1}(t^{n}+w)^{-1}$ at integer bases $t\ge2$, for every nonzero
rational $w$ with $w\ne-t^m$ for all $m\ge1$, by Pad\'e approximation rather
than by digit clearing; his estimates also show that these values are not
Liouville numbers~\cite[Thm.~4, pp.~257--258]{borwein1991}.  In 2001 Van
Assche recovered the integer-base irrationality and the bound
$\mu_{\rm irr}(F(p))\le 2\pi^2/(\pi^2-2)=2.50828\ldots$ using little $q$-Legendre
Pad\'e approximants
~\cite[Thm.~1, p.~10; proof pp.~10--11]{vanassche2001}.  His more general
Theorem~3 proves irrationality of
$\sum_{k\ge1}(cp^k-1)^{-1}$ for an integer $p>1$ and a fixed nonzero
rational $c$ with $cp^k\ne1$ for every $k\ge1$
~\cite[Thm.~3, p.~14]{vanassche2001}; it does not cover a
rational noninteger base $t$ in $F(t)$, because the multiplier needed to
write $t^k$ over an integer base varies with $k$.

\paragraph{Related series and their different recurrences.}
The same Lambert value was already the target of Amdeberhan and
Zeilberger's $q$-WZ construction~\cite{az1998}.  Here $p>1$ is the integer-base parameter, $q=p^{-1}$, and
$P_n(x\mid q)$ denotes the little $q$-Legendre polynomial used by these
authors, normalised by $P_n(0\mid q)=1$.  The two constructions
share that bivariate little-$q$-Legendre Pad\'e kernel, but Van Assche's
diagonal does not satisfy the Amdeberhan--Zeilberger scalar recurrence.
Their Theorems~1 and~2 do prove, respectively, the irrationality of the
non-alternating $h_p(1)$ and alternating $q$-logarithm values for their stated
integer parameters, with reported irrationality measure $4.80$; those are
external integer-parameter results and do not transfer to the rational
noninteger base $3/2$ studied here.

Amdeberhan--Zeilberger use the moving diagonal
$P_n(p^{n+1}\mid p^{-1})$, whereas Van Assche uses
$P_n(p^n\mid p^{-1})$.  Van Assche also records, citing
Borwein's 1992 Lemma~2, the neighbouring evaluation
$P_{n-1}(c p^{n+1}\mid p^{-1})$
\cite[(23) and the following paragraph, p.~6]{vanassche2001};
see \cite[Lemma~2, p.~143]{borwein1992}.
The shared kernel therefore does not license transfer of recurrence,
endpoint, lattice, or valuation claims between the diagonals.  Indeed, if
$A_n(p)=P_n(p^n\mid p^{-1})$, exact substitution at $n=0$ into the
Amdeberhan--Zeilberger operator \cite[Sec.~1.5, p.~2]{az1998} leaves
\[
 -p(p-1)^2(p+1)(p^5+2p^4+2p^3+2p^2+2),
\]
which is nonzero for every real $p>1$.  One nonzero residual is decisive for
non-transfer of that recurrence.

Lean checks the
\lword{Erdos1049/QAperyDiagonalNonEquivalence.lean}{67}{vanAsscheQAperyResidualAtZero_eq}{exact
residual factorisation} and its
\lword{Erdos1049/QAperyDiagonalNonEquivalence.lean}{94}{vanAsscheQAperyResidualAtZero_ne_zero}{nonvanishing
for $p>1$}.  These are finite statements at $n=0$; they do not supply a
recurrence for either diagonal.  No recurrence, endpoint, lattice, or
valuation statement for one diagonal is used for the other; the displayed
exact residual is the sole claim made here about their incompatibility.

Coussement and Smet use little $q$-Jacobi polynomials for further Lambert
series whose base parameter is the reciprocal of an integer
\cite[Thms.~1.2--1.3, p.~2; Sec.~2]{coussementsmet2007}, and Koizumi and
Yokoi identify one of their three-parameter Ap\'ery-type approximations with
the Coussement--Smet Pad\'e approximants
\cite[Sec.~7, Prop.~7.1, p.~33]{koizumiyokoi2026}.
That identification includes the source's moving evaluation point
$z_n=p_2^{n+1}$ and its normalization; it is not an identification of
all little-$q$-Jacobi diagonals or their scalar recurrences.

At a noninteger rational base, convergence of the rational approximants
is not enough. If $P_n/Q_n\to\xi$, the cleared error is
$Q_n\xi-P_n=Q_n(\xi-P_n/Q_n)$, whose size also depends on $Q_n$.
For example, $1/n\to0$, but the corresponding cleared error is always
$-1$. Thus convergence alone does not give the small integer linear
forms required for irrationality.

In 2013 Vandehey proved that
$\sum_{n\ge1}d(n)a_n/b^n$ is irrational whenever $b>1$ is an integer and
$(a_n)$ ranges in a finite integer alphabet excluding zero; taking
$a_n=(-1)^n$ completes the elementary digit method for integer bases
$b\le-2$~\cite[Thm.~1.2, p.~2]{vandehey2013}.  His companion theorem permits
a finite alphabet of nonnegative integers containing zero, provided the
coefficient sequence is not eventually zero
\cite[Thm.~1.1, p.~2]{vandehey2013}.  All
of these statements require an integer base. Luca and Tachiya's periodic-coefficient theorem
likewise proves irrationality at every integer base of absolute value greater
than one; their full-support example strengthens this to joint linear
independence for finite families of iterated divisor functions
\cite[Theorem~A, p.~139; Example~1, p.~140]{lucatachiya2017}.  Duverney and
Tachiya refine the Chowla--Erd\H{o}s method to linear independence results for
Lambert series at an integer base $q$ with $|q|>1$
\cite[Thms.~1.1--1.2, pp.~2--3]{duverneytachiya2019}; their Corollary~1.1
generalises Vandehey's Theorem~1.2 \cite[p.~3]{duverneytachiya2019}.  Each of these results assumes an integer base.
For the underlying small-integral-tail criterion, see also Duverney
\cite[Thm.~2, p.~4 of the author manuscript]{duverney2011}.
We use that elementary criterion only as methodological context, not his
separate square-index linear-independence argument.


In 2004 Zudilin obtained the uniform bound
$\mu_{\rm irr}(F(p))\le2.46497868\ldots$ for every integer $p\ge2$,
by way of Heine's basic transform and a permutation group
~\cite[Thm.~1, p.~154; Secs.~4--5, pp.~159--162]{zudilin2004}.
His source theorem also covers negative integer inverse bases; our
notation $F(t)$ and positive-remainder arguments are restricted to $t>1$.
The ordinary-hypergeometric antecedent is Rhin and Viola's $S_5$ action and
twelve-coset denominator reduction for rational forms in
$\zeta(2)$~\cite[Sec.~3, pp.~38--42; Sec.~4, pp.~46--51]{rhinviola1996}.

These sources suggest looking for common polynomial factors before
specialisation. The congruence lemmas in
Section~\ref{long1049:sec:endpoints} describe what the first and last
coefficients can imply at $3/2$. Their hypotheses have not been proved
there for either cited coefficient family.

The neighbouring problem of Lambert
subseries $\sum_{n\in A}(t^{n}-1)^{-1}$ over a restricted index set $A$ is
treated by Kova\v{c} and Tao~\cite[Sec.~2.1.2 and Thm.~2.3, pp.~4--5]{kovactao2024}.

The rational non-integer
progress relevant here is instead the height criterion of Bundschuh and
V\"a\"an\"anen~\cite[Thm.~2, p.~177]{bv1994}, which gives the strongest explicit
numerical threshold for this value among the previously published sources
compared here; Theorem~\ref{long1049:res:region} enlarges the region that
criterion covers.  Zudilin's later Pad\'e-and-Hankel treatment of the generalized
$q$-logarithm also comments on non-integral rational bases: it remarks that its results can be
given for non-integer $p=r/s$ with $|p|>1$, under an assumption
$\log|r|>c\log|s|$ for some computable
$c>0$~\cite[Sec.~2, p.~4]{zudilin2016}.  No value of $c$ is computed there.  That remark
announces the shape of a rational-base extension for the generalized
$q$-logarithm of that paper.  Section~\ref{long1049:sec:region} proves an explicit extension for
$F$ using the different, 2004 construction. Its polynomial cancellation
and denominator calculation show that
$c=\mu=2.464978683574975\ldots$, the 2004 bound quoted above, is admissible,
and gives the resulting region and its application to $31/4$.

A third published rational-base region for this same series is Duverney's
Th\'eor\`eme~2, whose hypothesis is
$\log|s|/\log|r|<\frac13\bigl(1-3/\pi^{2}\bigr)=0.2320\ldots$
\cite[Th\'eor\`eme~2, p.~174]{duverney1996}; his Th\'eor\`eme~1 for general
numerators is weaker still.  Both cutoffs lie strictly inside the Bundschuh and
V\"a\"an\"anen region.  Matala-aho, V\"a\"an\"anen and Zudilin's combined
treatment of $q$-logarithms keeps the integer hypothesis $p=1/q\in\Z\mathbin{\backslash}\{0,\pm1\}$
throughout, and states in print that its methods do not sharpen the
$q$-harmonic case of~\cite{zudilin2004}
\cite[abstract p.~879; introduction p.~880]{matalaaho2006}.  So none of the
printed regions for $F$ compared here reaches $\log4/\log31$.


\pdfbookmark[1]{Statements and declarations}{statements-declarations}
\subsection*{Statements and declarations}

\paragraph{Artefact and data availability.} The
\href{\sourceurl}{pinned formal-source revision} contains the Lean sources, the
fixed toolchain, and the library manifest used in the verification.
Formal-verification claims refer to those pinned sources. The mathematical
proofs are given in the text; ordinary arguments and finite computations
are identified separately.

\paragraph{Funding and competing interests.} This work received no external
funding.  The author declares no competing interests.

\paragraph{Acknowledgements.} The problem numbering and status follow the
Erd\H{o}s Problems catalogue maintained by Thomas Bloom~\cite{erdosproblems}.
I thank Wouter van Doorn for advice on writing for a first-time reader,
using fewer names and symbols, and explaining the force of a theorem's
hypotheses. His advice was given on another note and does not constitute
mathematical review or endorsement of the results here.

\appendix
\section{Guide to the formal sources}\label{long1049:app:index}

\subsection*{Formalisation coverage and remaining dependencies}
\label{long1049:sec:coverage}
\papersectiontarget{coverage}

Every theorem, lemma, proposition and corollary of this record other than the
six named here has a Lean statement of the same assertion, with the same
hypotheses, checked by the Lean kernel using only the axioms
\texttt{propext}, \texttt{Classical.choice} and \texttt{Quot.sound}, in the
development at revision \texttt{181078b6b009}.  Two of the six,
Propositions~\ref{long1049:res:nomahler} and~\ref{long1049:res:finite-pencil},
were proved later: each has a Lean statement of the same assertion, checked
by the Lean kernel using only those three axioms, at revision
\texttt{e4ba6841dfdb}, as described below.  The other four are proved by
ordinary arguments and have no Lean statement at that revision:
Theorem~\ref{long1049:thm:geometric-universality},
Proposition~\ref{long1049:prop:rogers-factorisation} and
Theorem~\ref{long1049:thm:sharp-fixed-base} in
Section~\ref{long1049:sec:sharp-fixed-base}, and
Proposition~\ref{long1049:res:tail-lattice}.

Proposition~\ref{long1049:res:nomahler} is
\href{\laterepobase/lean/ErdosProblems/Erdos1049/PaperCompleteR21/SimultaneousMahlerSystemUnconditional.lean\#L298}{checked in Lean}
with no additional hypothesis.  Its Lean proof uses stability under
$z\mapsto z^2$ alone, without the theorem of Adamczewski and
Bell~\cite[Thm.~1.1, p.~6]{adamczewskibell2013} cited in the printed proof,
as explained after that proof.

Proposition~\ref{long1049:res:finite-pencil} is
\href{\laterepobase/lean/ErdosProblems/Erdos1049/PaperR20/FinitePencilProposition.lean\#L43}{checked in Lean}
for every real $p>1$ and every rank $N\le8$.  The positivity of the eight
unshifted determinants $D_{k,0}$ is proved in Lean by the certificate
described in Section~\ref{long1049:sec:coefficient-questions}.  The identity
$F(p)A_N-B_N=(v^*_{i+j})$ and the positive definiteness of this matrix, the
reality of the pencil roots, their position below $F(p)$ and their
interlacing are formalised as well.

% BEGIN GENERATED CONCORDANCE (claude_lane/build_concordance.py; do not edit by hand)
\paragraph{Concordance of statements and Lean declarations.}
Each result of this record that has a kernel-checked Lean statement of the same assertion is listed
below with the declarations that jointly state it.  Each name links to its declaration at revision
\texttt{181078b6b009}.  Where the Lean statement is stronger than the printed one and implies it by
an immediate specialisation, the entry says so.  An entry marked \emph{compared} was also checked
independently: a restatement of the same declarations against Mathlib alone, together with its
proof, was verified by Comparator (\texttt{leanprover/comparator}) in a clean continuous
integration environment, in the run named by its number.  Comparator trusts the restated
statement, so the correspondence between the printed result and that statement is the one this
concordance records.

\begingroup\small\raggedright
\par\noindent\hangindent=1.5em Lemma~\ref{long1049:res:omega-indicator}: \href{https://github.com/wcook04/plectis-erdos/blob/181078b6b009d905809cf2e007ad309386a3d1e0/lean/ErdosProblems/Erdos1049/PaperOmegaIndicatorR7.lean\#L1297}{\nolinkurl{omega_indicator}}.  \emph{Compared}, run \href{https://github.com/wcook04/plectis-erdos-lean/actions/runs/35624228171}{\texttt{35624228171}}.
\par\noindent\hangindent=1.5em Theorem~\ref{long1049:res:region}: \href{https://github.com/wcook04/plectis-erdos/blob/181078b6b009d905809cf2e007ad309386a3d1e0/lean/ErdosProblems/Erdos1049/PaperCompleteR21/PrintedContourConstants.lean\#L418}{\nolinkurl{printed_contour}}, \href{https://github.com/wcook04/plectis-erdos/blob/181078b6b009d905809cf2e007ad309386a3d1e0/lean/ErdosProblems/Erdos1049/PaperCompleteR21/PrintedContourConstants.lean\#L487}{\nolinkurl{contour_enclosure}}, \href{https://github.com/wcook04/plectis-erdos/blob/181078b6b009d905809cf2e007ad309386a3d1e0/lean/ErdosProblems/Erdos1049/PaperCompleteR21/PrintedContourConstants.lean\#L334}{\nolinkurl{zudilinJ_enclosure}}, \href{https://github.com/wcook04/plectis-erdos/blob/181078b6b009d905809cf2e007ad309386a3d1e0/lean/ErdosProblems/Erdos1049/PaperCompleteR21/PrintedContourConstants.lean\#L383}{\nolinkurl{zudilinC0_enclosure}}, \href{https://github.com/wcook04/plectis-erdos/blob/181078b6b009d905809cf2e007ad309386a3d1e0/lean/ErdosProblems/Erdos1049/PaperCompleteR21/RationalBaseThreshold.lean\#L96}{\nolinkurl{zudilin_rpow_lt_iff_contourRegion}}, \href{https://github.com/wcook04/plectis-erdos/blob/181078b6b009d905809cf2e007ad309386a3d1e0/lean/ErdosProblems/Erdos1049/PaperCompleteR21/RationalBaseThreshold.lean\#L109}{\nolinkurl{rational_base_threshold}}, \href{https://github.com/wcook04/plectis-erdos/blob/181078b6b009d905809cf2e007ad309386a3d1e0/lean/ErdosProblems/Erdos1049/PaperCompleteR21/RationalBaseThreshold.lean\#L116}{\nolinkurl{rational_base_threshold_log}}.  \emph{Compared}, runs \href{https://github.com/wcook04/plectis-erdos-lean/actions/runs/35624228171}{\texttt{35624228171}}, \href{https://github.com/wcook04/plectis-erdos-lean/actions/runs/35674034595}{\texttt{35674034595}}.
\par\noindent\hangindent=1.5em Theorem~\ref{long1049:res:31over4}: \href{https://github.com/wcook04/plectis-erdos/blob/181078b6b009d905809cf2e007ad309386a3d1e0/lean/ErdosProblems/Erdos1049/PaperCompleteR21/PrintedLogConstants.lean\#L118}{\nolinkurl{printed_log_ratio}}, \href{https://github.com/wcook04/plectis-erdos/blob/181078b6b009d905809cf2e007ad309386a3d1e0/lean/ErdosProblems/Erdos1049/PaperCompleteR21/PrintedLogConstants.lean\#L195}{\nolinkurl{printed_bv_cutoff}}, \href{https://github.com/wcook04/plectis-erdos/blob/181078b6b009d905809cf2e007ad309386a3d1e0/lean/ErdosProblems/Erdos1049/PaperCompleteR21/PrintedLogConstants.lean\#L218}{\nolinkurl{printed_bvMu}}, \href{https://github.com/wcook04/plectis-erdos/blob/181078b6b009d905809cf2e007ad309386a3d1e0/lean/ErdosProblems/Erdos1049/PaperCompleteR21/PrintedLogConstants.lean\#L308}{\nolinkurl{printed_four_rpow_bvMu}}, \href{https://github.com/wcook04/plectis-erdos/blob/181078b6b009d905809cf2e007ad309386a3d1e0/lean/ErdosProblems/Erdos1049/PaperCompleteR21/PrintedLogConstants.lean\#L212}{\nolinkurl{paperBvMu_eq}}, \href{https://github.com/wcook04/plectis-erdos/blob/181078b6b009d905809cf2e007ad309386a3d1e0/lean/ErdosProblems/Erdos1049/PaperCompleteR21/PrintedContourConstants.lean\#L625}{\nolinkurl{printed_four_rpow_mu}}, \href{https://github.com/wcook04/plectis-erdos/blob/181078b6b009d905809cf2e007ad309386a3d1e0/lean/ErdosProblems/Erdos1049/PaperCompleteR21/PrintedContourConstants.lean\#L658}{\nolinkurl{four_rpow_mu_lt_thirtyOne}}, \href{https://github.com/wcook04/plectis-erdos/blob/181078b6b009d905809cf2e007ad309386a3d1e0/lean/ErdosProblems/Erdos1049/PaperCompleteR21/RationalBaseThreshold.lean\#L125}{\nolinkurl{thirtyoneFour_irrational}}, \href{https://github.com/wcook04/plectis-erdos/blob/181078b6b009d905809cf2e007ad309386a3d1e0/lean/ErdosProblems/Erdos1049/PaperCompleteR21/RationalBaseThreshold.lean\#L129}{\nolinkurl{thirtyoneFour_pow_irrational}}, \href{https://github.com/wcook04/plectis-erdos/blob/181078b6b009d905809cf2e007ad309386a3d1e0/lean/ErdosProblems/Erdos1049/PaperCompleteR21/RationalBaseThreshold.lean\#L134}{\nolinkurl{thirtyoneFour_ratio_chain}}, \href{https://github.com/wcook04/plectis-erdos/blob/181078b6b009d905809cf2e007ad309386a3d1e0/lean/ErdosProblems/Erdos1049/PaperCompleteR21/RationalBaseThreshold.lean\#L153}{\nolinkurl{inv_bvMu_eq}}, \href{https://github.com/wcook04/plectis-erdos/blob/181078b6b009d905809cf2e007ad309386a3d1e0/lean/ErdosProblems/Erdos1049/PaperCompleteR21/RationalBaseThreshold.lean\#L169}{\nolinkurl{thirtyoneFour_between_rpow}}, \href{https://github.com/wcook04/plectis-erdos/blob/181078b6b009d905809cf2e007ad309386a3d1e0/lean/ErdosProblems/Erdos1049/PaperCompleteR21/RationalBaseThreshold.lean\#L182}{\nolinkurl{thirtyoneFour_outside_bv_inside_contour}}.  \emph{Compared}, runs \href{https://github.com/wcook04/plectis-erdos-lean/actions/runs/35624228171}{\texttt{35624228171}}, \href{https://github.com/wcook04/plectis-erdos-lean/actions/runs/35674034595}{\texttt{35674034595}}.
\par\noindent\hangindent=1.5em Corollary~\ref{long1049:cor:rational-base-measure}: \href{https://github.com/wcook04/plectis-erdos/blob/181078b6b009d905809cf2e007ad309386a3d1e0/lean/ErdosProblems/Erdos1049/PaperCompleteR21/RationalBaseThreshold.lean\#L194}{\nolinkurl{rational_base_measure_uniform}}, \href{https://github.com/wcook04/plectis-erdos/blob/181078b6b009d905809cf2e007ad309386a3d1e0/lean/ErdosProblems/Erdos1049/PaperCompleteR21/RationalBaseThreshold.lean\#L205}{\nolinkurl{thirtyoneFour_power_measure_lt_301}}, \href{https://github.com/wcook04/plectis-erdos/blob/181078b6b009d905809cf2e007ad309386a3d1e0/lean/ErdosProblems/Erdos1049/PaperR17/SourceConsumers.lean\#L167}{\nolinkurl{rational_base_power_measure}}, \href{https://github.com/wcook04/plectis-erdos/blob/181078b6b009d905809cf2e007ad309386a3d1e0/lean/ErdosProblems/Erdos1049/PaperR17/SourceConsumers.lean\#L176}{\nolinkurl{thirtyone_four_power_measure_lt_301}}.  \emph{Compared}, runs \href{https://github.com/wcook04/plectis-erdos-lean/actions/runs/35544127144}{\texttt{35544127144}}, \href{https://github.com/wcook04/plectis-erdos-lean/actions/runs/35674034595}{\texttt{35674034595}}.
\par\noindent\hangindent=1.5em Theorem~\ref{long1049:res:archcap} (the Lean statement is stronger): \href{https://github.com/wcook04/plectis-erdos/blob/181078b6b009d905809cf2e007ad309386a3d1e0/lean/ErdosProblems/Erdos1049/PaperLongCapR9.lean\#L411}{\nolinkurl{long_record_archcap}}.  \emph{Compared}, run \href{https://github.com/wcook04/plectis-erdos-lean/actions/runs/35674034595}{\texttt{35674034595}}.
\par\noindent\hangindent=1.5em Corollary~\ref{long1049:cor:no-decay-below-square} (the Lean statement is stronger): \href{https://github.com/wcook04/plectis-erdos/blob/181078b6b009d905809cf2e007ad309386a3d1e0/lean/ErdosProblems/Erdos1049/PaperNoDecayR9.lean\#L69}{\nolinkurl{cleared_below_square_not_tendsto_zero}}.  \emph{Compared}, run \href{https://github.com/wcook04/plectis-erdos-lean/actions/runs/35544127144}{\texttt{35544127144}}.
\par\noindent\hangindent=1.5em Lemma~\ref{long1049:res:sourceheight}: \href{https://github.com/wcook04/plectis-erdos/blob/181078b6b009d905809cf2e007ad309386a3d1e0/lean/ErdosProblems/Erdos1049/PaperCompleteR21/SourceCoefficientHeights.lean\#L49}{\nolinkurl{exists_quadratic_source_height_bound}}, \href{https://github.com/wcook04/plectis-erdos/blob/181078b6b009d905809cf2e007ad309386a3d1e0/lean/ErdosProblems/Erdos1049/PaperCompleteR21/SourceCoefficientHeights.lean\#L40}{\nolinkurl{maxPairHeight_source_eq}}.
\par\noindent\hangindent=1.5em Theorem~\ref{long1049:res:powerbracket}: \href{https://github.com/wcook04/plectis-erdos/blob/181078b6b009d905809cf2e007ad309386a3d1e0/lean/ErdosProblems/Erdos1049/PaperFiniteAssembliesR7.lean\#L24}{\nolinkurl{power_bracket}}.  \emph{Compared}, run \href{https://github.com/wcook04/plectis-erdos-lean/actions/runs/35674034595}{\texttt{35674034595}}.
\par\noindent\hangindent=1.5em Corollary~\ref{long1049:res:sharpgaps}: \href{https://github.com/wcook04/plectis-erdos/blob/181078b6b009d905809cf2e007ad309386a3d1e0/lean/ErdosProblems/Erdos1049/PaperFiniteAssembliesR7.lean\#L34}{\nolinkurl{height_and_hankel_deficits}}.  \emph{Compared}, run \href{https://github.com/wcook04/plectis-erdos-lean/actions/runs/35674034595}{\texttt{35674034595}}.
\par\noindent\hangindent=1.5em Theorem~\ref{long1049:res:chargeceilings}: \href{https://github.com/wcook04/plectis-erdos/blob/181078b6b009d905809cf2e007ad309386a3d1e0/lean/ErdosProblems/Erdos1049/PaperFiniteAssembliesR7.lean\#L44}{\nolinkurl{charge_ceilings}}.  \emph{Compared}, run \href{https://github.com/wcook04/plectis-erdos-lean/actions/runs/35674034595}{\texttt{35674034595}}.
\par\noindent\hangindent=1.5em Theorem~\ref{long1049:res:zudilin-sharp-qorder}: \href{https://github.com/wcook04/plectis-erdos/blob/181078b6b009d905809cf2e007ad309386a3d1e0/lean/ErdosProblems/Erdos1049/AllRow/Producer.lean\#L173}{\nolinkurl{order_zudilinNormalizedHankelDet_all}}, \href{https://github.com/wcook04/plectis-erdos/blob/181078b6b009d905809cf2e007ad309386a3d1e0/lean/ErdosProblems/Erdos1049/AllRow/Producer.lean\#L199}{\nolinkurl{coeff_zudilinNormalizedHankelDet_all_rat}}.  \emph{Compared}, runs \href{https://github.com/wcook04/plectis-erdos-lean/actions/runs/35544127144}{\texttt{35544127144}}, \href{https://github.com/wcook04/plectis-erdos-lean/actions/runs/35674034595}{\texttt{35674034595}}.
\par\noindent\hangindent=1.5em Lemma~\ref{long1049:res:allrowinitial}: \href{https://github.com/wcook04/plectis-erdos/blob/181078b6b009d905809cf2e007ad309386a3d1e0/lean/ErdosProblems/Erdos1049/PaperCompleteR21/AllRowInitialCoefficient.lean\#L213}{\nolinkurl{all_row_initial}}, \href{https://github.com/wcook04/plectis-erdos/blob/181078b6b009d905809cf2e007ad309386a3d1e0/lean/ErdosProblems/Erdos1049/PaperCompleteR21/AllRowInitialCoefficient.lean\#L242}{\nolinkurl{all_row_initial_reciprocal}}, \href{https://github.com/wcook04/plectis-erdos/blob/181078b6b009d905809cf2e007ad309386a3d1e0/lean/ErdosProblems/Erdos1049/PaperCompleteR21/AllRowInitialCoefficient.lean\#L255}{\nolinkurl{all_row_initial_dvd}}, \href{https://github.com/wcook04/plectis-erdos/blob/181078b6b009d905809cf2e007ad309386a3d1e0/lean/ErdosProblems/Erdos1049/PaperCompleteR21/AllRowInitialCoefficient.lean\#L116}{\nolinkurl{paperE_eq_rowExponent}}, \href{https://github.com/wcook04/plectis-erdos/blob/181078b6b009d905809cf2e007ad309386a3d1e0/lean/ErdosProblems/Erdos1049/PaperCompleteR21/AllRowInitialCoefficient.lean\#L165}{\nolinkurl{paperReciprocal_zero}}, \href{https://github.com/wcook04/plectis-erdos/blob/181078b6b009d905809cf2e007ad309386a3d1e0/lean/ErdosProblems/Erdos1049/PaperCompleteR21/AllRowInitialCoefficient.lean\#L171}{\nolinkurl{paperReciprocal_rec}}, \href{https://github.com/wcook04/plectis-erdos/blob/181078b6b009d905809cf2e007ad309386a3d1e0/lean/ErdosProblems/Erdos1049/PaperCompleteR21/AllRowInitialCoefficient.lean\#L73}{\nolinkurl{paperRatio_agree}}.  \emph{Compared}, runs \href{https://github.com/wcook04/plectis-erdos-lean/actions/runs/35624228171}{\texttt{35624228171}}, \href{https://github.com/wcook04/plectis-erdos-lean/actions/runs/35674034595}{\texttt{35674034595}}.
\par\noindent\hangindent=1.5em Theorem~\ref{long1049:res:content}: \href{https://github.com/wcook04/plectis-erdos/blob/181078b6b009d905809cf2e007ad309386a3d1e0/lean/ErdosProblems/Erdos1049/PaperFiniteAssembliesR7.lean\#L59}{\nolinkurl{integer_scalar_content}}.  \emph{Compared}, run \href{https://github.com/wcook04/plectis-erdos-lean/actions/runs/35674034595}{\texttt{35674034595}}.
\par\noindent\hangindent=1.5em Theorem~\ref{long1049:res:endpoints}: \href{https://github.com/wcook04/plectis-erdos/blob/181078b6b009d905809cf2e007ad309386a3d1e0/lean/ErdosProblems/Erdos1049/PaperFiniteAssembliesR7.lean\#L74}{\nolinkurl{endpoint_residues}}.  \emph{Compared}, run \href{https://github.com/wcook04/plectis-erdos-lean/actions/runs/35674034595}{\texttt{35674034595}}.
\par\noindent\hangindent=1.5em Proposition~\ref{long1049:res:commonmult}: \href{https://github.com/wcook04/plectis-erdos/blob/181078b6b009d905809cf2e007ad309386a3d1e0/lean/ErdosProblems/Erdos1049/ZudilinConeArithmetic.lean\#L398}{\nolinkurl{commonMultiplier_not_two_not_three_of_endpoint_units}}.  \emph{Compared}, run \href{https://github.com/wcook04/plectis-erdos-lean/actions/runs/35674034595}{\texttt{35674034595}}.
\par\noindent\hangindent=1.5em Corollary~\ref{long1049:res:nomult}: \href{https://github.com/wcook04/plectis-erdos/blob/181078b6b009d905809cf2e007ad309386a3d1e0/lean/ErdosProblems/Erdos1049/PaperFiniteAssembliesR7.lean\#L220}{\nolinkurl{endpoint_scalar_content_exclusion}}.  \emph{Compared}, run \href{https://github.com/wcook04/plectis-erdos-lean/actions/runs/35674034595}{\texttt{35674034595}}.
\par\noindent\hangindent=1.5em Proposition~\ref{long1049:res:cyclounit}: \href{https://github.com/wcook04/plectis-erdos/blob/181078b6b009d905809cf2e007ad309386a3d1e0/lean/ErdosProblems/Erdos1049/ZudilinConeArithmetic.lean\#L302}{\nolinkurl{cyclotomicHomEval_isCoprime_mul}}.  \emph{Compared}, run \href{https://github.com/wcook04/plectis-erdos-lean/actions/runs/35674034595}{\texttt{35674034595}}.
\par\noindent\hangindent=1.5em Theorem~\ref{long1049:res:jetkernel}: \href{https://github.com/wcook04/plectis-erdos/blob/181078b6b009d905809cf2e007ad309386a3d1e0/lean/ErdosProblems/Erdos1049/PaperFiniteAssembliesR7.lean\#L248}{\nolinkurl{fourJet_paper_statement}}.  \emph{Compared}, run \href{https://github.com/wcook04/plectis-erdos-lean/actions/runs/35674034595}{\texttt{35674034595}}.
\par\noindent\hangindent=1.5em Corollary~\ref{long1049:res:rankfortyone}: \href{https://github.com/wcook04/plectis-erdos/blob/181078b6b009d905809cf2e007ad309386a3d1e0/lean/ErdosProblems/Erdos1049/PaperFiniteAssembliesR7.lean\#L149}{\nolinkurl{rank_fortyone}}.  \emph{Compared}, run \href{https://github.com/wcook04/plectis-erdos-lean/actions/runs/35674034595}{\texttt{35674034595}}.
\par\noindent\hangindent=1.5em Theorem~\ref{long1049:res:boundedfibre}: \href{https://github.com/wcook04/plectis-erdos/blob/181078b6b009d905809cf2e007ad309386a3d1e0/lean/ErdosProblems/Erdos1049/AdelicHeightBridge.lean\#L1832}{\nolinkurl{exists_ne_map_eq_map_ne_of_card_mul_lt}}.  \emph{Compared}, run \href{https://github.com/wcook04/plectis-erdos-lean/actions/runs/35544127144}{\texttt{35544127144}}.
\par\noindent\hangindent=1.5em Theorem~\ref{long1049:res:plucker-collapse}: \href{https://github.com/wcook04/plectis-erdos/blob/181078b6b009d905809cf2e007ad309386a3d1e0/lean/ErdosProblems/Erdos1049/PaperFiniteAssembliesR7.lean\#L269}{\nolinkurl{plucker_paper_statement}}.
\par\noindent\hangindent=1.5em Theorem~\ref{long1049:res:scalar} (the Lean statement is stronger): \href{https://github.com/wcook04/plectis-erdos/blob/181078b6b009d905809cf2e007ad309386a3d1e0/lean/ErdosProblems/Erdos1049/PaperFiniteAssembliesR7.lean\#L177}{\nolinkurl{scalar_margin}}.  \emph{Compared}, run \href{https://github.com/wcook04/plectis-erdos-lean/actions/runs/35674034595}{\texttt{35674034595}}.
\par\noindent\hangindent=1.5em Theorem~\ref{long1049:res:corridorbound}: \href{https://github.com/wcook04/plectis-erdos/blob/181078b6b009d905809cf2e007ad309386a3d1e0/lean/ErdosProblems/Erdos1049/RationalBaseLambert.lean\#L121}{\nolinkurl{coordinatewiseCorridor_implies_pow_lt_linear}}.  \emph{Compared}, run \href{https://github.com/wcook04/plectis-erdos-lean/actions/runs/35624228171}{\texttt{35624228171}}.
\par\noindent\hangindent=1.5em Proposition~\ref{long1049:res:exp}: \href{https://github.com/wcook04/plectis-erdos/blob/181078b6b009d905809cf2e007ad309386a3d1e0/lean/ErdosProblems/Erdos1049/RationalBaseLambert.lean\#L142}{\nolinkurl{three_mul_lt_two_pow_succ}}.  \emph{Compared}, run \href{https://github.com/wcook04/plectis-erdos-lean/actions/runs/35624228171}{\texttt{35624228171}}.
\par\noindent\hangindent=1.5em Theorem~\ref{long1049:res:nocorridor}: \href{https://github.com/wcook04/plectis-erdos/blob/181078b6b009d905809cf2e007ad309386a3d1e0/lean/ErdosProblems/Erdos1049/RationalBaseLambert.lean\#L155}{\nolinkurl{threeHalves_no_coordinatewiseCorridor}}.  \emph{Compared}, run \href{https://github.com/wcook04/plectis-erdos-lean/actions/runs/35544127144}{\texttt{35544127144}}.
\par\noindent\hangindent=1.5em Theorem~\ref{long1049:res:tailrec}: \href{https://github.com/wcook04/plectis-erdos/blob/181078b6b009d905809cf2e007ad309386a3d1e0/lean/ErdosProblems/Erdos1049/RationalBaseLambert.lean\#L187}{\nolinkurl{rationalBaseClearedTailQ_succ}}.  \emph{Compared}, run \href{https://github.com/wcook04/plectis-erdos-lean/actions/runs/35544127144}{\texttt{35544127144}}.
\par\noindent\hangindent=1.5em Theorem~\ref{long1049:res:forcing}: \href{https://github.com/wcook04/plectis-erdos/blob/181078b6b009d905809cf2e007ad309386a3d1e0/lean/ErdosProblems/Erdos1049/PaperFiniteAssembliesR7.lean\#L186}{\nolinkurl{forcing_term}}.  \emph{Compared}, run \href{https://github.com/wcook04/plectis-erdos-lean/actions/runs/35674034595}{\texttt{35674034595}}.
\par\noindent\hangindent=1.5em Theorem~\ref{long1049:res:sevenhalves}: \href{https://github.com/wcook04/plectis-erdos/blob/181078b6b009d905809cf2e007ad309386a3d1e0/lean/ErdosProblems/Erdos1049/RationalBaseLambert.lean\#L83}{\nolinkurl{sevenHalves_archimedean_height_condition}}.  \emph{Compared}, run \href{https://github.com/wcook04/plectis-erdos-lean/actions/runs/35624228171}{\texttt{35624228171}}.
\par\noindent\hangindent=1.5em Proposition~\ref{long1049:res:pade}: \href{https://github.com/wcook04/plectis-erdos/blob/181078b6b009d905809cf2e007ad309386a3d1e0/lean/ErdosProblems/Erdos1049/PaperFiniteAssembliesR7.lean\#L201}{\nolinkurl{pade_summand_bound_and_gap}}.  \emph{Compared}, run \href{https://github.com/wcook04/plectis-erdos-lean/actions/runs/35674034595}{\texttt{35674034595}}.
\par\endgroup
% END GENERATED CONCORDANCE

\paragraph{The Lambert identity and the integer case.}
The cited \emph{Formal Conjectures} file contains unproved declarations
for the conjecture and the integer-base irrationality theorem, as well
as a proved identity between the two series~\cite{formalconjectures1049}.
For $t>1$, that identity is an equality of convergent series. The file
also treats a different range using Lean's convention that an
unsummable series has totalised sum zero; that convention is not used
in the present argument. The proved identity should therefore not be
confused with a proved irrationality statement. The rational-base
formal proofs used here are in the separate sources described below.


\begingroup
\raggedright

The public sources supplied for this prose revision are at commit
\texttt{6b78209ab63a8c643281115f8628a3be79ff7ec7}; the parallel release is
at \texttt{52f29ad173b04e3bac941b3663f2b9aebe5de0bb}.
An earlier editorial audit used
\texttt{3d6d938d696fed0fb71dd55115a18a73738ff223}. All existing links
retain their original commits and line numbers; they have not been
retargeted to the supplied snapshot.
The supplied index distinguishes public CI-checked declarations,
release-only declarations and declarations outside the checked build.
Those categories must not be conflated.

The main irrationality and measure proofs for the constructed forms are in
\texttt{Erdos1049/PaperR17/SourceConsumers.lean}. The file constructs the
cancelled polynomial forms, then proves irrationality in the stated
region, irrationality for positive integral powers of $31/4$, and the
irrationality-exponent bound.
The all-rank determinant proofs are in
\texttt{Erdos1049/AllRow/Producer.lean}. That file proves the first
nonzero term of every transformed row and then deduces the order and
leading coefficient of the determinant. Both constructions therefore
supply their mathematical objects rather than assume their existence.

Other modules check the clearing inequalities and recurrence identities, degree
bookkeeping, endpoint congruences, finite collision principles and
specified exponent models. The ordinary theorem for unimodular rows must be distinguished from
the checked version requiring an invertible coordinate in every row.
The supplied scripts reproduce the source-polynomial checks in
Section~\ref{long1049:sec:receipts} and the finite coefficient-moment
tests in Section~\ref{long1049:sec:coefficient-questions}. The five complete polynomial contents and seventy-six cyclotomic
residue witnesses are also supplied with exact reproduction scripts
and direct polynomial determinant checks. None of these computations
is a Lean declaration. The positivity of the eight unshifted determinants
used in Proposition~\ref{long1049:res:finite-pencil} also has a separate
Lean proof, described in Section~\ref{long1049:sec:coefficient-questions}.

The supplied index reports earlier build results. The kernel checks and
axioms stated in the coverage subsection above belong to the revisions
named there. The build status in that index and
the scope of a theorem's hypotheses are separate matters. Historical
links remain unchanged and should be read at their own pinned revisions.

\par\endgroup

% The linked source manifest is retained in the authored TeX for automated
% validation, and kept outside the printed note.
% Historical hidden declaration table is preserved in the pass-one provenance.

\paragraph{The hypotheses in the modular row argument.}
The ordinary proof uses rowwise unimodularity and zero adjacent minors.
The linked Lean statements instead assume that each second coordinate is
a unit. The example $(2,3)$ modulo six satisfies the former condition but
not the latter. Thus the ordinary generalisation is not silently included
in the scope of the checked unit-coordinate statements. Neither version
establishes the required hypotheses for an actual all-tail approximation
family, or proves nonzero analytic remainder for a selector difference.
% Part: part family_catalogue
\section{Publication scope}
\label{sec:erdos-1049-complete-family-map}
The project's archival catalogue groups the results for publication.
Those classifications do not change any theorem's hypotheses or show
that its proof was included in a particular Lean build. The mathematical
statements in this paper, the proofs linked beside them, and the build
information described in Appendix~\ref{long1049:app:index} must be read
separately. The introduction locates the principal arguments.
% Part: part back
\begin{thebibliography}{99}
\small
\raggedright
\setlength{\itemsep}{2pt}
\bibitem{erdos1948} P.~Erd\H{o}s,
  \href{https://users.renyi.hu/~p_erdos/1948-04.pdf}{\emph{On arithmetical
  properties of Lambert series}},
  J. Indian Math. Soc. (N.S.) \textbf{12} (1948), 63--66.
\bibitem{erdos1988} P.~Erd\H{o}s, \emph{On the irrationality of certain series:
  problems and results}, in A.~Baker (ed.), \emph{New Advances in Transcendence
  Theory}, Cambridge UP, 1988, pp.~102--109,
  doi:\href{https://doi.org/10.1017/CBO9780511897184.009}{10.1017/CBO9780511897184.009}.
\bibitem{borwein1991} P.~B. Borwein, \emph{On the irrationality of
  $\sum1/(q^{n}+r)$},
  J. Number Theory \textbf{37} (1991), no.~3, 253--259,
  doi:\href{https://doi.org/10.1016/S0022-314X(05)80041-1}{10.1016/S0022-314X(05)80041-1}.
\bibitem{borwein1992} P.~B. Borwein, \emph{On the irrationality of certain
  series}, Math. Proc. Cambridge Philos. Soc. \textbf{112} (1992), no.~1,
  141--146,
  doi:\href{https://doi.org/10.1017/S030500410007081X}{10.1017/S030500410007081X}.
  Van Assche cites Lemma~2 for the neighbouring little-$q$-Legendre
  evaluation used above.
\bibitem{az1998} T.~Amdeberhan and D.~Zeilberger, \emph{$q$-Ap\'ery
  irrationality proofs by $q$-WZ pairs}, Adv. Appl. Math. \textbf{20} (1998),
  no.~2, 275--283,
  doi:\href{https://doi.org/10.1006/aama.1997.0565}{10.1006/aama.1997.0565};
  arXiv:\href{https://arxiv.org/abs/math/9804122v1}{math/9804122v1}.
  Page references are to arXiv:math/9804122v1.
\bibitem{bv1994} P.~Bundschuh and K.~V\"a\"an\"anen,
  \href{https://numdam.org/item/CM_1994__91_2_175_0.pdf}{\emph{Arithmetical
  investigations of a certain infinite product}},
  Compositio Math. \textbf{91} (1994), no.~2, 175--199.
\bibitem{zudilin2002} W.~Zudilin, \emph{Remarks on irrationality of
  \(q\)-harmonic series}, Manuscripta Math. \textbf{107} (2002), no.~4,
  463--477,
  doi:\href{https://doi.org/10.1007/s002290200249}{10.1007/s002290200249}.
\bibitem{zudilin2004} W.~Zudilin,
  \href{https://geodesic.mathdoc.fr/articles/10.4064/aa111-2-4/}{\emph{Heine's
  basic transform and a permutation group for \(q\)-harmonic series}},
  Acta Arith. \textbf{111} (2004), no.~2, 153--164,
  doi:\href{https://doi.org/10.4064/aa111-2-4}{10.4064/aa111-2-4}.
  Page references are to the printed journal pages.
\bibitem{zudilin2016} W.~Zudilin,
  \href{https://arxiv.org/abs/1601.02688v2}{\emph{On the irrationality of
  generalized \(q\)-logarithm}}, arXiv:1601.02688; Res. Number Theory
  \textbf{2} (2016), Art.~15,
  doi:\href{https://doi.org/10.1007/s40993-016-0042-x}{10.1007/s40993-016-0042-x}.
  Page references are to arXiv:1601.02688v2.
  The remark that the results extend to non-integer \(p=r/s\), \(|p|>1\), under
  an assumption \(\log|r|>c\log|s|\) for a computable \(c>0\), is in Section~2,
  p.~4, in the paragraph beginning ``Finally, we remark''; no value of \(c\) is
  computed there, and the remark is made for the generalized \(q\)-logarithm of
  that paper.
\bibitem{zudilin2017det} W.~Zudilin, \emph{A determinantal approach to
  irrationality}, Constr. Approx. \textbf{45} (2017), no.~2, 301--310,
  doi:\href{https://doi.org/10.1007/s00365-016-9333-7}{10.1007/s00365-016-9333-7};
  arXiv:\href{https://arxiv.org/abs/1507.05697v1}{1507.05697v1}.
  Page and equation references are to arXiv:1507.05697v1.
\bibitem{bundschuhzudilin2008} P.~Bundschuh and W.~Zudilin, \emph{Rational
  approximations to a \(q\)-analogue of \(\pi\) and some other \(q\)-series},
  in H.~P. Schlickewei, K.~Schmidt and R.~F. Tichy (eds.), \emph{Diophantine
  Approximation}, Dev. Math. \textbf{16}, Springer, 2008, pp.~123--139,
  doi:\href{https://doi.org/10.1007/978-3-211-74280-8_6}{10.1007/978-3-211-74280-8\_6}.
\bibitem{krvz2009} C.~Krattenthaler, I.~Rochev, K.~V\"a\"an\"anen and
  W.~Zudilin, \emph{On the non-quadraticity of values of the
  \(q\)-exponential function and related \(q\)-series}, Acta Arith.
  \textbf{136} (2009), no.~3, 243--269,
  doi:\href{https://doi.org/10.4064/aa136-3-4}{10.4064/aa136-3-4};
  arXiv:\href{https://arxiv.org/abs/0812.2921v1}{0812.2921v1}.
  Page references are to arXiv:0812.2921v1.
\bibitem{duverney1996} D.~Duverney,
  \href{https://numdam.org/item/JTNB_1996__8_1_173_0.pdf}{\emph{\`A propos de la
  s\'erie $\sum_{n\ge1}x^{n}/(q^{n}-1)$}},
  J. Th\'eor. Nombres Bordeaux \textbf{8} (1996), no.~1, 173--181.
  Th\'eor\`eme~2 on p.~174 gives the rational-base region
  $\log|s|/\log|r|<\frac13(1-3/\pi^{2})=0.2320\ldots$ for this series;
  Th\'eor\`eme~1, for a general numerator, is weaker.
\bibitem{matalaaho2006} T.~Matala-aho, K.~V\"a\"an\"anen and W.~Zudilin,
  \href{https://doi.org/10.1090/S0025-5718-05-01812-0}{\emph{New irrationality
  measures for \(q\)-logarithms}}, Math. Comp. \textbf{75} (2006), no.~254,
  879--889, doi:10.1090/S0025-5718-05-01812-0.  The hypothesis
  \(p=1/q\in\Z\mathbin{\backslash}\{0,\pm1\}\) is carried in the abstract on p.~879 and in
  both theorem statements on p.~880, where the authors also record that their
  methods do not sharpen the \(q\)-harmonic case of~\cite{zudilin2004}.
\bibitem{adamczewskibell2013} B.~Adamczewski and J.~P. Bell,
  \href{https://arxiv.org/abs/1303.2019v1}{\emph{A problem about Mahler
  functions}}, Ann. Sc. Norm. Super. Pisa Cl. Sci. \textbf{17} (2017), no.~4,
  1301--1355; arXiv:\href{https://arxiv.org/abs/1303.2019v1}{1303.2019v1}, 2013.
  Theorem~1.1 on p.~6 of arXiv:1303.2019v1: over a field of characteristic
  zero, a power series is both \(k\)- and \(\ell\)-Mahler for multiplicatively
  independent \(k,\ell\) if and only if it is a rational function.  Bell and
  Smertnig cite it as Theorem~1.3 of the journal version
  \cite[Thm.~2.5, p.~5]{bellsmertnig2026}.
\bibitem{rhinviola1996} G.~Rhin and C.~Viola,
  \href{https://geodesic.mathdoc.fr/articles/10.4064/aa-77-1-23-56/}
  {\emph{On a permutation group related to $\zeta(2)$}},
  Acta Arith. \textbf{77} (1996), no.~1, 23--56,
  doi:10.4064/aa-77-1-23-56.
\bibitem{vanassche2001} W.~Van Assche,
  \href{https://arxiv.org/abs/math/0101187v1}{\emph{Little \(q\)-Legendre
  polynomials and irrationality of certain Lambert series}},
  Ramanujan J. \textbf{5} (2001), no.~3, 295--310,
  doi:\href{https://doi.org/10.1023/A:1012930828917}{10.1023/A:1012930828917}.
  Page references are to arXiv:math/0101187v1.
\bibitem{coussementsmet2007} J.~Coussement and C.~Smet, \emph{Irrationality
  proof of certain Lambert series using little \(q\)-Jacobi polynomials},
  J. Comput. Appl. Math. \textbf{233} (2009), no.~3, 680--690,
  doi:\href{https://doi.org/10.1016/j.cam.2009.02.036}{10.1016/j.cam.2009.02.036};
  arXiv:\href{https://arxiv.org/abs/math/0701345v1}{math/0701345v1}.
  Page references are to arXiv:math/0701345v1.
\bibitem{koizumiyokoi2026} J.~Koizumi and A.~Yokoi, \emph{Ap\'ery-type
  approximations and irrationality measures for certain \(q\)-series},
  arXiv:\href{https://arxiv.org/abs/2608.26918v1}{2608.26918v1},
  27 August 2026.
\bibitem{vandehey2013} J.~Vandehey,
  \href{https://arxiv.org/abs/1206.0340v1}{\emph{On an incomplete argument of
  Erd\H{o}s on the irrationality of Lambert series}}, Integers
  \textbf{13} (2013), Paper A58. Page references are to
  arXiv:1206.0340v1 (2012).
\bibitem{lucatachiya2017} F.~Luca and Y.~Tachiya,
  \href{https://www.kurims.kyoto-u.ac.jp/~kyodo/kokyuroku/contents/pdf/2014-14.pdf}
  {\emph{Linear independence results for the values of divisor functions
  series}}, RIMS K\^oky\^uroku No.~2014 (2017), 138--150.
  Theorem~A on p.~139 restates the periodic-coefficient irrationality theorem;
  Example~1 on p.~140 gives the divisor-function specialization.
\bibitem{duverneytachiya2019} D.~Duverney and Y.~Tachiya, \emph{Refinement of
  the Chowla--Erd\H{o}s method and linear independence of certain Lambert
  series}, Forum Math. \textbf{31} (2019), no.~6, 1557--1566,
  doi:\href{https://doi.org/10.1515/forum-2018-0299}{10.1515/forum-2018-0299}.
  Page references are to the
  \href{https://danielduverney.fr/documents/theorie-des-nombres/DuverneyTachiya190522.pdf}{authors'
  version}.
\bibitem{rivin2026} I.~Rivin,
\href{https://arxiv.org/abs/2604.25151v1}{\emph{Zero Coefficients of Rational
  Power Series and Rational Lambert Series}}, arXiv:2604.25151v1,
  28 April 2026.
  Theorem~1.1 is on p.~2 and proved on pp.~6--7; the periodic-coefficient
  Corollary~6.4 is on p.~9.
\bibitem{kovactao2024} V.~Kova\v{c} and T.~Tao,
  \href{https://arxiv.org/abs/2406.17593v4}{\emph{On several irrationality
  problems for Ahmes series}}, Acta Math. Hungar. \textbf{175} (2025), no.~2,
  572--608,
  doi:\href{https://doi.org/10.1007/s10474-025-01528-0}{10.1007/s10474-025-01528-0};
  arXiv:\href{https://arxiv.org/abs/2406.17593v4}{2406.17593v4}.
  Page references are to arXiv:2406.17593v4.
\bibitem{erdosproblems} T.~F. Bloom,
  \href{https://www.erdosproblems.com/1049}{\emph{Erd\H{o}s Problem \#1049}},
  \texttt{erdosproblems.com/1049}. Historical snapshot cited in the supplied
  manuscript: accessed 28 July 2026, displaying ``last edited
  28 September 2025''.
\bibitem{bellsmertnig2026} J.~Bell and D.~Smertnig,
  \href{https://arxiv.org/abs/2603.23456v1}{\emph{Mahler series with
  multiplicative coefficient sequences}}, arXiv:2603.23456v1,
  24 March 2026.  Theorem~1.3 is on pp.~2--3; its stated consequences on
  p.~3 include that the divisor and totient generating series are not
  $k$-Mahler for any $k\ge2$.
\bibitem{dlmf} F.~W.~J. Olver et al.\ (eds.), \emph{NIST Digital Library of
  Mathematical Functions},
  \href{https://dlmf.nist.gov/17.2.E37}{\texttt{dlmf.nist.gov/17.2.E37}},
  Eq.~17.2.37 in \S17.2(iii), accessed 15 September 2026.
\bibitem{formalconjectures1049} The Formal Conjectures Authors,
  \href{https://github.com/google-deepmind/formal-conjectures/blob/f776d2f2039351b00737ffcafb9d7d7666e1d9af/FormalConjectures/ErdosProblems/1049.lean}
  {\emph{FormalConjectures.ErdosProblems.\texttt{1049}}},
  Lean source at commit \texttt{f776d2f}, 2026, accessed 28 July 2026.  The
  irrationality declarations are unproved; the Lambert-series identity is
  proved.
\bibitem{postelmansvanassche2007} K.~Postelmans and W.~Van Assche,
  \href{https://arxiv.org/abs/math/0604312v1}{\emph{Irrationality of
  $\zeta_q(1)$ and $\zeta_q(2)$}}, J. Number Theory \textbf{126} (2007),
  no.~1, 119--154,
  doi:\href{https://doi.org/10.1016/j.jnt.2006.11.011}{10.1016/j.jnt.2006.11.011}.
  Page references are to arXiv:math/0604312v1 (2006).
\bibitem{krattenthaler2021} C.~Krattenthaler,
  \href{https://arxiv.org/abs/2103.03969v1}{\emph{A determinant identity for
  moments of orthogonal polynomials that implies Uvarov's formula for the
  orthogonal polynomials of rationally related densities}},
  arXiv:2103.03969v1 (2021). Page references are to this version.
\bibitem{wangzhu2016} Y.~Wang and B.-X.~Zhu,
  \href{https://arxiv.org/abs/1612.04114v1}{\emph{Log-convex and Stieltjes
  moment sequences}}, Adv. Appl. Math. \textbf{81} (2016), 115--127,
  doi:\href{https://doi.org/10.1016/j.aam.2016.06.008}{10.1016/j.aam.2016.06.008}.
  Page references are to arXiv:1612.04114v1.
\bibitem{stanley2016} R.~P. Stanley, \emph{Smith normal form in combinatorics},
  J. Combin. Theory Ser. A \textbf{144} (2016), 476--495,
  doi:\href{https://doi.org/10.1016/j.jcta.2016.06.013}{10.1016/j.jcta.2016.06.013};
  arXiv:\href{https://arxiv.org/abs/1602.00166v1}{1602.00166v1}.
  Page references are to arXiv:1602.00166v1.
\bibitem{duverney2011} D.~Duverney,
  \href{https://doi.org/10.1063/1.3630035}{\emph{Arithmetical functions and
  irrationality of Lambert series}}, AIP Conf. Proc. \textbf{1385} (2011),
  5--16, doi:10.1063/1.3630035. Theorem and page references are to the
  12-page author manuscript supplied with the research archive.
\bibitem{berg2007} C.~Berg, \emph{On powers of Stieltjes moment sequences, II},
  J. Comput. Appl. Math. \textbf{199} (2007), 23--38;
  arXiv:\href{https://arxiv.org/abs/math/0412340v1}{math/0412340v1}.
  Theorem references use the preprint; Theorem~5.1 treats factorial powers.
\bibitem{sw2024} A.~D. Sokal and J.~Walrad,
  \emph{Continued-fraction characterization of Stieltjes moment sequences
  with support in $[\xi,\infty)$},
  \href{https://arxiv.org/abs/2404.12131v1}{arXiv:2404.12131v1}, 2024.
  The classical Stieltjes criterion is recalled on pp.~1--2.
\bibitem{lrz2017} H.~Liang, J.~Remmel and S.~Zheng,
  \emph{Stieltjes moment sequences of polynomials},
  \href{https://arxiv.org/abs/1710.05795v1}{arXiv:1710.05795v1}, 2017.
\bibitem{psz2023} M.~P\'etr\'eolle, A.~D. Sokal and B.-X.~Zhu,
  \emph{Lattice paths and branched continued fractions: An infinite sequence
  of generalizations of the Stieltjes--Rogers and Thron--Rogers polynomials,
  with coefficientwise Hankel-total positivity}, Mem. Amer. Math. Soc.
  \textbf{291} (2023), no.~1450,
  doi:\href{https://doi.org/10.1090/memo/1450}{10.1090/memo/1450};
  \href{https://arxiv.org/abs/1807.03271v2}{arXiv:1807.03271v2}.
  Theorem and page references use that preprint version.
\bibitem{middeke2021} J.~Middeke, D.~J. Jeffrey and C.~Koutschan,
  \emph{Common Factors in Fraction-Free Matrix Decompositions},
  Math. Comput. Sci. \textbf{15} (2021), no.~4, 589--608,
  doi:\href{https://doi.org/10.1007/s11786-020-00495-9}{10.1007/s11786-020-00495-9};
  \href{https://arxiv.org/abs/2005.12380v1}{arXiv:2005.12380v1}.
  Section and theorem references use the preprint.
\bibitem{golubwelsch1969} G.~H. Golub and J.~H. Welsch,
  \emph{Calculation of Gauss Quadrature Rules}, Math. Comp. \textbf{23}
  (1969), no.~106, 221--230,
  doi:\href{https://doi.org/10.1090/S0025-5718-69-99647-1}{10.1090/S0025-5718-69-99647-1}.
\bibitem{vinroot2010} C.~R. Vinroot,
  \href{https://arxiv.org/abs/1011.0984}{\emph{Multivariate Rogers--Szeg\H{o}
  polynomials and flags in finite vector spaces}}, arXiv:1011.0984v1,
  3 November 2010, abstract accessed 20 September 2026.  Cited for the
  identification of the sum of all $q$-multinomial coefficients of fixed
  degree and length with a flag count, and for its recursion generalising the
  Galois numbers.  The factorisation of the moment weights is proved here.
\end{thebibliography}
